\documentclass[11pt,openany]{book}
\usepackage[letterpaper,margin=0.9in,headheight=15pt]{geometry}
\usepackage[T1]{fontenc}
\usepackage{lmodern}
\usepackage{microtype}
\usepackage{xcolor}
\usepackage{booktabs,longtable,array}
\newcolumntype{L}[1]{>{\raggedright\arraybackslash}p{#1}}
\usepackage{upquote}
\usepackage{fvextra}
\usepackage{fancyhdr}
\usepackage{tikz}
\usetikzlibrary{arrows.meta,positioning,shapes.geometric,calc,fit,backgrounds}
\usepackage[shortlabels]{enumitem}
\usepackage[most]{tcolorbox}
\usepackage{titlesec}
\usepackage{caption}
\usepackage{needspace}
\usepackage{float}
\usepackage{adjustbox}
\usepackage{xurl}
\usepackage{makeidx}
\makeindex
\usepackage{hyperref}% load last

% ---------------------------------------------------------------- palette
\definecolor{navy}{HTML}{16324A}
\definecolor{teal}{HTML}{007F83}
\definecolor{paper}{HTML}{F3F6F8}
\definecolor{muted}{HTML}{536675}
\definecolor{rust}{HTML}{9C472B}
\definecolor{gold}{HTML}{8A6A12}
\definecolor{leaf}{HTML}{3D6F2E}

\hypersetup{colorlinks=true,linkcolor=teal,urlcolor=teal,
  pdftitle={Modern Perl Concurrency},
  pdfauthor={Scott Muller},
  pdfsubject={Synchronization, Distributed Coordination and Application Architecture; engineering handbook, edition 2.1.1},
  pdfkeywords={Perl, concurrency, locking, MySQL, Redis, Valkey, Memcached}}

% ---------------------------------------------------------------- layout
\pagestyle{fancy}\fancyhf{}
\newcommand{\currentpart}{Handbook}
\renewcommand{\chaptermark}[1]{\markboth{#1}{}}
\fancyhead[L]{\scriptsize\sffamily\color{muted}%
  \makebox[\headwidth][l]{%
    \makebox[.27\headwidth][l]{\begin{adjustbox}{max width=.25\headwidth}\currentpart\end{adjustbox}}%
    \makebox[.65\headwidth][c]{\begin{adjustbox}{max width=.62\headwidth}\leftmark\end{adjustbox}}%
    \makebox[.08\headwidth][r]{\thepage}}}
\fancyfoot[C]{\scriptsize\sffamily\color{muted}Modern Perl Concurrency \textperiodcentered{} Scott Muller}
\renewcommand{\headrulewidth}{0.4pt}
\fancypagestyle{plain}{\fancyhf{}\fancyfoot[C]{\small\sffamily\color{muted}\thepage\quad\textperiodcentered\quad Scott Muller}\renewcommand{\headrulewidth}{0pt}}
\raggedbottom
\setlength{\emergencystretch}{3em}
\setlength{\parindent}{0pt}\setlength{\parskip}{6pt}
\setcounter{tocdepth}{2}
\makeatletter
\renewcommand*\l@section{\@dottedtocline{1}{1.5em}{3.1em}}
\renewcommand*\l@subsection{\@dottedtocline{2}{4.6em}{3.5em}}
\makeatother
\setlength{\LTpre}{6pt}
\setlength{\LTpost}{7pt}
\renewcommand{\arraystretch}{1.10}
\clubpenalty=10000
\widowpenalty=10000
\setlength{\textfloatsep}{14pt plus 3pt minus 3pt}
\setlength{\intextsep}{12pt plus 3pt minus 3pt}
\setlist{itemsep=2pt,topsep=4pt,parsep=0pt}

\titleformat{\chapter}[display]
  {\sffamily\bfseries\color{navy}}
  {\large\color{teal}CHAPTER \thechapter}{4pt}{\huge}
  [\vspace{4pt}{\color{teal}\titlerule[1.2pt]}]
\titlespacing*{\chapter}{0pt}{-12pt}{16pt}
\titleformat{\section}{\sffamily\bfseries\Large\color{navy}}{\thesection}{0.75em}{}
\titleformat{\subsection}{\sffamily\bfseries\large\color{navy}}{\thesubsection}{0.75em}{}
\titlespacing*{\section}{0pt}{14pt}{4pt}

\captionsetup{font={small,sf},labelfont={bf,color=navy},textfont={color=muted},
  singlelinecheck=false,skip=4pt}

% ---------------------------------------------------------------- figures
\tikzset{
  every picture/.append style={font=\sffamily\small},
  every node/.append style={execute at begin node={\hyphenpenalty=10000\exhyphenpenalty=10000}},
  box/.style={draw=navy,fill=paper,rounded corners=2pt,align=center,inner sep=4pt,minimum height=7mm},
  okbox/.style={box,draw=teal,fill=teal!8},
  badbox/.style={box,draw=rust,fill=rust!8},
  decision/.style={diamond,aspect=2.4,draw=navy,fill=white,align=center,inner sep=1pt,font=\sffamily\footnotesize},
  state/.style={circle,draw=navy,fill=paper,minimum size=13mm,align=center,font=\sffamily\footnotesize},
  arr/.style={-{Latex[length=2mm]},thick,draw=teal},
  badarr/.style={-{Latex[length=2mm]},thick,draw=rust},
  thinarr/.style={-{Latex[length=1.6mm]},draw=navy},
  lbl/.style={font=\sffamily\scriptsize,text=muted,align=center},
  lane/.style={draw=muted!60,dashed},
  span/.style={line width=5pt,draw=teal},
  badspan/.style={line width=5pt,draw=rust},
  waitspan/.style={line width=5pt,draw=muted!35},
}
\captionsetup[figure]{font={small,sf},labelfont={bf,color=navy},textfont={color=muted},skip=6pt}
\newcommand{\fig}[2]{\begin{figure}[!htbp]\centering\begin{adjustbox}{max width=\linewidth}\input{figures/#1}\end{adjustbox}\caption{#2}\label{fig:#1}\end{figure}}

% ---------------------------------------------------------------- code presentation
% Shared code presentation. Content is highlighted offline; no shell escape.

\makeatletter
\def\PG@reset{\let\PG@it=\relax \let\PG@bf=\relax%
    \let\PG@ul=\relax \let\PG@tc=\relax%
    \let\PG@bc=\relax \let\PG@ff=\relax}
\def\PG@tok#1{\csname PG@tok@#1\endcsname}
\def\PG@toks#1+{\ifx\relax#1\empty\else%
    \PG@tok{#1}\expandafter\PG@toks\fi}
\def\PG@do#1{\PG@bc{\PG@tc{\PG@ul{%
    \PG@it{\PG@bf{\PG@ff{#1}}}}}}}
\def\PG#1#2{\PG@reset\PG@toks#1+\relax+\PG@do{#2}}

\@namedef{PG@tok@}{\def\PG@tc##1{\textcolor[rgb]{0.09,0.20,0.29}{##1}}}
\@namedef{PG@tok@w}{\def\PG@tc##1{\textcolor[rgb]{0.09,0.20,0.29}{##1}}}
\@namedef{PG@tok@err}{\def\PG@tc##1{\textcolor[rgb]{0.09,0.20,0.29}{##1}}}
\@namedef{PG@tok@c}{\def\PG@tc##1{\textcolor[rgb]{0.33,0.40,0.46}{##1}}}
\@namedef{PG@tok@k}{\def\PG@tc##1{\textcolor[rgb]{0.00,0.42,0.44}{##1}}}
\@namedef{PG@tok@kt}{\def\PG@tc##1{\textcolor[rgb]{0.00,0.42,0.44}{##1}}}
\@namedef{PG@tok@n}{\def\PG@tc##1{\textcolor[rgb]{0.09,0.20,0.29}{##1}}}
\@namedef{PG@tok@nf}{\def\PG@tc##1{\textcolor[rgb]{0.14,0.36,0.58}{##1}}}
\@namedef{PG@tok@nc}{\def\PG@tc##1{\textcolor[rgb]{0.14,0.36,0.58}{##1}}}
\@namedef{PG@tok@nb}{\def\PG@tc##1{\textcolor[rgb]{0.00,0.42,0.44}{##1}}}
\@namedef{PG@tok@nv}{\def\PG@tc##1{\textcolor[rgb]{0.15,0.29,0.44}{##1}}}
\@namedef{PG@tok@s}{\def\PG@tc##1{\textcolor[rgb]{0.60,0.25,0.13}{##1}}}
\@namedef{PG@tok@m}{\def\PG@tc##1{\textcolor[rgb]{0.46,0.35,0.07}{##1}}}
\@namedef{PG@tok@o}{\def\PG@tc##1{\textcolor[rgb]{0.21,0.34,0.43}{##1}}}
\@namedef{PG@tok@p}{\def\PG@tc##1{\textcolor[rgb]{0.09,0.20,0.29}{##1}}}
\@namedef{PG@tok@g}{\def\PG@tc##1{\textcolor[rgb]{0.09,0.20,0.29}{##1}}}
\@namedef{PG@tok@kc}{\def\PG@tc##1{\textcolor[rgb]{0.00,0.42,0.44}{##1}}}
\@namedef{PG@tok@kd}{\def\PG@tc##1{\textcolor[rgb]{0.00,0.42,0.44}{##1}}}
\@namedef{PG@tok@kn}{\def\PG@tc##1{\textcolor[rgb]{0.00,0.42,0.44}{##1}}}
\@namedef{PG@tok@kp}{\def\PG@tc##1{\textcolor[rgb]{0.00,0.42,0.44}{##1}}}
\@namedef{PG@tok@kr}{\def\PG@tc##1{\textcolor[rgb]{0.00,0.42,0.44}{##1}}}
\@namedef{PG@tok@na}{\def\PG@tc##1{\textcolor[rgb]{0.09,0.20,0.29}{##1}}}
\@namedef{PG@tok@bp}{\def\PG@tc##1{\textcolor[rgb]{0.00,0.42,0.44}{##1}}}
\@namedef{PG@tok@no}{\def\PG@tc##1{\textcolor[rgb]{0.09,0.20,0.29}{##1}}}
\@namedef{PG@tok@nd}{\def\PG@tc##1{\textcolor[rgb]{0.09,0.20,0.29}{##1}}}
\@namedef{PG@tok@ni}{\def\PG@tc##1{\textcolor[rgb]{0.09,0.20,0.29}{##1}}}
\@namedef{PG@tok@ne}{\def\PG@tc##1{\textcolor[rgb]{0.09,0.20,0.29}{##1}}}
\@namedef{PG@tok@fm}{\def\PG@tc##1{\textcolor[rgb]{0.14,0.36,0.58}{##1}}}
\@namedef{PG@tok@py}{\def\PG@tc##1{\textcolor[rgb]{0.09,0.20,0.29}{##1}}}
\@namedef{PG@tok@nl}{\def\PG@tc##1{\textcolor[rgb]{0.09,0.20,0.29}{##1}}}
\@namedef{PG@tok@nn}{\def\PG@tc##1{\textcolor[rgb]{0.09,0.20,0.29}{##1}}}
\@namedef{PG@tok@nx}{\def\PG@tc##1{\textcolor[rgb]{0.09,0.20,0.29}{##1}}}
\@namedef{PG@tok@nt}{\def\PG@tc##1{\textcolor[rgb]{0.09,0.20,0.29}{##1}}}
\@namedef{PG@tok@vc}{\def\PG@tc##1{\textcolor[rgb]{0.15,0.29,0.44}{##1}}}
\@namedef{PG@tok@vg}{\def\PG@tc##1{\textcolor[rgb]{0.15,0.29,0.44}{##1}}}
\@namedef{PG@tok@vi}{\def\PG@tc##1{\textcolor[rgb]{0.15,0.29,0.44}{##1}}}
\@namedef{PG@tok@vm}{\def\PG@tc##1{\textcolor[rgb]{0.15,0.29,0.44}{##1}}}
\@namedef{PG@tok@sa}{\def\PG@tc##1{\textcolor[rgb]{0.60,0.25,0.13}{##1}}}
\@namedef{PG@tok@sb}{\def\PG@tc##1{\textcolor[rgb]{0.60,0.25,0.13}{##1}}}
\@namedef{PG@tok@sc}{\def\PG@tc##1{\textcolor[rgb]{0.60,0.25,0.13}{##1}}}
\@namedef{PG@tok@dl}{\def\PG@tc##1{\textcolor[rgb]{0.60,0.25,0.13}{##1}}}
\@namedef{PG@tok@sd}{\def\PG@tc##1{\textcolor[rgb]{0.60,0.25,0.13}{##1}}}
\@namedef{PG@tok@s2}{\def\PG@tc##1{\textcolor[rgb]{0.60,0.25,0.13}{##1}}}
\@namedef{PG@tok@se}{\def\PG@tc##1{\textcolor[rgb]{0.60,0.25,0.13}{##1}}}
\@namedef{PG@tok@sh}{\def\PG@tc##1{\textcolor[rgb]{0.60,0.25,0.13}{##1}}}
\@namedef{PG@tok@si}{\def\PG@tc##1{\textcolor[rgb]{0.60,0.25,0.13}{##1}}}
\@namedef{PG@tok@sx}{\def\PG@tc##1{\textcolor[rgb]{0.60,0.25,0.13}{##1}}}
\@namedef{PG@tok@sr}{\def\PG@tc##1{\textcolor[rgb]{0.60,0.25,0.13}{##1}}}
\@namedef{PG@tok@s1}{\def\PG@tc##1{\textcolor[rgb]{0.60,0.25,0.13}{##1}}}
\@namedef{PG@tok@ss}{\def\PG@tc##1{\textcolor[rgb]{0.60,0.25,0.13}{##1}}}
\@namedef{PG@tok@mb}{\def\PG@tc##1{\textcolor[rgb]{0.46,0.35,0.07}{##1}}}
\@namedef{PG@tok@mf}{\def\PG@tc##1{\textcolor[rgb]{0.46,0.35,0.07}{##1}}}
\@namedef{PG@tok@mh}{\def\PG@tc##1{\textcolor[rgb]{0.46,0.35,0.07}{##1}}}
\@namedef{PG@tok@mi}{\def\PG@tc##1{\textcolor[rgb]{0.46,0.35,0.07}{##1}}}
\@namedef{PG@tok@il}{\def\PG@tc##1{\textcolor[rgb]{0.46,0.35,0.07}{##1}}}
\@namedef{PG@tok@mo}{\def\PG@tc##1{\textcolor[rgb]{0.46,0.35,0.07}{##1}}}
\@namedef{PG@tok@ow}{\def\PG@tc##1{\textcolor[rgb]{0.21,0.34,0.43}{##1}}}
\@namedef{PG@tok@pm}{\def\PG@tc##1{\textcolor[rgb]{0.09,0.20,0.29}{##1}}}
\@namedef{PG@tok@ch}{\def\PG@tc##1{\textcolor[rgb]{0.33,0.40,0.46}{##1}}}
\@namedef{PG@tok@cm}{\def\PG@tc##1{\textcolor[rgb]{0.33,0.40,0.46}{##1}}}
\@namedef{PG@tok@cp}{\def\PG@tc##1{\textcolor[rgb]{0.33,0.40,0.46}{##1}}}
\@namedef{PG@tok@cpf}{\def\PG@tc##1{\textcolor[rgb]{0.33,0.40,0.46}{##1}}}
\@namedef{PG@tok@c1}{\def\PG@tc##1{\textcolor[rgb]{0.33,0.40,0.46}{##1}}}
\@namedef{PG@tok@cs}{\def\PG@tc##1{\textcolor[rgb]{0.33,0.40,0.46}{##1}}}
\@namedef{PG@tok@gd}{\def\PG@tc##1{\textcolor[rgb]{0.09,0.20,0.29}{##1}}}
\@namedef{PG@tok@ge}{\def\PG@tc##1{\textcolor[rgb]{0.09,0.20,0.29}{##1}}}
\@namedef{PG@tok@gr}{\def\PG@tc##1{\textcolor[rgb]{0.09,0.20,0.29}{##1}}}
\@namedef{PG@tok@gh}{\def\PG@tc##1{\textcolor[rgb]{0.09,0.20,0.29}{##1}}}
\@namedef{PG@tok@gi}{\def\PG@tc##1{\textcolor[rgb]{0.09,0.20,0.29}{##1}}}
\@namedef{PG@tok@go}{\def\PG@tc##1{\textcolor[rgb]{0.09,0.20,0.29}{##1}}}
\@namedef{PG@tok@gp}{\def\PG@tc##1{\textcolor[rgb]{0.09,0.20,0.29}{##1}}}
\@namedef{PG@tok@gs}{\def\PG@tc##1{\textcolor[rgb]{0.09,0.20,0.29}{##1}}}
\@namedef{PG@tok@gu}{\def\PG@tc##1{\textcolor[rgb]{0.09,0.20,0.29}{##1}}}
\@namedef{PG@tok@ges}{\def\PG@tc##1{\textcolor[rgb]{0.09,0.20,0.29}{##1}}}
\@namedef{PG@tok@gt}{\def\PG@tc##1{\textcolor[rgb]{0.09,0.20,0.29}{##1}}}

\def\PGZbs{\char`\\}
\def\PGZus{\char`\_}
\def\PGZob{\char`\{}
\def\PGZcb{\char`\}}
\def\PGZca{\char`\^}
\def\PGZam{\char`\&}
\def\PGZlt{\char`\<}
\def\PGZgt{\char`\>}
\def\PGZsh{\char`\#}
\def\PGZpc{\char`\%}
\def\PGZdl{\char`\$}
\def\PGZhy{\char`\-}
\def\PGZsq{\char`\'}
\def\PGZdq{\char`\"}
\def\PGZti{\char`\~}
% for compatibility with earlier versions
\def\PGZat{@}
\def\PGZlb{[}
\def\PGZrb{]}
\makeatother


% Pygments emits slot 39 for apostrophes; use the actual ASCII glyphs.
\renewcommand{\PGZsq}{\textquotesingle}
\newcommand{\PGZgrave}{\textasciigrave}
\newcounter{guideListing}[chapter]
\renewcommand{\theguideListing}{\thechapter.\arabic{guideListing}}
\definecolor{codebackground}{HTML}{FAFBFC}
\definecolor{codeborder}{HTML}{C9D5DD}
\definecolor{codeheader}{HTML}{E9EFF3}
\renewcommand{\theFancyVerbLine}{\sffamily\scriptsize\color{muted}\arabic{FancyVerbLine}}
\newcommand{\CodeLanguage}[1]{\ifcsname code@lang@#1\endcsname\csname code@lang@#1\endcsname\else#1\fi}
\makeatletter
\def\code@lang@perl{Perl}
\def\code@lang@sql{SQL}
\def\code@lang@c{C}
\def\code@lang@bash{Shell}
\def\code@lang@lua{Redis server-side operation (Lua)}
\def\code@lang@perllua{Perl application (embedded Redis script)}
\definecolor{serverheader}{HTML}{F4E9DB}
\newcommand{\CodeHeaderColor}[1]{\ifcsname code@header@#1\endcsname\csname code@header@#1\endcsname\else codeheader\fi}
\def\code@header@lua{serverheader}
\def\code@lang@yaml{YAML}
\def\code@lang@text{Output / transcript}
\makeatother
\tcbset{guidecode/.style={enhanced,breakable,sharp corners,
  colback=codebackground,colframe=codeborder,colbacktitle=codeheader,coltitle=navy,
  boxrule=.5pt,left=24pt,right=7pt,top=6pt,bottom=6pt,
  lefttitle=8pt,righttitle=8pt,toptitle=5pt,bottomtitle=5pt,
  fonttitle=\sffamily\small\bfseries,before skip=9pt,after skip=9pt,
  pad at break=4pt,lines before break=5,
  overlay unbroken={\draw[codeborder] ([xshift=21pt]interior.north west)--([xshift=21pt]interior.south west);},
  overlay first={\draw[codeborder] ([xshift=21pt]interior.north west)--([xshift=21pt]interior.south west);},
  overlay middle={\draw[codeborder] ([xshift=21pt]interior.north west)--([xshift=21pt]interior.south west);},
  overlay last={\draw[codeborder] ([xshift=21pt]interior.north west)--([xshift=21pt]interior.south west);}}}
\newcommand{\CodeFile}[3]{%
  \par\needspace{6\baselineskip}\refstepcounter{guideListing}\label{listing:#2}%
  \begin{tcolorbox}[guidecode,
    colbacktitle=\CodeHeaderColor{#1},
    title={Listing \theguideListing\enspace\textperiodcentered\enspace\CodeLanguage{#1}\\[-1pt]\mdseries #3},
    title after break={Listing \theguideListing\enspace\textperiodcentered\enspace\CodeLanguage{#1}\enspace(continued)\\[-1pt]\mdseries #3}]%
  \VerbatimInput[commandchars=\\\{\},fontsize=\small,baselinestretch=1.10,numbers=left,numbersep=6pt,
    breaklines=true,breakanywhere=true,breakautoindent=true,breakindent=1em,
    breaksymbolleft={\tiny\color{muted}\ensuremath{\hookrightarrow}},
    breaksymbolright={},tabsize=4,obeytabs=true,showspaces=false,showtabs=false,
    formatcom=\color{navy}]{generated/code/#2.tex}%
  \end{tcolorbox}%
}


% ---------------------------------------------------------------- macros
% Inline code: verbatim-safe, with straight quotes and backticks so that
% copied Perl and shell snippets work (T1 typewriter fonts curl ' by default).
\usepackage{xstring}
\makeatletter
\DeclareRobustCommand{\code}[1]{\texttt{\expandafter\code@quotes\expandafter{\detokenize{#1}}}}
\newcommand{\code@quotes}[1]{\saveexpandmode\expandarg
  \StrSubstitute{#1}{'}{\noexpand\textquotesingle}[\code@tmp]%
  \StrSubstitute{\code@tmp}{`}{\noexpand\textasciigrave}[\code@tmp]%
  \StrSubstitute{\code@tmp}{::}{::\noexpand\allowbreak}[\code@tmp]%
  \restoreexpandmode\code@tmp}
\makeatother
% A single callout system: dedicated title, clear spacing, readable body.
\newtcolorbox{callout}[2]{enhanced,breakable,sharp corners,
  colback=#1!3!white,colframe=#1!45!white,boxrule=.45pt,
  leftrule=2.6pt,left=9pt,right=9pt,top=6pt,bottom=7pt,
  colbacktitle=#1!8!white,coltitle=#1,
  fonttitle=\sffamily\bfseries\small,title={#2},
  lefttitle=9pt,righttitle=9pt,toptitle=5pt,bottomtitle=5pt,
  title after break={#2\space(continued)},
  fontupper=\small\raggedright,
  before skip=11pt,after skip=11pt,pad at break=5pt}
\newcommand{\note}[1]{\begin{callout}{teal}{Engineering note}#1\end{callout}}
\newcommand{\warn}[1]{\begin{callout}{rust}{Correctness warning}#1\end{callout}}
\newcommand{\version}[1]{\begin{callout}{gold}{Version note}#1\end{callout}}
\newcommand{\tested}[1]{\begin{callout}{navy}{Historical laboratory evidence}#1\end{callout}}
\newcommand{\tip}[1]{\begin{callout}{leaf}{Practical tip}#1\end{callout}}
% Strategy card: key, name, plain terms, use when, avoid when, Perl tools
\newtcolorbox{strategybox}[1]{enhanced,colback=white,colframe=navy,boxrule=0.6pt,arc=2pt,
  colbacktitle=navy,coltitle=white,fonttitle=\sffamily\bfseries\small,title={#1},
  fontupper=\small,left=7pt,right=7pt,top=5pt,bottom=5pt,before skip=10pt,after skip=10pt}
\newcommand{\strategy}[6]{\begin{strategybox}{STRATEGY \textperiodcentered\ #2}\phantomsection\label{str:#1}%
\textbf{\sffamily\color{teal}In plain terms.} #3\par\smallskip
\textbf{\sffamily\color{navy}Use it when:} #4\par\smallskip
\textbf{\sffamily\color{rust}Avoid it when:} #5\par\smallskip
\textbf{\sffamily\color{navy}In Perl:} #6\end{strategybox}}
\newcommand{\refdoc}[2]{\par\hangindent=1.5em\hangafter=1\textbf{#1.} \url{#2}\par}


\newcommand{\chapterbrief}[1]{\par\noindent{\sffamily\color{teal}\bfseries Reading objective}\par
{\small\itshape #1}\par\medskip}
\newcommand{\partstart}[3]{%
  \clearpage\refstepcounter{part}\renewcommand{\currentpart}{#1}%
  \phantomsection\addcontentsline{toc}{part}{\thepart\quad #2}%
  \markboth{}{}\thispagestyle{empty}%
  \vspace*{.9in}{\sffamily\bfseries\Large\color{teal}PART \thepart\par}
  \vspace{10pt}{\color{teal}\rule{\linewidth}{2pt}}\par
  \vspace{18pt}{\sffamily\bfseries\Huge\color{navy}\raggedright\hyphenpenalty=10000\exhyphenpenalty=10000 #2\par}
  \vspace{20pt}{\large #3\par}\vfill
  {\small\sffamily\color{muted}Modern Perl Concurrency \textperiodcentered{} Edition 2.1.1}
  \clearpage}
\begin{document}
% ================================================================ title page
\begin{titlepage}
\color{navy}\vspace*{0.45in}
{\sffamily\bfseries\large ENGINEERING HANDBOOK\par}\vspace{0.25in}
{\color{teal}\rule{\linewidth}{3pt}}\vspace{0.4in}
{\Huge\sffamily\bfseries Modern Perl\\[8pt]Concurrency\par}
\vspace{0.28in}{\Large\sffamily Synchronization, Distributed Coordination\\
and Application Architecture\par}
\medskip{\normalsize\sffamily\color{muted}Threads, processes, event loops, MySQL, DBIx::Class, Redis, Valkey, Catalyst and Memcached\par}
\vspace{0.45in}
\begin{center}\scalebox{1.12}{\begin{tikzpicture}[x=1cm,y=1cm,font=\sffamily\small]
% workers
\foreach \i/\y in {1/1.5,2/0.5,3/-0.5,4/-1.5} {
  \node[box,minimum width=2.3cm] (w\i) at (0,\y) {worker \i};}
\node[lbl] at (0,-2.3) {threads, processes, pods};
% coordination layer
\node[draw=teal,fill=teal!8,rounded corners=3pt,minimum width=5.2cm,minimum height=4.0cm] (co) at (5.3,0) {};
\node[font=\sffamily\small\bfseries,text=teal] at (5.3,1.6) {coordination};
\foreach \t/\y in {{lock() / Future::Mutex}/0.85,{flock / semaphores}/0.15,{transactions / SKIP LOCKED}/-0.55,{leases + fencing tokens}/-1.25} {
  \node[box,fill=white,minimum width=4.6cm,minimum height=5.5mm,font=\sffamily\footnotesize] at (5.3,\y) {\t};}
\foreach \i in {1,2,3,4} \draw[arr] (w\i.east) -- (co.west |- w\i.east);
% resources
\node[box,minimum width=2.9cm,minimum height=1.1cm] (r1) at (10.7,0.9) {shared state\\\footnotesize rows, files, keys};
\node[okbox,minimum width=2.9cm,minimum height=1.1cm] (r2) at (10.7,-0.9) {fenced resource\\\footnotesize rejects stale writers};
\draw[arr] (co.east |- r1) -- (r1.west); \draw[arr] (co.east |- r2) -- (r2.west);
\end{tikzpicture}
}\end{center}
\vfill
{\LARGE\sffamily\bfseries Scott Muller\par}
\smallskip
{\sffamily Software Engineer, Matsch Systems\\
Owner \& Independent Consultant, 21h Systems\\
\href{mailto:scott@21h.io}{scott@21h.io}\par}
\vspace{0.3in}
{\sffamily November 2026\quad /\quad Edition 2.1.1\quad /\quad Technical reference and self-verifying laboratory\par}
\end{titlepage}

% ---------------------------------------------------------------- copyright page
\thispagestyle{empty}
\vspace*{\fill}
{\small\sffamily
\textbf{Modern Perl Concurrency}\\
Synchronization, Distributed Coordination and Application Architecture\\
Engineering Handbook, Edition 2.1.1 (November 2026)\par\medskip
\textbf{Scott Muller}\par
Software Engineer, Matsch Systems\\
Owner \& Independent Consultant, 21h Systems\par
Contact: \href{mailto:scott@21h.io}{scott@21h.io}.\par
\textit{This work is an independent technical publication. The views and recommendations expressed herein are solely those of the author and do not necessarily reflect those of his employer or affiliated organizations. Professional affiliations do not imply institutional endorsement.}\par
\copyright\ 2026 Scott Muller.\par\smallskip
Author-owned prose and illustrations: \href{https://creativecommons.org/licenses/by/4.0/}{CC BY 4.0}. Author-owned example code and build tooling: MIT. See the source package's license notices; third-party material retains its own terms.\par\bigskip
\color{muted}
Documentation baseline: Perl 5.44; MySQL 8.4 LTS and 9.7 LTS; Redis 7.2--8.4; Valkey 8.1--9.0; Catalyst 5.90; DBIx::Class 0.0828.\\
Current release verification: Appendix~\ref{sec:current-verification}; historical results are retained separately.\par\medskip
The original laboratory verification record is retained in Appendix~\ref{ch:refs}. The application-architecture laboratories are separately identified in the verification chapter. Current release logs distinguish executed tests from optional skips; full publication requires the configured laboratories.
Examples that create, flush or delete data must be run only against disposable, isolated services.\par}
\clearpage

% ================================================================ front matter
\chapter*{Reading this handbook}
\addcontentsline{toc}{chapter}{Reading this handbook}
This guide distinguishes concurrency (multiple activities in progress) from parallelism (multiple activities executing simultaneously), and synchronization from durability. It is written for engineers maintaining production Perl systems: long-running daemons, batch processors, web services and database-backed workers. It favors established primitives, explicit ownership and constraints enforced by the protected resource over elaborate locking frameworks.

The source package includes executable Perl laboratories, a native C futex demonstration, a complete Catalyst application, configuration, SQL and protocol transcripts. Complete programs check their invariants and exit non-zero on failure; short excerpts illustrate individual operations. Perl is the application language throughout. SQL and embedded Redis Lua express operations executed by the corresponding server, not alternative application implementations. The original integration results are retained in Appendix~\ref{ch:refs}. The reader/writer lock and barrier are educational implementations, not audited libraries.

\section*{How to read this handbook}
\textbf{New to concurrency?} Read Part I, then read each chapter's \emph{strategy cards} and figures before the code. Each strategy card starts with a plain-language explanation (\emph{In plain terms}) and says when to use and avoid the strategy. The glossary (Appendix~\ref{ch:glossary}) defines every term in one or two sentences. Run the examples: each one prints what it proved, and fails loudly if the guarantee breaks.

\textbf{Experienced engineers} can go straight to the correctness warnings, version notes and verified results, which record exact behaviors, error codes and the measurements taken on the laboratory host. Appendix~\ref{ch:refs} lists what was and was not tested.

The examples use \code{use strict; use warnings;} so that they run on the older Perls still common in production; on Perl 5.36 or later, \code{use v5.36} enables both, plus subroutine signatures. Example 24 shows the 5.36+ style.

\section*{Reading paths}
\begin{longtable}{L{.28\linewidth}L{.63\linewidth}}
\toprule Reader & Suggested path\\ \midrule\endhead
New Perl developer & Parts I and II, then the worked examples in Part III.\\
Application developer & Parts III and V, especially transaction ownership and request coordination.\\
Data-processing engineer & Parts II, III and V; follow the capstone from admission to delivery.\\
Systems engineer & Parts IV and VI, including failure boundaries, deployment and verification.\\
\bottomrule
\end{longtable}
Use the subject index, API entries and example catalog for problem-oriented lookup.
The appendices retain the complete application sources and native implementation details.

\section*{Conventions}
\begin{longtable}{L{.28\linewidth}L{.63\linewidth}}
\toprule Label & Meaning \\ \midrule\endhead
\textcolor{teal}{\textbf{Engineering note}} & Design guidance and operational practice.\\
\textcolor{rust}{\textbf{Correctness warning}} & A failure mode that can corrupt state or violate exclusion.\\
\textcolor{gold}{\textbf{Version note}} & Behavior dependent on a server, module or Perl release.\\
\textcolor{navy}{\textbf{Historical evidence}} & What earlier laboratories actually verified.\\
\textcolor{leaf}{\textbf{Practical tip}} & A useful operational practice or shortcut.\\
\bottomrule\end{longtable}

\strategy{example}{How a strategy card reads}{A one-paragraph explanation with an everyday analogy, for readers meeting the idea for the first time.}
{the situations the strategy is designed for.}
{the situations where it fails or a different strategy is better.}
{the modules, functions and laboratory examples that implement it.}

\section*{Recommended defaults}
\begin{itemize}
\item Let the protected resource enforce correctness: a unique key, a guarded \code{UPDATE} or a fencing check beats any lock.
\item Transfer work through queues instead of sharing mutable state; make every external effect idempotent.
\item Use lexical \code{lock} for small in-process invariants, processes for CPU-heavy isolated work and an event loop for high-concurrency I/O.
\item Use InnoDB transactions for persistent shared state; use advisory locks only where the resource does not fit naturally into a transaction.
\item Keep web sessions small and immutable in practice; put state that concurrent requests change into atomic Redis structures or the database.
\item Treat every Redis, Valkey or Memcached lock as a lease that can be lost. Use it for efficiency; pair it with destination-enforced fencing when a violation would cause harm.
\end{itemize}

\section*{Locking strategies at a glance}
\begingroup\small\renewcommand{\arraystretch}{1.03}
\begin{longtable}{L{.30\linewidth}L{.24\linewidth}L{.24\linewidth}L{.10\linewidth}}\toprule Strategy & Scope & Typical use & Section \\ \midrule\endhead
Mutex (\code{lock}) & threads & shared variables & \ref{str:mutex} \\
Atomic operations, futexes & native code & how locks work & \ref{str:futex} \\
Snapshot publication & threads & read-mostly config & \ref{str:snapshot} \\
Counting semaphore & threads, processes & cap concurrency & \ref{str:semaphore} \\
Condition variable & threads & wait for a state & \ref{str:condvar} \\
Barrier & threads & phased computation & \ref{str:barrier} \\
Work queue & threads, processes & worker pools & \ref{str:queue} \\
Async mutex & event loop & state across \code{await} & \ref{str:asyncmutex} \\
Reader/writer lock & threads, processes & read-heavy data & \ref{str:rwlock} \\
Lock ordering, striping & any & multi-lock operations & \ref{str:ordering} \\
Release on every path & any & cleanup discipline & \ref{str:release} \\
File lock (\code{flock}) & processes, one host & cron, single writer & \ref{str:flock} \\
System~V semaphore & processes, one host & shared permits & \ref{str:sysvsem} \\
Optimistic concurrency & database & forms, APIs & \ref{str:optimistic} \\
Pessimistic row locks & database & read-then-update & \ref{str:rowlocks} \\
Named advisory lock & database server & cross-host jobs & \ref{str:advisory} \\
Database work queue & database & durable jobs & \ref{str:dbqueue} \\
Redis/Valkey lease & cluster & avoid duplicate work & \ref{str:redislease} \\
Fencing token & cluster & reject stale writers & \ref{str:fencing} \\
Distributed semaphore & cluster & shared quotas & \ref{str:redissem} \\
Redis Streams & cluster & job distribution & \ref{str:streams} \\
Memcached add/CAS lease & cluster & cache refresh & \ref{str:memcas} \\
Cache stampede control & cluster & hot cache keys & \ref{str:stampede} \\
Per-session coordination & web sessions & parallel requests & \ref{str:session} \\
Idempotency key & web requests & payments, orders & \ref{str:idempotency} \\
Rate limiting & web requests & login, quotas & \ref{str:ratelimit} \\
Hashed API keys & API clients & integrations, scripts & \ref{str:apikey} \\
Short-lived token with revocation list & API clients & JWT access tokens & \ref{str:jwtrevocation} \\
Refresh-token rotation & API clients & mobile apps, SPAs & \ref{str:refreshrotation} \\
Signed request with nonce & partner APIs & webhooks, B2B calls & \ref{str:signedrequest} \\
Consensus coordinator & cluster & correctness-critical & \ref{str:consensus} \\
Worker-owned resources & process & client lifetime & \ref{str:app-objects}\\
Service transaction & database & atomic operation & \ref{str:app-transactions}\\
Durable job claim & database & recoverable work & \ref{str:app-pipelines}\\
Transactional outbox & services & durable delivery & \ref{str:app-outbox}\\
Invariant migration & application & mixed writers & \ref{str:app-migration}\\
Deterministic race test & tests & reproduce failure & \ref{str:app-testing}\\
Bounded admission & deployment & finite capacity & \ref{str:app-operations}\\
\bottomrule\end{longtable}

\endgroup
\needspace{4\baselineskip}\section*{Coverage and verification}
The laboratories demonstrate successful coordination and deliberately reproduced races. Appendix~\ref{ch:refs} identifies the original integration environment and the limits of subsequent verification. The source archive contains the complete examples, build instructions and revision history.

\clearpage
\section*{Continuously verified technical publishing}
This handbook and its executable laboratories are maintained as a unified engineering project. Canonical example sources feed both automated tests and generated syntax-highlighted listings. GitHub Actions exercises the laboratories against disposable services, records benchmark observations, compiles the LaTeX documents and checks the resulting PDFs before publication.

\begin{center}
\begin{adjustbox}{max width=\linewidth}
\begin{tikzpicture}[node distance=6mm and 8mm]
\node[box,text width=48mm] (repo) {Versioned repository};
\node[box,text width=43mm,below left=of repo] (src) {Canonical examples};
\node[box,text width=43mm,below right=of repo] (tex) {LaTeX manuscript};
\node[box,text width=43mm,below=of src] (test) {Formatting and laboratory tests};
\node[box,text width=43mm,below=of tex] (listing) {Generated listings};
\node[box,text width=43mm,below=of test] (evidence) {Verification and benchmarks};
\node[box,text width=43mm,below=of listing] (pdf) {Compilation and PDF preflight};
\node[okbox,text width=70mm,below=12mm of $(evidence.south)!0.5!(pdf.south)$] (release) {Versioned release and documentation portal};
\draw[arr] (repo) -- (src);
\draw[arr] (repo) -- (tex);
\draw[arr] (src) -- (test);
\draw[arr] (test) -- (listing);
\draw[arr] (test) -- (evidence);
\draw[arr] (tex) -- (listing);
\draw[arr] (listing) -- (pdf);
\draw[arr] (evidence) -- (pdf);
\draw[arr] (evidence) -- (release);
\draw[arr] (pdf) -- (release);
\end{tikzpicture}
\end{adjustbox}
\captionof{figure}{The handbook as a continuously verified release artifact.}
\end{center}

Release artifacts include the source snapshot, verification manifest, benchmark results, build log and checksums. These provide inspectable evidence of the tested behavior, not a formal proof of concurrency correctness or a substitute for independent review. See Appendix~\ref{sec:current-verification} and the \href{https://github.com/interrupt21h/Modern-Perl-Concurrency/actions}{public build history}.

\clearpage
\begingroup\footnotesize\setlength{\parskip}{0pt}\tableofcontents\endgroup

% ================================================================ chapter 1

\partstart{I / Foundations}{Foundations and mental models}{Begin with the invariant and owner of state before selecting an implementation.}
\chapter{Choosing a concurrency architecture}
\chapterbrief{Select an execution model and identify the owner of shared state.}
\par\smallskip\noindent\textbf{Abbreviations used in this chapter:} CPU --- central processing unit; DBI --- Database Interface; IPC --- interprocess communication; MCE --- Many-Core Engine; DBD --- Database Driver.\par\smallskip

\index{Execution models}
\index{Ownership}
\index{Invariants}

\section{Execution models}
A Perl interpreter thread (ithread) runs its own interpreter. Creating one clones much of the parent's interpreter state; ordinary variables are private unless explicitly shared. Ithreads may run in parallel on separate cores; there is no global interpreter lock. The cost of interpreter cloning, object duplication and synchronization makes thread-per-request designs unattractive, and the \code{threads} documentation states that interpreter-based threads are officially discouraged for new code [1,27]. Ithreads remain appropriate for maintaining existing threaded applications and for small, fixed worker pools.

A forked worker has its own address space, initially sharing pages copy-on-write on Unix-like systems. A process crash can be isolated from its supervisor, although inherited descriptors and database connections demand care. An event loop instead multiplexes nonblocking I/O in one interpreter. It avoids memory-sharing races, but logical races still occur between callbacks and across every \code{await}: state read before a suspension point may be stale after it (Section~\ref{sec:async}). Blocking calls stall the whole loop unless dispatched to a worker.
\fig{execmodels}{The three execution models available to Perl programs, and how each shares and coordinates state.}
\begingroup\small\renewcommand{\arraystretch}{1.02}
\begin{longtable}{L{.22\linewidth}L{.36\linewidth}L{.34\linewidth}}
\toprule Model & Strength & Trade-off \\ \midrule\endhead
Ithreads & Existing shared-state applications, fixed worker pools & Interpreter memory, XS compatibility, complex lifecycle; officially discouraged for new code \\
Fork / MCE & CPU-intensive batch work, fault isolation & Serialization, inherited handles, process supervision \\
Event loop & Many network connections, timers, fan-out to services & Nonblocking clients required; CPU work must be offloaded \\
Database workers & Durable coordination across hosts & Transaction design, latency, deadlock recovery \\
Redis Streams workers & Low-latency distribution across hosts with acknowledgment & At-least-once delivery; durability depends on server configuration \\
\bottomrule
\end{longtable}
\endgroup

\section{Practical module selection}
Core facilities include \code{threads}, \code{threads::shared}, \code{Thread::Queue}, \code{Thread::Semaphore}, \code{IPC::Semaphore}, pipes and \code{flock}. MCE offers persistent worker pools and chunked parallel computation [9]; \code{MCE::Shared} provides managed shared objects across workers [16]. \code{Parallel::ForkManager} is useful for bounded process fan-out. For event-driven code, \code{IO::Async} supplies the loop, timers and process management [10]; \code{Future} represents eventual results; \code{Future::AsyncAwait} provides \code{async}/\code{await} syntax [29]; and \code{Future::Mutex} serializes asynchronous critical sections [30]. \code{Mojo::IOLoop} with \code{Mojo::Promise} is a coherent alternative within the Mojolicious ecosystem. DBI with \code{DBD::mysql} or \code{DBD::MariaDB} provides the database layer (Section~\ref{sec:dbi}); Redis client selection is covered in Section~\ref{sec:redis-clients}.

Prefer a pinned, tested dependency set (for example a \code{cpanfile} with \code{carton} or a vendored \code{local::lib}) to assumptions based on a module's age.
\CodeFile{bash}{snippets/guide-01.sh}{Inspect the runtime}

\clearpage\section{Ownership and scope}
Choose the scope of the invariant before choosing a primitive. A mutex in one process cannot prevent a separate cron job from updating a row. A database advisory lock does not serialize an uncooperative client. A queue delivers items but does not itself make external effects exactly-once. Document the resource identity, participants, lifecycle, failure behavior and acquisition order for each synchronization domain.
\begin{figure}[H]\centering\begin{adjustbox}{max width=\linewidth}\begin{tikzpicture}[x=1cm,y=1cm]
\def\step{1.02}\def\top{5.45}
\foreach \i/\shade in {0/14,1/0,2/14,3/0,4/14} {
  \fill[paper!\shade!white] (0.3*\i,0.12*\i) rectangle (14.6-0.3*\i,\top-\step*\i);
  \draw[navy!70,rounded corners=2pt] (0.3*\i,0.12*\i) rectangle (14.6-0.3*\i,\top-\step*\i);
}
\foreach \i/\who/\how in {
  0/{Hosts without a shared database}/{Redis or Valkey lease for efficiency; consensus store plus fencing for correctness},
  1/{Hosts sharing a MySQL database}/{InnoDB transactions and constraints, SKIP LOCKED queues, GET\_LOCK},
  2/{Processes on one host}/{flock, System V semaphores, pipes, a single writer process},
  3/{Threads in one interpreter process}/{lock(), cond\_wait, Thread::Queue, Thread::Semaphore},
  4/{Tasks in one event loop}/{Future::Mutex, bounded fmap\_void, single-owner state}} {
  \node[anchor=west,font=\sffamily\small\bfseries,text=navy] at (0.3*\i+0.2,\top-\step*\i-0.3) {\who};
  \node[anchor=west,font=\sffamily\footnotesize,text=muted] at (0.3*\i+0.2,\top-\step*\i-0.68) {\how};
}
\draw[arr] (14.95,1.0) -- node[lbl,right,text width=1.7cm,align=left]{wider\\scope:\\more\\failure\\modes and\\latency} (14.95,5.2);
\end{tikzpicture}
\end{adjustbox}\caption{Coordination scopes. Each wider scope must assume every failure mode of the narrower ones, plus its own.}\label{fig:scope}\end{figure}

\tip{Keep a lock register: one table in the design notes (or one module in the code) that lists every lock, what it protects, its scope, its timeout and its acquisition order. Most deadlocks and orphaned leases come from locks nobody knew existed.}

\clearpage\section{Selecting a coordination mechanism}\label{sec:select}
Two questions decide the mechanism: \emph{who participates}, and \emph{what a violation costs}. A lock taken for \emph{efficiency} prevents duplicated work; an occasional violation wastes resources but corrupts nothing. A lock taken for \emph{correctness} protects an invariant; a violation corrupts state, double-charges a customer or emits a duplicate filing [43]. For correctness across hosts, prefer destination-enforced constraints (unique keys, conditional updates, version columns); use locks only to reduce conflicts.
\begin{figure}[H]\centering\begin{adjustbox}{max width=.79\linewidth}\begin{tikzpicture}[node distance=7mm and 4mm]
\node[decision,aspect=3.6] (q1) {Can the resource enforce the rule itself?\\(unique key, guarded UPDATE, version column)};
\node[okbox,right=6mm of q1,text width=4.0cm] (c) {Use the constraint. Add locks only to reduce conflicts, never as the sole guard.};
\draw[arr] (q1) -- node[lbl,above]{yes} (c);
\node[box,below=8mm of q1] (q2) {Who participates?};
\draw[arr] (q1) -- node[lbl,right]{no} (q2);
\node[box,below left=10mm and 22mm of q2,text width=2.9cm] (p1) {One process\\\footnotesize lock(), Thread::Queue, Future::Mutex};
\node[box,right=4mm of p1,text width=2.6cm] (p2) {One host\\\footnotesize flock, SysV semaphore};
\node[box,right=4mm of p2,text width=3.1cm] (p3) {Shared MySQL\\\footnotesize transaction, SKIP LOCKED, GET\_LOCK};
\node[box,right=4mm of p3,text width=2.6cm] (p4) {Cluster, no shared database};
\foreach \n in {p1,p2,p3,p4} \draw[arr] (q2.south) -- (\n.north);
\node[decision,below=8mm of p4] (q3) {Is a violation\\harmful?};
\draw[arr] (p4) -- (q3);
\node[okbox,below left=8mm and -4mm of q3,text width=3.0cm] (e) {Redis/Valkey lease\\\footnotesize efficiency lock};
\node[badbox,below right=8mm and -4mm of q3,text width=4.3cm] (f) {Consensus coordinator\\plus a fencing token\\checked by the resource};
\draw[arr] (q3) -| node[lbl,pos=0.25,above]{no} (e);
\draw[badarr] (q3) -| node[lbl,pos=0.25,above]{yes} (f);
\end{tikzpicture}
\end{adjustbox}\caption{Selecting a coordination mechanism. Start from the resource, not from the lock.}\label{fig:decide}\end{figure}
\begin{longtable}{L{.20\linewidth}L{.34\linewidth}L{.38\linewidth}}
\toprule Participants & Efficiency (duplicates tolerable) & Correctness (violation corrupts state) \\ \midrule\endhead
Tasks in one event loop & \code{Future::Mutex}; bounded \code{fmap} & Same, plus invariant tests across \code{await} points \\
Threads in one interpreter & \code{lock}, \code{Thread::Queue} & Same; keep invariants inside one critical section \\
Processes on one host & \code{flock}, System~V semaphore & \code{flock} with atomic file replacement, or a database transaction \\
Hosts sharing MySQL & \code{GET_LOCK}, \code{SKIP LOCKED} queue & Transactions with constraints; attempt-fenced completion (Section~\ref{sec:dbqueue}) \\
Hosts without a shared database & Redis/Valkey lease; Memcached \code{add} for cache refresh & Consensus-backed coordinator plus destination-enforced fencing (Section~\ref{sec:authorities}) \\
\bottomrule
\end{longtable}


% ================================================================ chapter 2

\chapter{Synchronization primitives and semantics}
\chapterbrief{Distinguish exclusion, capacity, signalling and ordering.}
\index{Mutex}
\index{Semaphore}
\index{Condition variable}
\index{Barrier}

Each primitive below has a strategy card with a plain-language explanation and a diagram; the last column gives its section.
\begin{longtable}{L{.19\linewidth}L{.34\linewidth}L{.29\linewidth}L{.07\linewidth}}
\toprule Primitive & Meaning & Typical use & See \\ \midrule\endhead
Mutex & Exclusive ownership; one critical section at a time & Cache or compound state update & \ref{str:mutex} \\
Semaphore & Count of available permits; acquire consumes, release restores & Bound concurrent connections & \ref{str:semaphore} \\
Condition variable & Wait for a predicate while temporarily releasing a mutex & Startup readiness, queue internals & \ref{str:condvar} \\
Queue & Transfer messages between producers and consumers & Worker pools and ownership transfer & \ref{str:queue} \\
Reader/writer lock & Concurrent readers or a single writer & Read-heavy mutable configuration & \ref{str:rwlock} \\
Barrier & Wait until all participants arrive at a phase & Parallel iterative computation & \ref{str:barrier} \\
Atomic operation & Indivisible update or compare-and-exchange with memory ordering & Native counters or lock internals & \ref{str:futex} \\
Futex & Linux kernel wait/wake on a shared 32-bit word & Implementation of native synchronization & \ref{str:futex} \\
Spinlock & Repeatedly poll an atomic flag & Very short native critical sections & \ref{str:futex} \\
Lease & Time-limited claim that can expire while the holder still runs & Distributed coordination & \ref{str:redislease} \\
Fencing token & Monotonic generation checked by the protected resource & Rejecting stale lease holders & \ref{str:fencing} \\
\bottomrule
\end{longtable}
\section{Mutex versus semaphore}
\strategy{mutex}{Mutex (lock)}{A mutex is the key to a single-occupancy room. Whoever holds the key may change what is inside; everyone else waits at the door. When the holder leaves the room (the scope ends), the key is handed on automatically.}
{several threads update the same shared variables and the update takes more than one step (read, change, write).}
{the work inside would be slow (network calls, sleeps), or when the participants are separate processes or hosts, which cannot see a Perl lock at all.}
{\code{lock($var)} from \code{threads::shared}; examples 01 and 26.}
A plain \code{++$count} is three steps (read, add, write), and two threads can interleave them (Figure~\ref{fig:lostupdate}). Example 01 shows the cure.
\fig{lostupdate}{A lost update and its prevention. Both threads in (a) read 41 before either writes.}
A mutex protects a resource invariant and normally has ownership semantics. A counting semaphore expresses capacity; a binary semaphore can emulate exclusivity but typically has no ownership enforcement and can be over-released. A semaphore initialized to three does not guarantee that three threads will run; it only bounds concurrent permit holders if all participants follow the protocol. Neither primitive intrinsically guarantees fairness. A rate limiter controls operations per unit time, which is different from a concurrency limit.
\section{Conditions, barriers and queues}
A condition notification is not stored as a durable event. Waiters must check the predicate with the protecting lock held and recheck after waking, because wakeups can be spurious or consumed by another waiter. A barrier is a collective rendezvous, not mutual exclusion; one failed participant can strand everyone. A queue can encapsulate synchronization while conveying ownership. Bounded queues apply backpressure, but cyclic dependencies between bounded queues can deadlock.
\section{Native mechanisms}
Atomic instructions and futexes underlie many operating-system locks.
Perl applications ordinarily use higher-level primitives; Appendix~\ref{app:native}
explains the implementation, memory ordering and C laboratory.

\partstart{II / Local synchronization}{Perl and operating-system synchronization}{Apply primitives within interpreters, between threads and across process boundaries.}
\chapter{Perl interpreter threads}
\chapterbrief{Understand explicitly shared state and scoped locking.}
\index{Threads}
\index{Shared memory}
\index{API!threads::shared}
\index{API!lock}

\section{Shared memory and lexical mutexes}
\index{API!threads::shared}\index{API!lock}
Load \code{threads} before \code{threads::shared}. Declare shared variables before creating workers. A shared container may hold scalars and references to shareable values; use \code{shared_clone} for nested structures and verify object compatibility. Lock the same stable shared variable in every participant; a dedicated, never-reassigned lock variable avoids confusion about what is protected. Locking a container does not recursively lock nested values. A plain increment is a read-modify-write and is not atomic [2].
\fig{ithreads}{The ithreads memory model: separate interpreters, one explicitly shared store.}
\CodeFile{perl}{examples/01-mutex.pl}{\nolinkurl{01-mutex.pl}}
The lock is recursive for the owning thread and released when its enclosing scope exits, including exit by exception. Holding a lock across a slow network operation expands contention and risks lock inversion; keep the critical section focused on the invariant. There is no timed acquisition for \code{lock}; use \code{Thread::Semaphore::down_timed} or a condition variable with a deadline when acquisition must be bounded.
Shared nested data needs a publication discipline. Example 22 builds each configuration version privately, clones it into shared memory \emph{outside} the lock, then swaps a single shared reference under the lock. A published structure is never modified again, so a reader that grabs the current reference sees a self-consistent version for as long as it holds it, with no further locking. A reader that wants to edit its view takes a private deep copy.
\strategy{snapshot}{Immutable snapshot publication}{Instead of editing a notice while people are reading it, print a new notice and swap it onto the board in one motion. Readers holding the old copy keep a complete, consistent version; the next reader picks up the new one.}
{data is read often and changed rarely: configuration, routing tables, feature flags, reference data.}
{many writers change the data frequently, or readers must always see the very latest value at the instant it changes.}
{build privately, \code{shared_clone} outside the lock, swap one shared reference under \code{lock}; example 22.}
\fig{snapshot}{Snapshot publication: the expensive work happens outside the lock, and only a reference swap is protected.}
\CodeFile{perl}{examples/22-config-snapshot.pl}{\nolinkurl{22-config-snapshot.pl}}
\tested{Four readers took 1{,}880 snapshots while 499 new versions were published; every snapshot was internally consistent and versions never went backwards.}
A shallow copy (\texttt{\{ \%\$config \}}) is not isolation for nested state: the inner hashes are still shared. Do not share DBI handles, sockets with ambiguous ownership, or XS-backed objects without documented thread support.
\section{lock() semantics in detail}
\code{lock} is easy to misuse because several plausible calls either die or do less than they appear to. Every row below was verified on Perl 5.38.
\begin{longtable}{L{.36\linewidth}L{.56\linewidth}}\toprule Call or situation & Behavior \\ \midrule\endhead
\code{lock($s)}, \code{lock(@a)}, \texttt{lock(\%h)} on shared variables & Advisory lock, recursive for the holder, released at scope exit (including exit by exception). \\
\code{lock($h{key})}, \code{lock($a[0])} & Dies: \emph{lock can only be used on shared values}. Elements cannot be locked; lock the container or a dedicated lock variable. \\
\code{lock($plain)} on an unshared variable & Dies with the same message; it is not a silent no-op. \\
\code{lock($ref)} where \code{$ref} holds a shared reference & Follows one level of reference and locks the referent: another thread's \texttt{lock(\%\$ref)} waits. \\
Another thread writes \code{$h{k}} without locking while you hold \texttt{lock(\%h)} & Succeeds. Locks constrain only code that also locks. \\
\code{cond_wait($v)} without holding \code{lock($v)} & Dies: \emph{You need a lock before you can cond\_wait}. \\
\code{lock} inside a loop body & Released at the end of each iteration; to hold a variable number of locks, recurse (example 23). \\
Bounded wait & Not available; use \code{Thread::Semaphore->down_timed} or \code{cond_timedwait} with a deadline. \\
\bottomrule\end{longtable}
\section{Thread lifecycle}
Use a fixed pool, join each worker and inspect failures. A thread exception does not propagate into the parent: \code{join} returns an empty list and \code{$thread->error()} holds the exception. Every example in this edition returns a true value from the thread and fails loudly when \code{join} does not deliver it. Design application-level result envelopes for richer outcomes. Detached threads are appropriate only when their lifetime and failure handling are intentionally independent. Graceful shutdown stops admission, drains or cancels queued work, joins workers and closes resources. Do not kill threads while they hold locks, and keep the main interpreter alive until shutdown completes.

% ================================================================ chapter 4

\chapter{Semaphores and condition variables}
\chapterbrief{Bound concurrency and wait on predicates.}
\par\smallskip\noindent\textbf{Abbreviations used in this chapter:} API --- application programming interface; CPU --- central processing unit; IPC --- interprocess communication.\par\smallskip

\index{Semaphore}
\index{Condition variable}
\index{API!Thread::Semaphore}
\index{API!cond_wait@\texttt{cond\_wait}}

\section{Counting semaphores}
\index{API!Thread::Semaphore}\index{API!down}\index{API!up}
\strategy{semaphore}{Counting semaphore}{A car park with a fixed number of spaces and a counter at the gate. Each car entering takes a ticket (\code{down}); each car leaving returns one (\code{up}). When the counter reaches zero, the next car waits at the barrier.}
{you need to cap how many workers use something at once: database connections, API calls, CPU-heavy jobs.}
{you need exclusive ownership of a specific resource (use a mutex), or you need to limit requests per second (that is a rate limit, Section~\ref{str:ratelimit}).}
{\code{Thread::Semaphore} (threads, example 02); \code{IPC::Semaphore} (processes on one host, example 07); a Redis sorted set across hosts (example 16).}
\fig{semaphore}{A counting semaphore with three permits: the fourth worker waits until a permit is returned.}
\code{Thread::Semaphore} supports blocking, immediate and timed acquisition. \code{down_timed} takes a relative timeout in seconds; contrast this with the absolute epoch deadline used by \code{cond_timedwait}. A successful \code{down} must have exactly one matching \code{up} on every path, including exceptions. The lab tracks peak usage under a separate shared mutex [3].
\CodeFile{perl}{examples/02-semaphore.pl}{\nolinkurl{02-semaphore.pl}}
A permit is not a database connection. Create a private connection per worker or use a separately designed pool. Do not use \code{down_force} as a normal acquisition operation: it may drive the permit count negative. When acquiring multiple permits, request them together rather than one at a time, which risks partial-capacity deadlock.
\section{Condition-variable rendezvous}
\index{API!cond_wait@\texttt{cond\_wait}}\index{API!cond_signal@\texttt{cond\_signal}}\index{API!cond_broadcast@\texttt{cond\_broadcast}}
\strategy{condvar}{Condition variable}{A waiting room with a bell. You check the notice board (the predicate) while holding the room key; if your number is not up, you hand back the key and doze until the bell rings, then take the key again and re-read the board.}
{one thread must wait until another has prepared something: startup readiness, ``queue not empty'', ``all results in''.}
{a queue already expresses the hand-off (prefer \code{Thread::Queue}), or the waiting crosses process boundaries.}
{\code{cond_wait}, \code{cond_timedwait}, \code{cond_signal}, \code{cond_broadcast} from \code{threads::shared}; example 03.}
\fig{condvar}{A condition-variable rendezvous. The waiter re-checks the predicate after every wakeup.}
The predicate is protected by the mutex; the wait atomically releases the lock while sleeping and reacquires it before returning. Publish the state before notifying. Broadcast when many waiters become eligible; signal when exactly one can make progress. The startup example uses a deadline and verifies worker results [2].
\CodeFile{perl}{examples/03-condition.pl}{\nolinkurl{03-condition.pl}}
\note{A timeout is not proof that work was cancelled or that a resource was never acquired elsewhere. Treat timeout outcomes as explicit state transitions. For elapsed-time metrics use \code{CLOCK_MONOTONIC}; \code{cond_timedwait} alone takes a wall-clock epoch deadline and is therefore sensitive to clock steps.}
\section{Reusable barriers}
\strategy{barrier}{Barrier}{A tour group that regroups at each landmark: nobody walks on to the next stop until everyone has arrived.}
{parallel steps must finish together before the next step starts, as in iterative computations or phased batch jobs.}
{participants can fail or hang independently; one missing member stalls everyone unless the barrier supports timeouts or a broken state.}
{a generation counter with \code{lock} and \code{cond_wait}; example 09.}
\fig{barrier}{A reusable barrier: each phase starts only when every participant has arrived.}
A generation counter prevents a fast thread entering the next phase from confusing a waiter from the previous phase. This example assumes a fixed cohort and no participant failures. A production barrier needs cancellation or a broken-barrier state, deadlines and supervision.
\CodeFile{perl}{examples/09-barrier.pl}{\nolinkurl{09-barrier.pl}}

% ================================================================ chapter 5

\chapter{Queues, worker pools and event loops}
\chapterbrief{Transfer ownership and bound asynchronous work.}
\par\smallskip\noindent\textbf{Abbreviations used in this chapter:} API --- application programming interface; DBI --- Database Interface; MCE --- Many-Core Engine.\par\smallskip

\index{Queues}
\index{Worker pools}
\index{Event loop}
\index{Backpressure}
\index{API!Thread::Queue}
\index{API!MCE}
\index{API!Future}

\section{Ownership instead of shared state}
\index{API!Thread::Queue}\index{API!enqueue}\index{API!dequeue}
\strategy{queue}{Work queue (ownership transfer)}{A conveyor belt between a kitchen and the waiters. The cook puts plates on the belt; each waiter takes one plate and is then its only owner. Nobody fights over a plate, so nobody needs a lock.}
{work can be split into independent jobs; this is the default design for worker pools.}
{jobs must survive a crash (an in-memory queue loses them: use a database queue or Redis Streams), or the consumer needs results from many jobs at once.}
{\code{Thread::Queue} with \code{limit} and \code{end} (example 04); database queues (example 12); Redis Streams (example 17).}
\fig{queue}{A bounded work queue: each job has exactly one owner at a time, and a full queue slows the producer down.}
\code{Thread::Queue} supports blocking, nonblocking and timed dequeue, a capacity limit and end-of-stream signaling. Use an explicit message envelope with an identifier and payload. With a bounded input queue the producer blocks until consumers make space. In the example the results queue is unbounded because the main thread consumes only after joining; bounding both queues without a concurrent result consumer could deadlock [4].
\CodeFile{perl}{examples/04-queue.pl}{\nolinkurl{04-queue.pl}}
\section{Operational patterns}
A single writer thread can own a file or database session and apply updates from many compute workers. Batch writes to amortize overhead, but cap batch size and maximum delay. A dedicated result consumer permits bounded outputs. Track queue depth, oldest-item age, throughput and service time; depth alone is not a latency measure. If workers fail, the supervisor must detect outstanding work and avoid waiting indefinitely for queue capacity. In-memory queues lose pending messages on process failure; durable workloads need a database queue (Section~\ref{sec:dbqueue}) or a broker such as Redis Streams (Section~\ref{sec:streams}), with explicit acknowledgment and idempotent effects.
\section{Process pools}
\index{API!MCE}\index{API!gather}
\code{MCE::Loop} supplies persistent workers and chunk distribution [9]. Create database and Redis connections inside each worker, never before the fork.
Example 27 shows the three habits that matter in practice: per-worker setup in \code{user_begin} (where connections belong), chunking so that each worker does a meaningful batch per round trip, and one \code{gather} per chunk rather than per item.
\CodeFile{perl}{examples/27-mce-loop.pl}{\nolinkurl{27-mce-loop.pl}}
\tested{The gathered result matched the closed-form sum exactly, with every chunk received once.}
\warn{For $N=2{,}000{,}000$, the sum exceeds $2^{53}$; a floating-point computation can round both the result and the reference value alike. Use exact arithmetic or reject inputs beyond the exact integer range, as example 27 does.}
\section{Event loops and async/await}\label{sec:async}
\index{API!Future::AsyncAwait}\index{API!Future::Mutex}\index{API!await}
\strategy{asyncmutex}{Async mutex in an event loop}{One cashier serving many customers: while one customer looks for their wallet (\code{await}), the cashier helps someone else. If both customers are editing the same order, the cashier must not switch between them halfway through.}
{a read-modify-write of shared state spans an \code{await}.}
{the critical section has no \code{await} in it (it is already atomic in a single-threaded loop), or the state is shared with other processes.}
{\code{Future::Mutex} with \code{Future::AsyncAwait}; bounded concurrency with \code{Future::Utils::fmap_void}; example 19.}
\fig{asyncrace}{A race in a single-threaded event loop, and the same code serialized with Future::Mutex.}
\code{Future::AsyncAwait} turns callback chains into sequential-looking code [29]. The single interpreter removes data races but not logical races: every \code{await} is a suspension point at which other tasks run. A read-modify-write that spans an \code{await} needs serialization, provided by \code{Future::Mutex} [30]. Bounded fan-out (an asynchronous semaphore) is provided by \code{fmap_void} and its siblings in \code{Future::Utils} with the \code{concurrent} option.
\CodeFile{perl}{examples/19-async-await.pl}{\nolinkurl{19-async-await.pl}}
\tested{With the mutex removed, the same program lost 8 of 40 increments on the verification host. Concurrency peaked at exactly the configured bound of 5.}
Use nonblocking clients or offload DBI work to \code{IO::Async::Function}. Race a task-owned future against a timeout future using \code{Future->wait_any}; inspect which outcome won and propagate cancellation. Do not cancel shared work merely because one observer stops waiting. Re-read state after an await or protect it with \code{Future::Mutex}.

% ================================================================ chapter 6



\section{Cancellation and task lifetime}
\index{Cancellation}\index{API!Future!cancel}
A parent operation should own its tasks and join or cancel them before returning.
Cancellation requests do not prove that remote effects did not occur.
Future cancellation callbacks [76] depend on underlying I/O cleanup support.
Bound outstanding work with an overall deadline rather than a new timeout
for every retry. Example 38 demonstrates cooperative cancellation.

\subsection{Propagating an end-to-end deadline}
\index{Deadline propagation}\index{API!clock_gettime@\texttt{clock\_gettime}}\index{API!Future!wait_any@\texttt{wait\_any}}
Use a monotonic clock [83] for an absolute deadline within one process.
If an operation has a two-second budget, a 400~ms admission wait leaves
1.6~s, a 600~ms database operation leaves 1.0~s, and a 250~ms backoff
leaves 750~ms for any retry and response. The same deadline is passed
to each stage, not reset at each call. Reserve time for cleanup; if no
budget remains, refuse to start new work. The following excerpt is an
adapter contract, not a complete network client.
\CodeFile{perl}{snippets/deadline-budget.pl}{Deadline propagation through admission, storage and response (illustrative excerpt)}
\begin{figure}[H]\centering
\begin{tikzpicture}[x=5.7cm,y=1cm]
\draw[navy,thick] (0,0)--(2,0);
\foreach \x/\lab in {0/0,.4/400,1/1000,1.25/1250,2/2000}{\draw[navy] (\x,-.07)--(\x,.07);\node[lbl,below] at (\x,-.08){\lab};}
\foreach \a/\b/\txt/\col in {0/.4/Admission/teal,.4/1/Database/navy,1/1.25/Backoff/rust,1.25/2/Remaining/leaf}{\fill[\col!14] (\a,.12) rectangle (\b,.55);\node[lbl,text=navy] at ({(\a+\b)/2},.33){\txt};}
\node[lbl,anchor=east] at (2,-.55){Elapsed milliseconds};
\end{tikzpicture}
\caption{One operation budget: each stage consumes the same deadline, including backoff.}
\end{figure}
Map remaining time to documented adapter-specific connection, acquisition,
statement and response timeouts; an adapter may have only second-level
resolution or no cancellation API. A timeout on the waiting caller is not
a guarantee that database or remote work stopped. For a write with an
unknown result, reconcile by durable operation identity instead of retrying
blindly. Clean up task-owned resources and wait for children to terminate
where required. For cross-host calls transmit a remaining relative budget
or a protocol-defined deadline, and account for transport time and clock
skew; monotonic clock readings cannot be compared between hosts.

\chapter{Reader/writer locking and lock design}
\chapterbrief{Choose granularity, acquisition order and cleanup.}
\par\smallskip\noindent\textbf{Abbreviations used in this chapter:} FIFO --- first in, first out; PID --- process identifier.\par\smallskip

\index{Reader/writer locks}
\index{Lock ordering}
\index{Lock striping}
\index{API!defer}

\section{Reader/writer compatibility}
\strategy{rwlock}{Reader/writer lock}{A museum exhibit: any number of visitors may look at once, but when the curator rearranges it, the room is closed to everyone.}
{reads are frequent, take noticeable time, and can safely overlap; writes are rare.}
{critical sections are tiny (a plain mutex is cheaper), or you would need to upgrade a read lock to a write lock.}
{built from \code{lock} and \code{cond_wait} (example 05); for processes, \code{flock} \code{LOCK_SH} and \code{LOCK_EX}.}
\begin{center}\begin{tabular}{lcc}\toprule Held / requested & Read & Write\\\midrule Read & Compatible & Wait\\Write & Wait & Wait\\\bottomrule\end{tabular}\end{center}
A reader/writer lock can improve throughput when read sections are substantial and frequently overlap. It adds bookkeeping and wakeup costs, so a simple mutex often wins for small operations. The following writer-preferring algorithm blocks new readers once a writer is queued; this mitigates writer starvation but may starve readers. It is not FIFO, reentrant, upgradeable or cancellation-safe. Never upgrade while holding a read lock: release, reacquire for writing, then revalidate.
\fig{rwtimeline}{Writer preference in example 05: once a writer queues, later readers wait behind it, so the writer cannot starve.}
\CodeFile{perl}{examples/05-rwlock.pl}{\nolinkurl{05-rwlock.pl}}
\section{Granularity and ordering}
\strategy{ordering}{Lock ordering and lock striping}{Two people each holding one end of a rope, each waiting for the other to let go, wait forever. Agree that everyone always grabs the left end first, and the standoff cannot happen. Striping means using a fixed set of locks shared by many items, like a building with one key per floor.}
{an operation needs more than one lock (transfers, multi-key updates), or one global lock has become a bottleneck.}
{a single lock is fast enough: every extra lock adds ordering rules that every future maintainer must follow.}
{acquire \code{lock}s in a documented global order; hold a variable number by recursion; examples 23 and 11 (MySQL rows in primary-key order).}
Coarse locks are easy to audit but serialize unrelated operations. Fine-grained locks can improve parallelism while introducing ordering constraints and memory overhead. Lock striping hashes keys onto a fixed set of mutexes; unrelated keys can collide, but lock storage stays bounded. Multi-key operations must acquire unique stripes in sorted order and release in reverse order. Avoid calling user callbacks under a lock unless that is an intentional, documented contract.
Example 23 applies this to 64 accounts guarded by 8 stripe locks. A transfer locks its two stripes in ascending order; an auditor locks \emph{all} stripes, also in ascending order, to read a consistent total while transfers continue. Because \code{lock} releases at scope exit, a variable number of locks is held by recursion, one stripe per stack frame.
\fig{striping}{Lock striping: a transfer locks only its two stripes, in ascending order, so unrelated transfers run in parallel.}
\CodeFile{perl}{examples/23-lock-striping.pl}{\nolinkurl{23-lock-striping.pl}}
\tested{Four threads completed 19{,}589 transfers while the auditor took 17{,}771 consistent totals; money was neither created nor destroyed and no balance went negative.}
\fig{deadlock}{A wait-for cycle and how a global acquisition order removes it.}
A deadlock requires mutual exclusion, hold-and-wait, no preemption and a circular wait. Enforce a total ordering or avoid nested acquisitions. Livelock occurs when competing workers keep retrying without progress; randomized backoff helps. Starvation is indefinite denial despite overall progress. Priority inversion occurs when a lower-priority owner blocks a higher-priority waiter; Perl-level locks provide no priority-inheritance contract.
\section{Releasing on every path}\label{sec:release}
\index{API!defer}
\strategy{release}{Release on every path (scope-based cleanup)}{Whatever happens in the room (you finish, you are called away, the fire alarm goes off), you must hand back the key on the way out. Tie the return of the key to leaving the room, not to remembering to do it.}
{anything is acquired that is not released automatically: file locks, semaphore permits, database locks, Redis leases.}
{never; the only question is which mechanism to use.}
{\code{defer} and \code{try}/\code{catch} (example 24), guard objects with \code{DESTROY} (example 20), or \code{eval} plus rethrow (examples 02 and 13).}
\fig{releasepaths}{Every exit from a critical section must pass through the release.}
\code{lock} releases at scope exit automatically. Every other resource (a file lock, a semaphore permit, a database lock, a Redis lease) needs explicit release on success, on exception and on early return. Three idioms are available.
\begin{itemize}
\item \textbf{\code{eval} with captured error.} Portable to every supported Perl. Capture \code{$@} immediately, release, then rethrow, as in examples 02 and 13.
\item \textbf{\code{try}/\code{catch}.} Stable since Perl 5.40 and enabled by \code{use v5.40}; \code{return} inside \code{try} returns from the enclosing subroutine, unlike \code{eval}. The \code{finally} block and \code{defer} remain experimental through 5.44 [28,46].
\item \textbf{Scope guards.} An object whose destructor releases the resource. Deterministic for lexicals in Perl, but guard against inherited copies after \code{fork}, reference cycles and global destruction, where destruction order is undefined.
\end{itemize}
Example 24 packages the modern idiom as a reusable helper: \code{with_flock} acquires the lock, registers its release with \code{defer}, runs the caller's code inside \code{try}, and on failure rethrows with context. The example verifies release after a normal return, after an exception and when calls are nested.
\version{Example 24 enables experimental \code{defer} on Perl 5.36 or later. Prefer lexical guards where experimental syntax is not permitted.}
\CodeFile{perl}{examples/24-try-defer.pl}{\nolinkurl{24-try-defer.pl}}
Where \code{defer} is not acceptable, because experimental features are disallowed or older Perls must be supported, a guard object gives the same guarantee on any Perl:
\CodeFile{perl}{examples/20-scope-release.pl}{\nolinkurl{20-scope-release.pl}}
\warn{A destructor cannot safely perform network I/O during global destruction, and a guard copied into a forked child would release the parent's lock when the child exits. The guard above records its creating PID and releases only there. Release remote leases explicitly; rely on expiry only as the last resort.}

% ================================================================ chapter 7

\chapter{Cross-process synchronization on Linux}
\chapterbrief{Coordinate processes and handle descriptor lifetimes.}
\par\smallskip\noindent\textbf{Abbreviations used in this chapter:} API --- application programming interface; DBI --- Database Interface; IPC --- interprocess communication; MCE --- Many-Core Engine; PID --- process identifier; POSIX --- Portable Operating System Interface.\par\smallskip

\index{File locking}
\index{Fork}
\index{System V IPC}
\index{API!flock}
\index{API!IPC::Semaphore}

\section{File locks}
\index{API!flock}\index{API!LOCK_EX@\texttt{LOCK\_EX}}\index{API!LOCK_NB@\texttt{LOCK\_NB}}
\strategy{flock}{File lock (flock)}{A ``room occupied'' sign on a door that every process on the machine agrees to respect. It is advisory: a program that ignores signs can still walk in.}
{processes on one host must not run a section at the same time: cron jobs, deploy scripts, single-writer files.}
{the processes run on different hosts or containers that do not share a local filesystem, or the file is on NFS/SMB without testing.}
{\code{flock} with \code{LOCK_EX}/\code{LOCK_SH}/\code{LOCK_NB} (example 06); a deadline (example 25); util-linux \code{flock -n} for cron.}
\code{flock} is advisory on ordinary local Linux filesystems. An exclusive lock blocks conflicting shared and exclusive requests; multiple shared locks can coexist. Linux associates \code{flock} locks with open file descriptions: a descriptor inherited through \code{fork} or duplicated with \code{dup} shares the lock. Open the file independently after \code{fork} when workers must compete. Close all duplicates, or call \code{LOCK_UN}, to release. Do not unlink an active lockfile: a replacement inode creates an independent coordination domain [5].
\fig{ofd}{flock ownership follows the open file description, not the process.}
\CodeFile{perl}{examples/06-flock.pl}{\nolinkurl{06-flock.pl}}
The lab demonstrates exclusion, not crash-safe persistence: a crash between truncate and write can corrupt the file. For durable replacement, hold a separate stable lockfile, write a temporary file on the same filesystem, flush and \code{fsync} it, rename atomically, then \code{fsync} the directory if the durability contract requires it.
\code{flock} has no timeout. Example 25 adds one: it polls with \code{LOCK_NB} against a monotonic deadline, with capped exponential backoff and jitter, and distinguishes contention (\code{EWOULDBLOCK}) from real errors. A child process holds the lock so that the timeout path is actually exercised. Then the example confirms that shared locks coexist while an exclusive request is refused.
\fig{flockpoll}{Bounded acquisition: non-blocking attempts with growing, jittered pauses until a deadline.}
\CodeFile{perl}{examples/25-flock-deadline.pl}{\nolinkurl{25-flock-deadline.pl}}
Use a private administrator-controlled directory, restrict permissions and avoid symlink races in shared writable directories. For cron, the util-linux wrapper \code{flock -n /run/lock/myapp.lock command} is often sufficient. Traditional POSIX record locks (\code{fcntl}) are byte-range locks owned by the process, and closing \emph{any} descriptor for the file drops them; Linux open-file-description (OFD) record locks avoid that hazard [17]. On NFS, Linux emulates \code{flock} with byte-range locks, so \code{flock} and \code{fcntl} locks on the same file then interact; never assume local semantics on NFS or SMB without testing the actual mount and server.
\tip{Lock a dedicated \code{.lock} file, not the data file itself. Code that rewrites data with write-to-temporary-then-\code{rename} replaces the inode, and a lock taken on the old file then protects nothing.}

\section{System V and POSIX semaphores}
\index{API!IPC::Semaphore}
\strategy{sysvsem}{System V semaphore}{The car-park counter again, but kept by the kernel, so unrelated processes can share it. With \code{SEM_UNDO}, the kernel returns a ticket if a car is towed away (the process dies).}
{forked processes on one host need a shared permit count.}
{you can use \code{flock} (simpler) or a database (cross-host), or you cannot guarantee someone removes the semaphore set afterwards.}
{\code{IPC::Semaphore} with \code{IPC::SysV} constants; example 07.}
\fig{sysvsem}{A System V semaphore set lives in the kernel, outside every process that uses it.}
System V semaphore sets are kernel objects identified by IDs or keys; \code{IPC::Semaphore} wraps operations on them [7]. \code{SEM_UNDO} requests kernel adjustment at process exit: it restores a permit after abnormal exit but cannot repair a partially updated file or database. Creation and removal need a clear owner. The lab creates a private set before \code{fork}, removes it from an \code{END} block that runs only in the creating process, and has children leave through \code{POSIX::_exit} so they never run it.
\CodeFile{perl}{examples/07-sysv.pl}{\nolinkurl{07-sysv.pl}}
\code{END} blocks do not run on \code{SIGKILL} or when the parent is killed by an unhandled signal, so supervisors must still account for orphaned IPC objects and kernel limits. Use \code{ipcs -s} to inspect and \code{ipcrm} to remove only objects owned by your application. POSIX named semaphores use \code{sem_open}/\code{sem_unlink}; unnamed process-shared semaphores require shared memory; neither has \code{SEM_UNDO} semantics [8]. Perl has no core POSIX semaphore API; use a validated CPAN binding if one is needed.
\section{Pipes, sockets and shared memory}
A pipe provides byte-stream communication between processes, not a shared object. Close unused ends after \code{fork} so EOF is observable. With multiple writers, only writes of at most \code{PIPE_BUF} bytes (4096 on Linux) are atomic; buffered \code{print} calls are not necessarily single writes. Use framing, length limits, partial-I/O handling and backpressure for production protocols.
\CodeFile{perl}{examples/08-pipe.pl}{\nolinkurl{08-pipe.pl}}
Unix domain sockets support bidirectional IPC and can separate privilege boundaries. Shared memory (System~V, POSIX or \code{mmap}) reduces copying but does not serialize access. \code{IPC::Shareable} and \code{MCE::Shared} simplify shared structures [16]; check the locking semantics of multi-step operations. A read followed by a write is not one atomic transaction merely because each method is synchronized.

\section{Locks across fork, exec and signals}
Lock ownership differs by mechanism, and \code{fork} exposes the differences. The behaviors below were verified on Linux with Perl 5.38:
\begin{longtable}{L{.24\linewidth}L{.68\linewidth}}\toprule Mechanism & Across fork and close \\ \midrule\endhead
\code{flock} & Owned by the open file description. A child that inherits the descriptor shares the parent's lock (its own \code{LOCK_NB} request succeeds immediately). A child closing its copy does not release the parent's lock. \\
POSIX \code{fcntl} record lock & Owned by the process. A child does \emph{not} inherit it and competes like any other process. Closing \emph{any} descriptor for the file, even one opened elsewhere in the program, silently releases all of that process's locks on it. \\
Linux OFD record lock & Owned by the open file description, like \code{flock}; avoids the close hazard [17]. \\
System~V semaphore & Values are kernel-global; \code{SEM_UNDO} adjustments are per process and applied at that process's exit. \\
\code{threads::shared} & Meaningless across processes: after \code{fork} each process has its own copy. \\
DBI and Redis connections & Inherited sockets must never be used by both processes; create new connections in the child. \\
\bottomrule\end{longtable}
Perl's \code{alarm} does interrupt a blocking \code{flock}: with \texttt{local \$SIG\{ALRM\} = sub \{ die "timeout\textbackslash n" \}; alarm 1;} a contended lock gave up after 1.0\,s. Prefer the polling deadline of example 25 in libraries and servers anyway. There is only one alarm timer per process, so nested users overwrite each other. Other code in the process may rely on \code{alarm} too, and signals do not mix predictably with interpreter threads.

\tip{Open database and Redis connections lazily in each process and check the PID before reuse, as \code{SessionLab::Redis} does. With DBI, set \code{AutoInactiveDestroy} so a child's exit cannot close the parent's connection.}

\needspace{25\baselineskip}\section{Preforking web servers}
Starman, Gazelle, Hypnotoad and uWSGI serve requests from a pool of forked workers (Figure~\ref{fig:prefork}). Each worker is a separate process with its own interpreter, so a Perl variable, a \code{lock} or an in-memory cache in one worker is invisible to the others, and to other pods and cron jobs. Anything that must be coordinated between requests belongs in Redis, the database or, for a single host, a file lock. Chapter~\ref{ch:catalyst} applies this to Catalyst sessions.
\begin{figure}[H]\centering
\begin{adjustbox}{max width=.85\linewidth}\begin{tikzpicture}[x=1cm,y=1cm]
\node[box,minimum width=3.4cm] (lb) at (0,3.3) {ingress / load balancer};
\node[box,minimum width=3.4cm] (m) at (0,1.8) {Starman master\\\footnotesize listens, forks, supervises};
\draw[arr] (lb) -- (m);
\foreach \i/\x in {1/-5.4,2/-1.8,3/1.8,4/5.4} {
  \node[box,minimum width=3.1cm,minimum height=1.3cm] (w\i) at (\x,0) {worker \i\\\footnotesize own interpreter,\\\footnotesize own lexicals and lock()};
  \draw[thinarr] (m) -- (w\i);
}
\node[badbox,text width=12.6cm] at (0,-1.45) {A lock() or Perl variable in worker 1 is invisible to workers 2--4, to other pods and to cron jobs.};
\node[okbox,minimum width=3.4cm] (r) at (-3.6,-3.0) {Redis / Valkey\\\footnotesize sessions, leases, queues};
\node[okbox,minimum width=3.4cm] (d) at (0,-3.0) {MySQL\\\footnotesize transactions, constraints};
\node[okbox,minimum width=3.4cm] (f) at (3.6,-3.0) {flock on a local file\\\footnotesize same host only};
\node[lbl] at (0,-4.0) {coordination between workers must live outside the interpreter};
\end{tikzpicture}
\end{adjustbox}
\caption{A preforking server. In-process state and locks are per worker; coordination lives outside the interpreter.}\label{fig:prefork}\end{figure}

% ================================================================ chapter 8

\partstart{III / Database concurrency}{Database concurrency and application transactions}{Persistent state demands storage-engine guarantees, business transaction boundaries and durable identities.}
\chapter{MySQL locking and transaction design}
\chapterbrief{Understand storage-engine locks, isolation, guarded updates and DBI/DBIx::Class ownership.}
\par\smallskip\noindent\textbf{Abbreviations used in this chapter:} API --- application programming interface; DBI --- Database Interface; DDL --- data definition language; MVCC --- multi-version concurrency control; SQL --- Structured Query Language; ID --- identifier; DBD --- Database Driver.\par\smallskip

\index{MySQL}
\index{InnoDB}
\index{Isolation levels}
\index{Deadlock}
\index{API!DBI}
\index{API!DBIx::Class}
\index{API!FOR UPDATE}
\index{API!SKIP LOCKED}
\index{API!GET_LOCK@\texttt{GET\_LOCK}}

\section{Lock taxonomy in InnoDB}
MySQL 8.4 LTS and 9.7 LTS are the documentation baselines; 9.7 is also the last release line with sequential version numbers, after which MySQL adopts calendar versions such as 26.7 [32]. InnoDB uses multi-version concurrency control (MVCC) for ordinary consistent reads and locks for conflicting writes and locking reads. A consistent read can see an older committed snapshot without blocking a writer; it is not a reservation. InnoDB locks index records, not application objects: the query plan and indexes determine which records or ranges are locked [11].
\fig{nextkey}{Locks observed in \texttt{performance\_schema.data\_locks} for one range query on a secondary index.}
The figure shows what MySQL 8.0.46 reported for a four-row table with a secondary index on \code{v}. Under REPEATABLE READ the range 15--25 took next-key locks covering $(10,20]$ and $(20,30]$, plus record locks on the matching primary-key rows. That blocks inserts anywhere from 11 to 30 and locks the row with $v=30$, which lies outside the range. An equality search for the absent value 25 took a gap lock before 30. Under READ COMMITTED the same query locked only the record $v=20$.

The plan matters as much as the isolation level. A \code{FOR UPDATE} range on the \emph{primary key} of the same table, \code{WHERE id BETWEEN 2 AND 3}, was executed as a full scan of the covering secondary index. It locked every index record and the supremum pseudo-record, blocking all inserts into the table. Check \code{EXPLAIN} and \code{data_locks} for every locking query before it reaches production.
\begin{longtable}{L{.24\linewidth}L{.68\linewidth}}\toprule Lock & Purpose \\ \midrule\endhead
Record S/X & Shared or exclusive access to index records \\
Gap & Prevents inserts into an index interval; gap locks from different transactions can coexist \\
Next-key & Record lock plus the preceding gap; used for range protection \\
Intention IS/IX & Table-level declaration of intended record locking \\
Insert intention & Coordinates concurrent inserts into one gap \\
AUTO-INC & Allocation coordination; behavior depends on \code{innodb_autoinc_lock_mode} \\
Metadata lock & Protects object definitions against conflicting DDL \\
Named advisory lock & Session-owned application-level name, independent of transaction rollback \\
\bottomrule\end{longtable}
\section{Isolation levels}
READ UNCOMMITTED may observe uncommitted data. READ COMMITTED takes a fresh snapshot for each consistent read and generally reduces gap locking, with exceptions such as foreign-key and duplicate-key checks. REPEATABLE READ (the InnoDB default) keeps one snapshot for consistent reads in a transaction and uses next-key locking for many range-locking operations. SERIALIZABLE adds stronger coordination at a corresponding concurrency cost. Locking reads and snapshot reads may see different versions within one transaction. Never use a snapshot check as a substitute for enforcing the constraint in the write itself.
\CodeFile{sql}{snippets/guide-03.sql}{Atomic guarded update, without an explicit SELECT lock}
Prefer a single conditional statement or a unique constraint to a check-then-act sequence; enforce idempotency with a unique request ID, not a \code{SELECT} followed by an unprotected \code{INSERT}. Optimistic concurrency uses a version column: \code{UPDATE ... WHERE id=? AND version=?} and increment the version; zero affected rows indicates a conflict or a missing object, and the caller must reread or fail.
\strategy{optimistic}{Optimistic concurrency (version column)}{Like editing a shared document that shows ``version 7''. You save only if it is still version 7; if someone saved first, you reload and reapply your change instead of overwriting theirs.}
{conflicts are rare and users ``think'' between reading and writing (forms, editors, APIs).}
{the same row is hot and contended; retries will dominate (use a pessimistic lock or redesign).}
{\code{UPDATE ... WHERE id = ? AND version = ?} and check the affected-row count; with DBIx::Class, \code{->search({ version => $seen })->update(...)}; example 29.}
\fig{optimistic}{Optimistic concurrency: the losing writer detects the conflict from a zero-row update.}
\section{Pessimistic locking reads}
\index{API!FOR UPDATE}
\strategy{rowlocks}{Pessimistic row locks (FOR UPDATE, NOWAIT, SKIP LOCKED)}{Reserving the exact items you are about to change, so nobody else changes them before you have finished. NOWAIT means ``if it is reserved, tell me immediately''; SKIP LOCKED means ``give me the next unreserved one''.}
{a transaction reads rows and then updates them based on what it read, or workers claim rows from a queue table.}
{the transaction would wait on users or remote APIs while holding the lock, or a single guarded \code{UPDATE} can express the rule.}
{\code{SELECT ... FOR UPDATE} through DBI (examples 11, 12) or DBIx::Class \code{for => 'update'} (example 29).}
\fig{rowlocks}{Three ways to request a row that another transaction has locked.}
\CodeFile{sql}{snippets/guide-04.sql}{Session A: lock a row until transaction completion}
\CodeFile{sql}{snippets/guide-05.sql}{Session B: run while session A holds its lock}
\code{FOR SHARE} permits compatible shared locking reads; \code{FOR UPDATE} acquires exclusive locks on the examined records. \code{NOWAIT} fails immediately on a conflicting row lock; \code{SKIP LOCKED} omits locked rows and suits queue tables, not consistent reports. Neither modifier avoids other sources of blocking such as metadata locks. Locking reads need an explicit transaction to hold their locks beyond the statement [12].
\section{DBI ownership and driver selection}\label{sec:dbi}
Create a DBI connection inside the worker process or thread that uses it, and never use an inherited connection after \code{fork}. If pre-fork connections cannot be avoided, set \code{AutoInactiveDestroy} so a child's exit does not close the parent's session; it does not make sharing a live connection safe. Disable automatic reconnection (\code{mysql_auto_reconnect => 0}) for lock-dependent sessions: a reconnect silently discards session locks and transaction state [15,18].

Set \code{innodb_lock_wait_timeout} per session to a value consistent with the caller's retry budget. The 50-second server default otherwise governs how long a blocked worker waits before error 1205, which dwarfs any client-side backoff.
\version{\code{DBD::mysql} 5.x builds only against MySQL 8.x (or newer) client libraries, although the resulting driver can still connect to older servers [18]. Hosts that have MariaDB client libraries should use \code{DBD::MariaDB}, which renames driver attributes with a \code{mariadb_} prefix and always treats strings as Unicode [31]. Pin the 4.x series only as a deliberate, time-limited migration measure.}
Inject credentials through a secret store or the runtime environment, and never log them.
\CodeFile{bash}{snippets/guide-06.sh}{Laboratory setup (disposable database only)}
\CodeFile{sql}{examples/schema.sql}{Disposable laboratory schema}
\section{Transactions with deterministic lock ordering}
The transfer example locks both accounts in primary-key order, enforces the balance constraint in the \code{UPDATE}, and commits both changes together. It retries only deadlock and lock-wait victims, after rollback, with exponential backoff and jitter. Integer cents avoid floating-point rounding. Real financial transfers additionally need an idempotency ledger and a strategy for uncertain commit outcomes.
\CodeFile{perl}{examples/11-mysql-transfer.pl}{\nolinkurl{11-mysql-transfer.pl}}
Error 1213 marks a deadlock victim; InnoDB has rolled back the whole transaction. Error 1205 by default rolls back only the failing statement (\code{innodb_rollback_on_timeout} changes this); roll back explicitly before retrying either. Capture error codes before another DBI call can replace them. A lost connection during \code{COMMIT} is ambiguous: never blindly repeat a non-idempotent transfer. A failed rollback means the connection must be discarded. Use bounded retries and an overall operation budget, and never retry validation or integrity errors [13].
\tested{Four processes each performed 100 transfers in alternating directions against the same two rows. The total balance remained exactly 20{,}000 cents and no attempt exhausted its retry budget.}
\tip{Set \code{innodb_print_all_deadlocks = ON}. \code{SHOW ENGINE INNODB STATUS} keeps only the latest deadlock; with this setting every one is written to the error log, where your log pipeline can count and alert on them.}

\section{Named advisory locks}
\index{API!GET_LOCK@\texttt{GET\_LOCK}}\index{API!RELEASE_LOCK@\texttt{RELEASE\_LOCK}}
\strategy{advisory}{Named advisory lock (GET\_LOCK)}{A named token kept at the database's front desk. Your database connection signs it out; if the connection drops, the desk silently takes the token back, even if you are still working.}
{several hosts must not run a job at once and they already share a MySQL server, but the job does not fit in one transaction.}
{the work's correctness depends on the lock (use constraints and fencing), or connections are pooled and may be swapped.}
{\code{SELECT GET_LOCK(?, ?)} and \code{RELEASE_LOCK} on a pinned, non-reconnecting DBI handle; example 10.}
\fig{advisory}{MySQL named locks belong to a session, not to a transaction.}
\code{GET_LOCK(name, timeout)} returns 1 on acquisition, 0 on timeout and \code{NULL} on error; names are limited to 64 characters. Locks are session-scoped, may be acquired recursively and need matching releases; \code{COMMIT} and \code{ROLLBACK} do not release them. They coordinate only sessions on the same server, not independent primaries or across a failover. The connection must stay pinned; a pool must never reuse a connection that still holds a lock [14].
\CodeFile{perl}{examples/10-mysql-advisory.pl}{\nolinkurl{10-mysql-advisory.pl}}
\note{A session lock disappears on disconnect while the client may still believe it holds it. For external side effects, use durable idempotency and, where appropriate, fencing tokens validated by the target. A named lock is not a substitute for a transaction.}
\section{Database work queues}\label{sec:dbqueue}
\index{API!SKIP LOCKED}\index{Job claims}
\strategy{dbqueue}{Database work queue with leases}{A ticket rail in a kitchen. A cook takes a ticket and writes their name and a ``back by'' time on it. If they do not finish in time, the ticket returns to the rail for someone else, and the late cook's work is no longer accepted.}
{jobs must survive crashes and deploys, and workers run on several hosts that share the database.}
{volume or latency needs exceed what the database should carry; consider Redis Streams or a broker.}
{\code{FOR UPDATE SKIP LOCKED} claims, attempt numbers as fencing tokens, conditional completion; examples 12 and 29.}
The queue uses a \emph{visibility timeout}. A pending job is claimable once \code{visible_at} has passed. Claiming moves \code{visible_at} forward by the lease length, records an owner and increments \code{attempt}, all in one short transaction using \code{FOR UPDATE SKIP LOCKED}. A worker that dies simply stops renewing, and the job becomes claimable again without a separate sweeper. Work happens after the claim commits, so slow work never holds a row lock [12].

The attempt number is the fencing token. Heartbeats, failure reports and completion all carry \code{WHERE id=? AND owner=? AND attempt=?}, so a worker whose lease expired and was reclaimed can no longer change the job. Completion and the job's effect commit in one transaction, and the effect table's primary key makes a duplicate effect impossible even under a race. Jobs that exhaust \code{max_attempts} move to a dead-letter state for inspection.
\fig{queuestates}{Job states in the visibility-timeout queue of example 12.}
\CodeFile{perl}{examples/12-mysql-queue.pl}{\nolinkurl{12-mysql-queue.pl}}
\tested{Four concurrent workers completed 220 jobs; every job produced exactly one effect at attempt 1. A job whose lease had expired was reclaimed at attempt 2, an exhausted job was dead-lettered, and a completion attempted by the expired owner matched zero rows.}
\code{SKIP LOCKED} does not guarantee fairness or strict ordering under concurrency. Effects outside the database (an email, a payment API) cannot join the transaction; make them idempotent with a request key derived from the job ID and attempt, or stage them through an outbox table.
\section{Transactions and locking with DBIx::Class}\label{sec:dbic}
\index{API!txn_do@\texttt{txn\_do}}\index{API!txn_scope_guard@\texttt{txn\_scope\_guard}}
Most modern Perl applications reach MySQL through DBIx::Class rather than raw DBI. The same rules apply, but four DBIx::Class behaviors matter for locking. Each was confirmed in the source of version 0.082843 and exercised by example 29.
\begin{itemize}
\item \textbf{\code{txn_do} does not retry deadlocks.} It reruns the code block once, and only if the connection is gone after the first failure. A deadlock (1213) or lock-wait timeout (1205) propagates to the caller, so retry policy is the application's job (\code{with_deadlock_retry} in the example). Because of the reconnect retry, the block must also be safe to run twice if a connection can drop during \code{COMMIT}.
\item \textbf{\code{txn_scope_guard}} is the scope-based alternative: \code{my $guard = $schema->txn_scope_guard; ...; $guard->commit;}. If the scope is left without \code{commit}, by return or by exception, the transaction is rolled back. This is the release-on-every-path principle (Section~\ref{str:release}) applied to transactions.
\item \textbf{Locking reads} use the \code{for} attribute: \code{for => 'update'} emits \code{FOR UPDATE}. On MySQL, \code{for => 'shared'} emits the legacy \code{LOCK IN SHARE MODE}, and the MySQL SQL generator rejects the scalar-reference form that the ResultSet documentation describes. Example 29 subclasses it in nine lines so that \code{for => \'UPDATE SKIP LOCKED'} works.
\item \textbf{Fork awareness.} DBIx::Class checks the process ID and reconnects in a forked child, but connecting explicitly inside each worker keeps ownership obvious.
\end{itemize}
Guarded updates read naturally as resultset updates. \code{$rs->search({ id => $from, cents => { '>=' => $amt } })->update(...)} returns the affected-row count, so the rule ``never overdraw'' lives in the \code{WHERE} clause. Optimistic concurrency is the same idiom with a version column.
\CodeFile{perl}{examples/29-dbic-locking.pl}{\nolinkurl{29-dbic-locking.pl}}
\tested{Four processes made 200 opposing transfers inside \code{txn_do} with \code{FOR UPDATE} and preserved the total exactly. Four writers appended 80 lines to one document with optimistic version checks: no edit was lost, and about 200 conflicts were detected and retried, which shows the cost of optimism under heavy contention. Two workers drained 20 freshly seeded jobs with \code{FOR UPDATE SKIP LOCKED}, splitting them roughly evenly (10 and 10, or 9 and 11, across runs), and never claimed the same job twice.}

\tip{Write one small retry helper around \code{txn_do} and use it everywhere: retry the whole transaction on deadlock (1213), and on lock-wait timeout (1205) only if the work is idempotent, with capped exponential backoff and jitter. Never retry inside the transaction.}

\section{Metadata locks and legacy tables}
DDL may wait for transactions that retain metadata locks even while their queries are idle. Monitor metadata lock waits during migrations and ensure transactions end promptly. \code{LOCK TABLES} is a coarse session-level mechanism with implicit-commit interactions and is generally wrong for ordinary InnoDB application transactions. MyISAM relies on table locking and lacks transactions. Older MySQL versions may lack \code{NOWAIT}, \code{SKIP LOCKED}, enforced \code{CHECK} constraints, expression defaults or modern Performance Schema views; MariaDB behavior is not assumed identical.
\needspace{15\baselineskip}\section{Observability and contention diagnostics}
\CodeFile{sql}{snippets/guide-07.sql}{Inspect InnoDB and metadata lock contention}
Validate index selection with \code{EXPLAIN} before deploying locking queries. Examine blocking transactions, statement duration, lock hold time and transaction age. Performance Schema visibility depends on privileges and instrumentation. Killing a query does not necessarily end its transaction; verify the blocking session and the operational impact before terminating it. Never wait on a remote API while holding row locks when an outbox or staged workflow can provide equivalent correctness.



% ================================================================ chapter 11
\chapter{Transaction boundaries in layered applications}
\label{ch:enterprise-transactions}
\chapterbrief{Place transactions at the business boundary; classify retries.}
\index{Transactions}
\index{Savepoints}
\index{API!txn_do}
\index{API!txn_scope_guard}

A service operation should define the atomic unit of work. The controller
validates the request; repository methods participate in the transaction
rather than committing independently. A database commit belongs to the
business operation, not to an individual ORM result object.

\begin{figure}[!htbp]\centering
\begin{tikzpicture}[node distance=5mm]
\node[box,text width=10cm] (c) {Controller or job handler: validate and assign operation ID};
\node[okbox,below=of c,text width=10cm] (s) {Service: begin transaction; define retry boundary};
\node[box,below left=8mm and 0mm of s.south,text width=4.2cm,anchor=north east] (a) {Repository A\\guarded stock update};
\node[box,below right=8mm and 0mm of s.south,text width=4.2cm,anchor=north west] (b) {Repository B\\audit or outbox insert};
\node[okbox,below=8mm of $(a.south)!0.5!(b.south)$,text width=10cm] (d) {Commit together or roll back together};
\draw[arr](c)--(s); \draw[arr](s)--(a); \draw[arr](s)--(b);
\draw[arr](a)--(d); \draw[arr](b)--(d);
\node[box,below=of d,text width=10cm] (r) {Return response; relay external effects separately};
\draw[arr](d)--(r);
\end{tikzpicture}
\caption{A layered application has one explicit business transaction boundary.}
\end{figure}

\strategy{app-transactions}{Service-owned transactions}{One receipt covers all items or none.}{multiple changes must commit together.}{remote effects cannot share the same transaction.}{DBIx::Class, DBI, InnoDB}

\section{DBIx::Class composition}
Use \code{txn_do} for a complete unit of work or \code{txn_scope_guard}
when a lexical transaction guard better expresses the lifecycle.
Nested \code{txn_do} calls do not independently commit: the outer transaction
owns the commit [65]. Savepoints, when supported and correctly configured,
allow partial rollback, not independent durability. Do not catch a failed
nested operation and continue as if the transaction were still valid.
Avoid mixing raw DBI transaction calls with ORM-managed transaction depth.

\section{Retry the operation, not the final statement}
InnoDB deadlocks roll back a transaction; a lock-wait timeout normally
rolls back only the failed statement [67]. Explicitly abandon the failed
unit before retrying it. Capture driver error details promptly, use a
bounded retry budget with jitter, and reacquire all data needed by the
operation. A connection failure during commit can leave the outcome
unknown. Resolve it by a durable request ID rather than blindly retrying
a charge, decrement or external notification.

\begin{longtable}{L{.30\linewidth}L{.62\linewidth}}
\toprule Failure & Service behavior \\ \midrule\endhead
Validation error & Return a business error; do not retry. \\
Deadlock & Roll back/reset the unit; retry from the beginning with bounded backoff. \\
Lock wait timeout & Roll back the entire business operation explicitly before retry. \\
Unique-key conflict & Determine whether this is a duplicate operation or a different payload. \\
Lost commit response & Reconcile by durable operation ID; do not assume rollback. \\
\bottomrule
\end{longtable}

\section{Laboratory: a reservation and its audit record}
The example below deliberately fails between a stock decrement and an audit
insert, verifies rollback, then verifies successful commit and duplicate
rejection. SQLite is used only to make this application-layer test
self-contained; the MySQL chapter covers InnoDB-specific lock semantics.
\CodeFile{perl}{examples/33-layered-transaction.pl}{Example 33: service-owned DBIx::Class transaction}

\tip{Keep expensive computation and network I/O outside the transaction.
If the computation depends on mutable database state, validate its version
again inside the transaction before committing.}


\section{Savepoints and retry boundaries}
\index{API!auto_savepoint}\index{Savepoints}
Enable \code{auto_savepoint} only with a supported storage driver and
test the nested failure contract [75]. An inner successful operation is
not durable until the outermost transaction commits. A failed savepoint
can be handled locally only when the driver and ORM have actually rolled
back to it. Example 37 runs against SQLite by default, or a disposable
MySQL DSN when configured (Appendix~\ref{app:capstone}). It tests
partial rollback and outer rollback of a successful inner call.

\section{Reusable retry policy}
Example 38 separates confirmed conflicts from unknown outcomes and checks
a monotonic deadline and cooperative cancellation between attempts.
Driver-specific code must classify errors before entering this helper;
a deadlock may permit a retry, whereas an ambiguous commit requires
reconciliation. Do not retry arbitrary exceptions or remote effects.

\partstart{IV / Distributed coordination}{Distributed coordination}{Learn how coordination and cache services differ from transactional storage.}
% ================================================================ chapter 9
\chapter{Redis and Valkey: leases, semaphores and streams}\label{ch:redis}
\chapterbrief{Coordinate distributed workers while respecting failure, durability and ownership boundaries.}
\par\smallskip\noindent\textbf{Abbreviations used in this chapter:} ACL --- access control list; API --- application programming interface; PEL --- pending entries list; SQL --- Structured Query Language; TLS --- Transport Layer Security; TTL --- time to live; AOF --- append-only file; ID --- identifier.\par\smallskip

\index{Redis}\index{Valkey}\index{Fencing}\index{Lease}\index{API!SET NX PX}\index{API!WATCH}\index{API!EVAL}\index{API!XAUTOCLAIM}
\section{Scope and correctness contract}
Redis serializes commands, and scripts, against a single authoritative primary. A key with a finite lifetime is a \emph{lease}, not a kernel mutex. Independent processes, interpreter threads and hosts can compete for the same key, provided they agree on endpoint, database, key namespace and resource identity. A disconnected or paused client can keep executing after its ownership has expired, and asynchronous replication can lose an acknowledged acquisition during failover. Redis coordination is therefore excellent for efficiency: suppressing duplicate work, coalescing cache refreshes, bounding load on a downstream service and distributing jobs. Persistent business invariants belong in database transactions, or behind a fencing check enforced by the protected resource (Section~\ref{sec:fencing}).

Everything in this chapter applies equally to Valkey unless a version note says otherwise.

\subsection{Execution boundary: Perl and Redis}
Application orchestration, retries, deadlines, connections and error handling are Perl. Redis executes its own commands. A short Lua script may be supplied as a Perl heredoc and sent through \code{EVAL} or \code{EVALSHA} for conditional multi-command updates. No separate Lua interpreter is installed by the Perl worker. Redis serializes scripts with respect to other clients: this is isolation, not rollback on runtime errors or durable replication.
\fig{redis-perl-boundary}{Perl owns application logic; Redis owns the atomic state transition. Embedded Lua is data until Redis executes it.}

\section{Server landscape and client selection}\label{sec:redis-clients}
\version{Redis was BSD-licensed through 7.2. Redis 7.4 moved to a dual RSALv2/SSPLv1 license, and Redis 8.0 (May 2025) added AGPLv3 as a third option [40]. Valkey, a Linux Foundation fork of Redis 7.2.4, remains BSD-licensed and is offered by several managed-cloud providers. Both speak the same protocol and support every pattern in this chapter. Confirm the license and support position that applies to your organization before standardizing on either.}
Three Perl clients cover most needs. All three accept arbitrary commands, so newer server commands are available without a client upgrade.
\begin{longtable}{L{.22\linewidth}L{.40\linewidth}L{.30\linewidth}}
\toprule Client & Character & Use when \\ \midrule\endhead
\code{Redis} [21] & Pure Perl; TLS (\code{ssl}, \code{SSL_verify_mode}), ACL \code{username}/\code{password}, Sentinel discovery & Portability matters; moderate request rates \\
\code{Redis::Fast} [41] & XS over hiredis with a \code{Redis}-compatible API; TLS since 0.37, certificate verification on by default since 0.40; Sentinel (not combined with TLS) & Synchronous code that needs throughput \\
\code{Net::Async::Redis} [42] & \code{IO::Async} and \code{Future} based; RESP3, Streams, client-side caching; cluster routing via \code{Net::Async::Redis::Cluster} & Event-loop applications; Redis Cluster \\
\bottomrule\end{longtable}
Whichever client you use, create the connection after \code{fork} or inside the thread, set connect, read and write timeouts, and disable automatic reconnection on coordination connections, or restrict it to idempotent reads. A transparent replay of \code{SET NX} or \code{INCR} after a lost reply can create a lease or consume a fence the caller never learns about.

\section{Command reference}
\begin{longtable}{L{.36\linewidth}L{.56\linewidth}}\toprule Command & Purpose and availability \\ \midrule\endhead
\code{SET key token NX PX ms} & Atomically create an absent key with a millisecond expiry. All supported versions. \\
\code{GET}, \code{PTTL} & Observe owner and remaining time; neither grants ownership. \\
\code{EVALSHA} / \code{EVAL} script & Owner-checked release or renewal, or several steps as one atomic unit. All supported versions. \\
\code{DELEX key IFEQ token} & Native compare-and-delete. Redis 8.4+ [33,34]. \\
\code{SET key token IFEQ token PX ms} & Native owner-checked renewal. Redis 8.4+; Valkey 8.1+ [20,36]. \\
\code{DELIFEQ key token} & Native compare-and-delete. Valkey 9.0+ [35]. \\
\code{INCR} & Fence counter on one primary; see Section~\ref{sec:fencing} for its failure semantics. \\
\code{WAIT}, \code{WAITAOF} & Wait for replica acknowledgment, or (7.2+) for AOF fsync locally and on replicas [37]. \\
\code{XREADGROUP}, \code{XACK}, \code{XAUTOCLAIM} & Consumer-group queues; \code{XAUTOCLAIM} since 6.2 [38]; \code{XREADGROUP ... CLAIM} since 8.4. \\
\code{WATCH}/\code{MULTI}/\code{EXEC} & Optimistic transaction; retry an aborted \code{EXEC}; an ambiguous reply is not a failure. \\
\bottomrule\end{longtable}

\section{Acquisition, renewal and release}
\index{API!SET NX PX}\index{API!EVAL}\index{API!DELEX}
\strategy{redislease}{Redis or Valkey lease}{A parking permit with an expiry time printed on it. Whoever holds a valid permit may park; if they do not renew it in time, it lapses and someone else can take the spot, even if the first driver has not left yet.}
{processes on different hosts must avoid duplicate work (a nightly report, a cache rebuild, one active scheduler) and an occasional overlap is tolerable.}
{an overlap would corrupt data and the protected resource cannot check a fencing token.}
{\code{SET key token NX PX ttl} with owner-checked renewal and release scripts; \code{Redis}, \code{Redis::Fast} or \code{Net::Async::Redis}; example 13.}
Acquire with a cryptographically random token that is unique to each acquisition. Never use \code{SETNX} followed by \code{EXPIRE}: a crash between the two commands leaves an immortal lock. Never release with \code{GET} followed by \code{DEL}: the lease can expire and pass to another client between the two. Renewal and release must compare the owner and act in one server-side step [19,20].
\CodeFile{bash}{snippets/distributed-cache-01.sh}{Single-primary lease operations in a disposable instance}
\fig{leasestates}{Lease lifecycle. Expiry is the dangerous transition: the former holder may still be running.}
The Lua form works on every server and is used in the examples. Prefer the native command where the whole fleet supports it; it is simpler to audit and needs no script cache. Send scripts with \code{EVALSHA} and fall back to \code{SCRIPT LOAD} on a \code{NOSCRIPT} error: the script cache is not persisted and is empty after a restart or failover. Redis Functions (\code{FUNCTION LOAD}/\code{FCALL}, 7.0+) persist and replicate with the dataset, at the cost of a deployment step. Keep scripts O(1) and short; a long-running script blocks every other client.

In Redis Cluster, all keys used by one script must hash to the same slot. A common hash tag (\code{lease:{report}} and \code{fence:{report}}) guarantees this; use a cluster-aware client.

\tip{Name lease keys after what they protect (\code{lock:invoice:42}, not \code{lock:1}) and log the owner token on acquire and release. During an incident, \code{SCAN MATCH lock:*} plus the logs then tell you who holds what, and since when.}

\subsection{Perl-only orchestration with optimistic transactions}\label{sec:redis-watch}
\index{API!WATCH}\index{API!MULTI}\index{API!EXEC}
When native conditional commands are unavailable, Perl can use \code{WATCH}, read the owner, then queue the update with \code{MULTI}/\code{EXEC}. Redis aborts if a watched key changes, including expiry on Redis 6.0.9 or later [61,62]. Retry only confirmed conflicts, with bounded randomized backoff. A lost reply is ambiguous; never assume an ambiguous renewal authorizes additional work.
\begin{longtable}{L{.25\linewidth}L{.32\linewidth}L{.35\linewidth}}\toprule Approach & Benefits & Trade-offs \\ \midrule\endhead
Native conditional command & One command; minimal implementation & Requires appropriate server version \\
Perl with WATCH/MULTI/EXEC & No embedded Lua; optimistic concurrency & More round trips; conflict retries; exclusive connection per attempt \\
Perl invoking Lua & One isolated server operation; multi-key invariants & Script review and ACLs; keys in one cluster slot \\
\bottomrule\end{longtable}
Example 31 demonstrates owner-checked renewal and release and forces a WATCH conflict. It does not allocate fencing tokens or make a lease a correctness guarantee. Use a dedicated connection with reconnection disabled. For multi-key fenced acquisition, watch every key participating in the decision and validate the full result. For more complex token and semaphore state machines, prefer the existing atomic scripts unless there is a compelling reason to redesign them.
\CodeFile{perl}{examples/31-redis-watch-lease.pl}{31-redis-watch-lease.pl -- owner checks without embedded Lua}
\note{Example 31 was documentation-reviewed against the Redis transaction reference and Perl Redis client documentation. It is not part of the historical integration results.}

\section{A fenced lease}\label{sec:fenced-lease}
Example 13 acquires the lease and allocates the fencing token in the \emph{same} script: \code{SET NX PX} and, only if it succeeded, \code{INCR} on the fence key. Because both happen in one atomic step, no other client can acquire in between, and no paused former owner can obtain a newer token. The local validity deadline is measured from \emph{before} the request and reduced by a safety margin. The heartbeat renews after a third of the TTL using the same owner token. Release is owner-checked, and a failed release is reported, never ignored.
\CodeFile{perllua}{examples/13-redis-lease.pl}{\nolinkurl{13-redis-lease.pl}}
\tested{Four processes competed for the lease. Critical sections never overlapped and fences increased strictly in acquisition order. The run was repeated under the restricted ACL user from Section~\ref{sec:redis-security}.}
The client-side deadline is a soft budget: an in-flight command can complete after it, and the margin does not bound scheduling delay or clock error. An exception during acquisition must not authorize work; a key that was possibly created will expire on its own.

\section{Worker integration and heartbeat design}
Create a client inside each thread or after \code{fork}; never use one connection concurrently from unrelated execution contexts. Keep acquisition retries bounded and randomized. If work outlives the TTL, renew before expiry (for example every third of the TTL) with the original token, and stop after a configured maximum total duration. A failed or ambiguous renewal means: stop issuing protected operations and cancel the worker where possible. A heartbeat cannot recall work already sent to another service; that is the job of fencing. An event loop must use an asynchronous client or offload blocking calls.

\section{Owner tokens versus fencing tokens}\label{sec:fencing}
\index{Fencing}
\strategy{fencing}{Fencing token}{Every new permit carries a higher serial number, and the gatekeeper at the protected resource remembers the highest number it has seen. A late arrival waving permit 41 after permit 42 has been used is turned away at the gate, whatever it believes.}
{a lease protects anything where a stale holder's late write would cause harm.}
{the resource cannot store and compare a number atomically; then you need a different design, not a lease.}
{allocate the token in the same atomic step as the lease (example 13) and check it in a guarded \code{UPDATE} (example 18) or a version column.}
\fig{fencing}{Why a lease alone is not enough: client A pauses past its expiry and writes late; the storage rejects its older fence.}
A random owner token identifies an acquisition; it does not establish order. A fencing token is a monotonically increasing generation that the \emph{destination} checks atomically. A resource that admits several writes by the current holder and rejects older generations uses:
\CodeFile{sql}{snippets/distributed-cache-02.sql}{SQL example}
Initialize \code{last_fence} to zero, use a 64-bit unsigned column, bind the same token twice and require one affected row. Example 18 runs the full scenario: client A acquires and writes, stalls past its TTL, client B acquires a newer token and writes, and A's late write and late release are both refused.
\CodeFile{perllua}{examples/18-fenced-write.pl}{\nolinkurl{18-fenced-write.pl}}
\tested{The stale write matched zero rows, A's release returned 0 and left B's lease intact, and the row kept B's payload and fence.}
\warn{A Redis \code{INCR} counter can regress after data loss, an unpersisted restart or asynchronous failover. With destination-enforced fencing this \emph{fails safe}: new tokens fall below \code{last_fence}, so writes are rejected rather than corrupting data. The laboratory reproduced this by deleting the counter. Recovery must re-seed the counter above the destination's maximum (for example \code{SET fence:{report} <max>} before reopening traffic), or allocate generations transactionally in the database that owns the resource.}

\section{Distributed semaphores}
\strategy{redissem}{Distributed counting semaphore}{The car-park counter, kept in Redis so that cars arriving at any of several entrances share it. Each ticket has its own expiry, so a car that broke down inside is eventually removed from the count.}
{several hosts must not exceed a concurrency limit on a shared dependency, such as a licence pool, a fragile service or an API quota.}
{you need requests per second (rate limiting), or a single host suffices (use \code{Thread::Semaphore} or System~V).}
{a sorted set scored by server-time expiry with one Lua script for acquire and one for renew; example 16.}
\fig{redissem}{A distributed semaphore as a sorted set: expired holders are purged atomically on every acquisition.}
A counting semaphore bounds concurrent use of a downstream resource (an API quota, a licence pool, a fragile legacy service) across hosts. Example 16 stores holders in a sorted set scored by expiry time. The acquisition script reads the \emph{server} clock with \code{TIME}, purges expired holders, checks the count and admits the caller, all atomically. A crashed holder's permit therefore returns after its TTL, and no client clock participates. Renewal refuses tokens that have already expired, so a resumed holder cannot silently reclaim a permit that has been reissued.
\CodeFile{perllua}{examples/16-redis-semaphore.pl}{\nolinkurl{16-redis-semaphore.pl}}
\tested{Ten processes competed for three permits; observed peak concurrency was exactly three.}
A concurrency limit is not a rate limit. For operations per unit time, use a token bucket or GCRA implemented in a script, or the server-side features of a managed offering, and test it against bursty arrivals.

\section{Prefer queues: Redis Streams consumer groups}\label{sec:streams}
\index{API!XREADGROUP}\index{API!XACK}\index{API!XAUTOCLAIM}
\strategy{streams}{Redis Streams consumer group}{A shared inbox where each letter is handed to one clerk, stays on that clerk's ``in progress'' tray until they stamp it done, and can be reassigned if the clerk has been gone too long.}
{jobs are produced and consumed by many processes across hosts with low latency, and at-least-once delivery with idempotent effects is acceptable.}
{losing an acknowledged enqueue during a failover is unacceptable; use a transactional database queue instead.}
{\code{XADD}, \code{XREADGROUP}, \code{XACK}, \code{XAUTOCLAIM} through \code{Redis} or \code{Net::Async::Redis}; example 17.}
Many locks exist only to stop two workers taking the same job. A consumer group removes the need: each entry is delivered to one consumer and stays in that consumer's \emph{pending entries list} (PEL) until acknowledged with \code{XACK}. If a consumer dies, its pending entries remain; any other consumer can adopt entries that have been idle longer than a threshold with \code{XAUTOCLAIM} [38]. Delivery is therefore \emph{at least once}: make the effect idempotent, then acknowledge.
\fig{streams}{A consumer group with crash recovery, as exercised by example 17.}
\CodeFile{perl}{examples/17-redis-stream-worker.pl}{\nolinkurl{17-redis-stream-worker.pl}}
\tested{One of four consumers exited holding ten unacknowledged entries; another consumer reclaimed them. All 200 effects were applied exactly once.}
Production practice: cap stream length with \code{XADD ... MAXLEN ~ n} or trim by minimum ID; read delivery counts with the extended form of \code{XPENDING} and move entries that exceed a retry limit to a dead-letter stream; keep the idle threshold well above normal processing time to avoid needless duplicate execution; and alert on PEL size and oldest pending age. Redis 8.4 adds a \code{CLAIM} option to \code{XREADGROUP} that returns idle pending entries and new entries in a single call [34]. Streams inherit Redis durability: they are not a substitute for a transactional queue when loss of an acknowledged enqueue is unacceptable.

\section{Durability, replication, Redlock and cluster topology}
Asynchronous replication allows a primary to acknowledge an acquisition, fail before replicating it and promote a replica without the key; a second client then acquires the same logical lock. Sentinel and Redis Cluster improve availability, not this safety property. \code{WAIT} reduces replication exposure, and \code{WAITAOF} (7.2+) can additionally require an AOF fsync locally and on replicas, but neither turns Redis into a consensus system [37]. With \code{appendonly yes}, \code{appendfsync always} narrows the crash window at a latency cost; \code{everysec} can lose roughly the last one to two seconds of acknowledged writes on a crash.

\fig{failover}{How asynchronous replication loses a lease during failover.}
Key expiry is stored as an absolute Unix time on the server, so a stepped server clock can shorten or extend every lease at once. Run time synchronization that slews rather than steps, and alert on clock corrections.

\fig{redlock}{Redlock acquires on a majority of independent primaries within a validity window; it reduces but does not remove the need for fencing.}
Redlock acquires the same key on a majority of independent primaries (typically five) within a validity window that subtracts elapsed time and a drift allowance from the TTL, releases partial acquisitions on failure, and retries with jitter [19]. Its safety rests on bounded process pauses, bounded clock drift and independent failure domains. This assumption set has been publicly debated [43]; the practical conclusion is that Redlock may reduce, but does not eliminate, the need for destination fencing, and that fencing on a single primary is often the better engineering trade. Never build a quorum from replicas of one primary.

\section{Key design and memory policy}
Every coordination key used in this handbook has an owner, a type and a TTL. A key without a TTL outlives the process that created it, so every write path below sets one atomically.
\begin{longtable}{L{.24\linewidth}L{.10\linewidth}L{.22\linewidth}L{.34\linewidth}}\toprule Key pattern & Type & TTL & Purpose \\ \midrule\endhead
\code{lease:{res}} & string & lease (ms) & Owner token for a fenced lease (13, 18) \\
\code{fence:{res}} & string & none; re-seed after loss & Monotonic fencing counter, same hash slot as the lease \\
\code{sem:{pool}} & zset & latest holder expiry & Distributed semaphore holders scored by expiry (16) \\
\code{jobs:{q}} & stream & trimmed by \code{MAXLEN} & Work queue with consumer group (17) \\
\code{sess:session:SID} & string & session idle timeout & Catalyst session data (Chapter~\ref{ch:catalyst}) \\
\code{cart:{SID}} & hash & session idle timeout & Mutable per-session state outside the blob \\
\code{idem:USER:KEY} & string & 2 min pending, 24 h done & Idempotency claim and stored response \\
\code{ratelimit:login:USER} & string & window & Fixed-window attempt counter \\
\code{slock:{SID}} & string & 10 s & Per-session request lock \\
\bottomrule\end{longtable}
Use hash tags (\texttt{\{...\}}) wherever a script touches more than one key, so that the keys share a cluster slot. Remember that \code{maxmemory-policy} applies to a whole instance, not to a logical database number. Keep leases, fences, queues and sessions on an instance with \code{noeviction}, and put caches on a separate instance with an LRU or LFU policy.

\section{Security and deployment}\label{sec:redis-security}
Use a dedicated instance, or at least a dedicated logical database, for coordination keys, and set \code{maxmemory-policy noeviction} so that memory pressure fails writes instead of silently deleting leases. Expiration, restart and failover remain possible even then. Treat out-of-memory errors and network errors as acquisition failure. Never expose the service publicly; keep \code{protected-mode} enabled, authenticate every client, and use TLS with certificate verification (\code{ssl => 1} and \code{SSL_verify_mode} in \code{Redis.pm}) on any network you do not fully control.

Since Redis 6, commands executed inside scripts are checked against the calling user's ACL, so a least-privilege lease user needs both \code{EVALSHA} and the commands its scripts call [39]:
\CodeFile{bash}{snippets/distributed-cache-03.sh}{Least-privilege ACL for examples 13 and 16 (tested)}
\tested{Under this user, examples 13 and 16 passed, while \code{FLUSHALL} and writes to keys outside the granted patterns were refused with \code{NOPERM}.}
Persist ACLs through \code{aclfile} or configuration management rather than ad hoc \code{ACL SETUSER} calls, and rotate secrets without restarts by temporarily granting two passwords.

\section{Monitoring and testing}
Track acquisition wait, contention, lost and failed renewals, release failures, rejected fences, expired and evicted keys (\code{INFO stats}), rejected commands, replication lag, failover events and \code{SLOWLOG} entries for scripts. Name connections (\code{CLIENT SETNAME}) so that \code{CLIENT LIST} identifies holders. Test process death, \code{kill -STOP} followed by resume after the TTL, loss of replies, replica promotion, restart without persistence (fence regression) and clock steps. Verify that an old owner cannot release or renew a successor's lease and that the destination rejects stale generations.



% ================================================================ chapter 10
\chapter{Memcached locking and cache coordination}
\chapterbrief{Use cache coordination for efficiency without confusing it with durable authority.}
\par\smallskip\noindent\textbf{Abbreviations used in this chapter:} CAS --- compare-and-swap; TLS --- Transport Layer Security; TTL --- time to live; UDP --- User Datagram Protocol.\par\smallskip

\index{Memcached}\index{Cache stampede}\index{API!CAS}
\section{What Memcached can and cannot guarantee}
Memcached is a volatile cache with atomic per-item operations, not a durable coordination authority. \code{add} stores a value only when the key is absent; \code{cas} updates an item only if its version is unchanged. These operations can suppress a cache stampede across threads, processes and hosts. Eviction, restart, flush, expiry or a change in client routing can erase a claim while its holder is still working. Never use Memcached as the sole mechanism for financial correctness, unique job execution, inventory allocation or irreversible effects [22,23].
\begin{longtable}{L{.24\linewidth}L{.68\linewidth}}\toprule Primitive & Semantics \\ \midrule\endhead
\code{add} & Create only when absent, with a finite expiry in seconds. \\
\code{gets} & Read the value and its opaque CAS identifier. \\
\code{cas} & Replace if the observed identifier is still current; can set a new expiry. \\
\code{delete} & Unconditional removal in the classic protocol; unsafe for lease release. \\
\code{touch} & Unconditional expiry change; unsafe for owner-checked renewal. \\
\code{ms key 3 ME} & Meta-protocol add (mode \code{E}) of a 3-byte value. \\
\code{md key C123} & Meta-protocol conditional delete: succeeds only if the current CAS is 123. \\
\bottomrule\end{longtable}
Use short relative TTLs. Zero means no expiry; values above 30 days are interpreted as absolute Unix timestamps. Expiry has one-second granularity, so a TTL of 1 can expire almost immediately. Use one fixed server or an identical, stable routing configuration across all workers: rehashing keys during a failure creates independent ownership domains. Never treat a cache miss as proof that previous work has stopped.

\section{Classic text protocol and Perl CAS release}
\index{API!CAS}\index{API!Cache::Memcached::Fast}
\strategy{memcas}{Memcached add/CAS lease}{A sticky note on a shared whiteboard saying ``I'm rebuilding this''. Anyone can see it and back off, but the cleaner (eviction) may wipe the board at any time.}
{many web workers would otherwise rebuild the same expensive cache entry at once.}
{anything other than cache efficiency depends on it.}
{\code{add} to claim, \code{gets}/\code{cas} to release with a tombstone (\code{Cache::Memcached::Fast}, example 14); meta-protocol conditional delete (transcript 15).}
\fig{memcas}{A best-effort Memcached lease: claim with add, release by compare-and-swap so another owner's entry is never touched.}
The classic protocol lacks conditional delete. After reading and verifying the owner with \code{gets}, replace the entry with a short-lived tombstone using \code{cas}. If the version has changed, delete or overwrite nothing. Leave the tombstone to expire, since deleting it would recreate the original race. This delays the next acquisition by up to the tombstone TTL. Alternatively, let the lease expire. CAS identifiers are not fencing tokens and can repeat across server restarts.
\CodeFile{perl}{examples/14-memcached-lease.pl}{\nolinkurl{14-memcached-lease.pl}}
\code{Cache::Memcached::Fast} returns an array reference of CAS and value from \code{gets}; \code{add} and \code{cas} return false for a negative reply and \code{undef} for a backend error [24]. The example counts backend errors separately and gives up after repeated failures rather than treating an outage as permanent contention. Renew by fetching owner and CAS, verifying the token and performing \code{cas} with the original token and a new TTL; never use \code{touch} or \code{gat} to extend a lease that might belong to someone else.
\tested{Four competing processes were serialized. Treat this as evidence that the protocol is implemented correctly, not as a mutual-exclusion guarantee.}

\section{Modern meta protocol}
\index{API!Memcached meta protocol}
The meta protocol provides conditional storage, conditional deletion and built-in stampede control [25,26]. Use a client that supports meta commands, or a carefully written adapter; do not mix raw protocol values with a client's serialization without checking flags and encoding. In the transcript, \code{ME} selects add mode, \code{C} supplies an expected CAS and \code{T} a TTL. \code{HD} means success, \code{NS} not stored, \code{EX} a CAS mismatch and \code{NF} not found.
\CodeFile{text}{examples/15-memcached-meta.txt}{Meta protocol acquisition, renewal and release (replayed against 1.6.24)}
Every successful mutation changes the CAS, so read the owner and the current CAS together immediately before a conditional delete; a mismatch means relinquish the attempt. The CAS check prevents the ordinary expiry-and-reacquisition race on one server incarnation, not eviction or server-reset failures.

\section{Cache stampede pattern}
\strategy{stampede}{Cache stampede control}{When the popular dish runs out, one cook makes a new batch while the others keep serving yesterday's portion, instead of every cook starting the same recipe at once.}
{a hot cache key is expensive to rebuild and its expiry sends many requests to the database at once.}
{the data must never be stale; then the source must be read directly, behind proper capacity limits.}
{Memcached meta flags \code{N}, \code{R}, \code{W}, \code{Z}; or a short Redis lease around the rebuild with stale-while-revalidate.}
\fig{stampede}{Without control, an expiring hot key triggers a rebuild per request; with recache election, one request rebuilds while others serve the stale value.}
Read the cache first. On a miss, compete for a short lease with \code{add}. The winner rechecks the cache, performs a bounded refresh, publishes the result and releases conditionally. Losers sleep with jitter and reread, serve permitted stale data, or return a controlled error after a deadline. Do not let every waiter query the database when the cache goes down: use a local semaphore, request coalescing, stale-while-revalidate or a circuit breaker to cap backend load. If a stale refresh could overwrite a newer result, derive versions from authoritative rows and write conditionally.

The meta protocol builds this in. On a miss, \code{mg} with \code{N}\textit{ttl} creates a placeholder and returns \code{W} to exactly one client, which wins the recache; others see \code{Z} (a recache is already in progress). \code{R}\textit{ttl} grants an early recache to one client when the remaining TTL falls below a threshold, and \code{md} with \code{I} marks an item stale instead of removing it, so readers receive the old value with \code{X} and \code{W} while one refresh runs. Prefer these flags for cache-specific coordination when the client supports them; they are still not durable locks.

\tip{Add random jitter to cache TTLs (for example $\pm$10\,\%). Keys filled together by a deploy or a warm-up job otherwise expire together and stampede together.}

\section{Operational comparison}
\begin{longtable}{L{.20\linewidth}L{.36\linewidth}L{.36\linewidth}}\toprule Property & Redis / Valkey & Memcached \\ \midrule\endhead
Acquisition & \code{SET NX PX} (optionally with fence \code{INCR} in one script) & \code{add} / meta add \\
Expiry & Millisecond lease & Second-granularity lease \\
Release & Owner-checked script; \code{DELEX IFEQ} (8.4+); \code{DELIFEQ} (Valkey 9.0+) & CAS tombstone or meta conditional delete \\
Renewal & Owner-checked script; \code{SET IFEQ PX} & Owner-checked CAS update \\
Fencing & Atomic \code{INCR} in the acquire script; can regress & None \\
Queues & Streams with consumer groups & None \\
Durability & Configurable; failover may lose writes & Volatile by design \\
Best use & Best-effort coordination, bounded admission, job distribution & Cache-refresh coalescing \\
Critical writes & Require destination fencing & Use another authority \\
\bottomrule\end{longtable}
Monitor evictions, reclaimed items, rejected stores, timeout rates and hit ratio. Restrict access to private trusted networks, disable UDP unless explicitly needed (\code{-U 0}), and use SASL or TLS where the build supports them. On a disposable instance, test concurrent \code{add}, expiry and reacquisition, stale CAS, memory pressure, restart and routing changes. Never run \code{flush_all} against a production cache as a test.

\needspace{12\baselineskip}\section{Laboratory setup and verification}
\CodeFile{bash}{snippets/distributed-cache-04.sh}{Shell commands}
Each service example forks its own competing workers and checks the result, so no manual orchestration is required. To explore failure modes, increase the work duration toward the TTL, pause a holder with \code{kill -STOP}, or restart Redis between runs of example 18 and observe fencing fail safe.

\chapter{Failure boundaries and authoritative coordination}\label{ch:failure-boundaries}
\chapterbrief{Separate safety from progress; choose a coordinator with the required failure contract.}
\par\smallskip\noindent\textbf{Abbreviations used in this chapter:} HTTP --- Hypertext Transfer Protocol; TTL --- time to live; ID --- identifier.\par\smallskip

\index{Consensus}
\index{Leader election}
\index{Fencing}
\index{Failure boundaries}

\section{Lock hierarchy and failure boundaries}
An application that holds a local mutex while waiting for a database lock may deadlock with another path that holds the database lock while waiting for the local mutex. Establish a cross-layer order, or design operations so locks at different layers are never held simultaneously. Make acquisition order a documented invariant, and prefer resource ownership and message passing when nested locking becomes difficult to reason about.
\begin{longtable}{L{.25\linewidth}L{.67\linewidth}}\toprule Failure & Required design response \\ \midrule\endhead
Worker exception & Scope-based release or explicit cleanup; propagate the result \\
Process termination & Kernel locks release; repair partially written state \\
Database disconnect & Session locks vanish; transaction outcome may be unknown \\
Process pause (GC, swap, SIGSTOP) & Lease expires while the holder still runs; fence writes \\
Network partition & Lease ownership diverges from client belief; fence writes \\
Coordinator data loss or failover & Lease and fence counter may regress; destination fencing fails safe \\
Queue overload & Bound admission, reject or shed load, alert on oldest-item age \\
Deadlock & Detect or time out, roll back, retry only safe operations \\
\bottomrule\end{longtable}
\section{Distributed leases and fencing}
A time-limited lease is not a mutex with perfect failure detection. A paused worker may resume after its lease expired and a new owner was elected. Assign monotonic fencing tokens and require the protected service to reject older tokens atomically, as examples 12 and 18 do. Release must verify the owner token atomically. Prefer database constraints when the database itself is the protected resource.
\section{Authoritative coordinators}\label{sec:authorities}
When correctness spans hosts and no shared database is available, use a coordinator built on consensus, and still fence at the destination.
\strategy{consensus}{Consensus-backed coordinator}{Instead of one receptionist who might forget, a committee of three or five that only acts when a majority agrees, and that numbers every decision in order. The protected service accepts orders only with a number higher than any it has seen.}
{a violation would cause real harm (money, filings, data loss) and there is no shared database to enforce constraints.}
{a database constraint can do the job, or the cost of running and operating the cluster outweighs the risk.}
{etcd, ZooKeeper or Kubernetes Lease objects through their HTTP APIs; use a generation only if its monotonicity and recovery contract are established; Kubernetes leader election alone is not fencing.}
\fig{consensus}{A consensus coordinator supplies a monotonic revision; the protected service enforces it.}
\begin{itemize}
\item \textbf{etcd.} Leases with TTL and keep-alive; keys attached to a lease vanish when it expires. Each key's \code{mod_revision} increases monotonically across the cluster and serves as a fencing token [45].
\item \textbf{ZooKeeper.} Ephemeral sequential znodes implement lock queues without polling, and the transaction ID (zxid) orders updates.
\item \textbf{Kubernetes Lease objects} (\code{coordination.k8s.io/v1}). Used for controller leader election, with holder identity, duration and renew time updated through optimistic concurrency on \code{resourceVersion}. \code{leaseTransitions} counts changes in the current record, not a durable monotonic generation. Deletion or restoration can reset history. Kubernetes leader election does not fence old leaders; obtain a durable generation from the protected resource or another authority and enforce it there [44,74].
\item \textbf{PostgreSQL advisory locks.} \code{pg_advisory_xact_lock} releases automatically at commit or rollback, unlike MySQL's session-scoped \code{GET_LOCK}; the session-scoped variants share MySQL's disconnect hazards [47].
\end{itemize}

\section{Failure modes: progress and correctness}
\index{Deadlock}\index{Livelock}\index{Starvation}\index{Fencing}
\begin{figure}[H]\centering
\begin{tikzpicture}[x=1cm,y=1cm]
\node[box,text width=4.6cm] at (0,1.8) {Deadlock\\A holds X, waits for Y; B holds Y, waits for X};
\node[box,text width=4.6cm] at (6.1,1.8) {Starvation\\Other owners keep winning while one waits};
\node[badbox,text width=4.6cm] at (0,0) {Livelock\\Both retry repeatedly; neither completes};
\node[badbox,text width=4.6cm] at (6.1,0) {Stale ownership\\A resumes after B takes over};
\end{tikzpicture}
\caption{Deadlock and livelock prevent collective progress; starvation prevents individual progress; stale ownership threatens correctness.}
\end{figure}
Use consistent lock order for deadlock prevention, bounded retry and jitter
for livelock, explicit fairness policies for starvation, and destination
fencing for stale ownership. Test safety separately from liveness.
\begingroup\small\renewcommand{\arraystretch}{1.02}\begin{longtable}{L{.20\linewidth}L{.22\linewidth}L{.22\linewidth}L{.22\linewidth}}
\toprule Failure & InnoDB transaction & Redis lease & Memcached claim\\ \midrule\endhead
Worker exits & Uncommitted work rolls back & TTL or release & Expiry or eviction\\
Worker pauses & May retain locks & May expire & May expire or be evicted\\
Response lost & Commit uncertain & Command uncertain & Command uncertain\\
Failover & Configuration matters & Lease history may regress & State may disappear\\
Duplicate effects & Constraints and transactions & Resource guard required & Not an authority\\
\bottomrule
\end{longtable}\endgroup

\partstart{V / Application architecture}{Application architecture and practical scenarios}{Combine the mechanisms into object-oriented applications, web requests and recoverable processing pipelines.}
\chapter{Concurrency in object-oriented Perl applications}
\label{ch:enterprise-objects}
\chapterbrief{Separate application, worker, request and durable resources.}
\index{Moose}
\index{Catalyst}
\index{Fork}
\index{Resource ownership}


A production Perl application spans many lifetimes: interpreter, worker,
request, transaction and durable service. In a synthetic inventory application,
a web controller, an import job and a notification worker may all operate on
the same record. A lock attached to a Moose object cannot coordinate
independent processes. Choose the owner of the invariant first.

\strategy{app-objects}{Worker-owned resources}{Each worker has its own toolbox.}{live handles must remain process-local.}{state must coordinate independent workers.}{Moose, DBI, Catalyst}

\section{State ownership and object lifetime}
\begin{figure}[!htbp]\centering
\begin{tikzpicture}[node distance=6mm and 12mm]
\node[box,text width=6cm] (p) {Parent: immutable configuration and class metadata};
\node[okbox,below left=of p,text width=3.5cm] (a) {Worker A\\PID-owned clients};
\node[okbox,below right=of p,text width=3.5cm] (b) {Worker B\\PID-owned clients};
\draw[arr](p)--(a); \draw[arr](p)--(b);
\node[box,below=of a,text width=3.5cm] (ra) {Request context A};
\node[box,below=of b,text width=3.5cm] (rb) {Request context B};
\draw[arr](a)--(ra); \draw[arr](b)--(rb);
\node[okbox,below=8mm of $(ra.south)!0.5!(rb.south)$,text width=6cm] (db) {Shared authoritative database};
\draw[arr](ra.south)--(db.north west); \draw[arr](rb.south)--(db.north east);
\end{tikzpicture}
\caption{Each process owns its clients; request contexts are transient and shared invariants are enforced at the resource.}
\end{figure}

\begin{longtable}{L{.22\linewidth}L{.33\linewidth}L{.37\linewidth}}
\toprule Lifetime & Appropriate state & Avoid \\ \midrule\endhead
Application & Immutable configuration, class metadata & Current user, mutable totals \\
Worker & Connection factories, bounded pools, metrics & Assuming other workers see local changes \\
Request / job & Identity, operation ID, validated input & Storing request context on a reusable singleton \\
Durable service & Inventory, claims, checkpoints, audit rows & Treating a cache as authoritative \\
\bottomrule
\end{longtable}

Moose lazy attributes defer construction; they are not synchronization
primitives [63]. Catalyst component factories and context-aware lookup
(\code{COMPONENT} and \code{ACCEPT_CONTEXT}) allow different lifetimes [64].
Pass request-specific values explicitly instead of storing the active user
or current transaction on a persistent controller. A process-local singleton
does not become a deployment-wide singleton. Interpreter threads clone ordinary
state; an event-loop task can race another task across an \code{await} even
without parallel CPU execution.

\section{Resource ownership after fork}
Prefer creating live clients after the server forks. Use documented
driver-specific fork handling and reconnect behavior. A PID check detects
inherited state, but cannot make an inherited socket safe to use or destroy.
DBI's \code{AutoInactiveDestroy} protects against accidental destructor
interference in a forked child [15]. It does not make a shared connection safe.

\warn{Never use an inherited database handle concurrently in a parent and
child. Do not call a protocol-level disconnect on an inherited client
unless its documentation establishes that this is safe.}

The following Moose laboratory uses a plain hash to model resource ownership,
not a real database connection. The child must initialize its own resource
and leave the parent unchanged.
\CodeFile{perl}{examples/32-worker-lifecycle.pl}{Example 32: process-local Moose resource ownership}

\section{Reusable services and invariants}
Give each service operation explicit inputs, a stable operation ID and
well-defined results: success, conflict, validation failure or unknown
outcome. Web, batch and administrative entry points should use the same
business rules. Prefer a resource-enforced guard or unique constraint so
that an alternate writer cannot silently bypass an advisory lock.


% ================================================================ Catalyst chapter
\chapter{Catalyst sessions and request coordination on Redis}\label{ch:catalyst}
\chapterbrief{Understand request-local state, concurrent sessions and safe web operations.}
\par\smallskip\noindent\textbf{Abbreviations used in this chapter:} ACL --- access control list; API --- application programming interface; CSRF --- cross-site request forgery; CSPRNG --- cryptographically secure pseudorandom number generator; DBI --- Database Interface; HTTP --- Hypertext Transfer Protocol; HTTPS --- Hypertext Transfer Protocol Secure; JWT --- JSON Web Token; PID --- process identifier; TLS --- Transport Layer Security; TTL --- time to live; AOF --- append-only file; ID --- identifier.\par\smallskip

\index{Catalyst}
\index{Sessions}
\index{Idempotency}
\index{API!Catalyst::Plugin::Session}

A Catalyst application in production is never one interpreter. Starman, Gazelle or uWSGI fork several workers per pod, and Kubernetes runs several pods (Figure~\ref{fig:prefork}). Any request can land on any worker, so sessions, rate limits, duplicate-submission guards and per-user locks must live outside the interpreter. Redis or Valkey is the usual home for that state. This chapter explains how \code{Catalyst::Plugin::Session} uses its store, why the commonly packaged Redis stores need hardening, and how to handle the concurrency problems that sessions create. It also covers idempotency keys, login rate limits, session revocation and, for API clients, hashed API keys, JWT access tokens, refresh-token rotation and signed requests described with OpenAPI~3.2.

Everything here is exercised by a complete application in \code{examples/catalyst-session/} and by two drivers, \code{examples/28-catalyst-sessions.pl} (29 checks) and \code{examples/30-api-tokens.pl} (33 checks). Each starts the application under Starman with four workers and checks its properties against live Redis (Section~\ref{sec:catalyst-lab}).

\section{How Catalyst::Plugin::Session uses its store}
The plugin keeps three records per session, each keyed by the session ID: \code{session:SID} holds the data hash, \code{expires:SID} the expiry time in epoch seconds, and \code{flash:SID} the flash data. The session state plugin (normally \code{Session::State::Cookie}) only carries the ID; everything else is in the store.

Loading is lazy and saving is all-or-nothing (Figure~\ref{fig:sessionlifecycle}). The first call to \code{$c->session} reads \code{session:SID} and records an \code{Object::Signature} digest of the data. In \code{finalize_body}, before the response is sent, the plugin recomputes the digest and, if anything changed, writes the \emph{entire} hash back. Expiry is extended on each request, which writes \code{expires:SID}; the \code{expiry_threshold} option limits those writes to once per threshold interval.
\fig{sessionlifecycle}{One request through Catalyst::Plugin::Session. Nothing is locked between the load and the save.}
Three consequences follow directly from this design:
\begin{itemize}
\item Every request that changes the session costs a full serialize-and-write of the session, so keep sessions small.
\item Two concurrent requests in the same session race: the last writer wins and the other's changes vanish silently (Section~\ref{sec:session-race}).
\item Everything stored must serialize cleanly. Objects, handles and closures do not belong in a session.
\end{itemize}

\section{Choosing a session store}
Store names below omit the common prefix \code{Catalyst::Plugin::Session::}.
\begin{longtable}{L{.26\linewidth}L{.66\linewidth}}\toprule Store & Assessment \\ \midrule\endhead
\code{Store::FastMmap}, \code{Store::File} & Single host only. With more than one pod, users are logged out whenever a request reaches another host. \\
\code{Store::DBIC}, \code{Store::DBI} & Durable and shared, but every request with a changed session writes to the primary database; an expired-session sweeper is required. \\
\code{Store::Redis} 0.900 & Shared and fast, but: writes \code{SET} then \code{EXPIRE} as two commands; deserializes with Storable \code{thaw}; sends \code{PING} before every operation; keeps one class-wide connection; and its configuration accepts only server, database, debug and reconnect, with no password, ACL user or TLS. \\
\code{Store::Cache} with \code{CHI::Driver::Redis} 0.10 & Accepts Redis options (credentials, TLS) and a configurable serializer, but also writes \code{SET} then \code{EXPIRE}. It adds every key to a namespace set that is never trimmed when keys expire, so that set grows without bound. \\
Hardened custom store (this chapter) & One atomic \code{SET ... EX}, Sereal with objects refused, fork-safe connection with ACL and TLS, no per-call \code{PING}. About 50 lines. \\
\bottomrule\end{longtable}
\version{Assessments are from the source of the Ubuntu 24.04 packages: \code{Catalyst} 5.90131, \code{Catalyst::Plugin::Session} 0.43, \code{Session::State::Cookie} 0.18, \code{Store::Redis} 0.900 and \code{CHI::Driver::Redis} 0.10. Check for newer CPAN releases before adopting any of them, and re-read their \code{store_session_data} before trusting the changelog.}
\warn{Storable \code{thaw} can instantiate arbitrary Perl classes and trigger their hooks and destructors. Any process that can write to the session database can then influence code paths in every application worker. Treat write access to session keys as equivalent to code execution unless the store refuses objects.}

\section{A hardened Redis store}
\index{API!Catalyst::Plugin::Session}
The store below implements the four methods that \code{Catalyst::Plugin::Session::Store} requires. Each write sets the value and its TTL in one command, so a crash can never leave a session that never expires. Records are encoded with Sereal, and the decoder refuses blessed data, so a forged record is discarded and logged rather than revived. Size limits bound the cost of a hostile payload.
See the complete implementation in Appendix~\ref{app:catalyst}.

The connection helper creates one Redis handle per process, lazily, and creates a new one when the PID changes. A preforking server can therefore load the application in its master process without the workers sharing an inherited socket, whose replies would interleave between processes. The helper disables automatic reconnection, so a lost reply is never followed by a silent replay of \code{SET NX} or \code{INCR}, and it caches scripts with an \code{EVALSHA} fallback.

See the complete implementation in Appendix~\ref{app:catalyst}.
Refusing objects has a design benefit beyond security: sessions hold plain data only. Store an epoch or an ISO date instead of a \code{DateTime}, and a user ID instead of a \code{DBIx::Class} row.

\section{Session security}
\subsection*{Session identifiers}
The plugin's default \code{generate_session_id} hashes a counter, the time, Perl's \code{rand}, the PID and a memory address. None of these is a cryptographic source, and perl5440delta repeats that \code{rand} is unsuitable for security [46]. The laboratory application overrides it with 256 bits from the kernel CSPRNG, and uses the same \code{random_hex} helper for CSRF tokens and lock tokens. Hex output satisfies the plugin's \code{validate_session_id}.

\subsection*{Fixation}
An identifier issued before authentication must not survive it, or an attacker who planted it can ride the victim's login. Call \code{$c->change_session_id} at login and at every privilege change; it copies the data to a fresh ID and deletes the old records.

\subsection*{Cookies and proxies}
\code{Session::State::Cookie} 0.18 defaults to \code{HttpOnly} and \code{SameSite=Lax}, but \code{cookie_secure} defaults to 0, so the cookie also travels over plain HTTP. Set it to 1 in production. Mode 2 (secure only for HTTPS requests) depends on \code{$c->req->secure}, which is false behind a TLS-terminating proxy unless \code{Plack::Middleware::ReverseProxy} is active. Catalyst's \code{apply_default_middlewares} enables that middleware only for requests whose \code{REMOTE_ADDR} is 127.0.0.1. Behind a Kubernetes ingress running in another pod, it is therefore off unless \code{using_frontend_proxy} is configured. Once on, \code{X-Forwarded-*} headers are trusted, so make sure only the proxy can reach the application.

\subsection*{Contents, lifetime and revocation}
Keep secrets out of the session. Pair the idle timeout (\code{expires}) with an absolute lifetime enforced by the application, for example a \code{created} timestamp checked in \code{auto}. IP binding (\code{verify_address}) logs out mobile and NAT users and is rarely worth it. To revoke every session for a user after a password change, keep an index set of that user's session IDs (\code{user-sessions:USER}) and delete their records. The laboratory's \code{logout_everywhere} action does exactly this.

\section{Concurrent requests in one session}\label{sec:session-race}
\index{Lost update}
\strategy{session}{Per-session coordination}{A shared shopping list on the fridge. If two people each take it down, add an item at their desk and pin their copy back up, the second copy hides the first person's addition. Either write each item directly on the list in one stroke (atomic), or let only one person hold the list at a time (lock).}
{a web session receives parallel requests that change session state: tabs, background fetches, double clicks.}
{the session only identifies the user and holds rarely changed data; then nothing needs coordinating.}
{atomic Redis structures outside the session (\code{HINCRBY}, \code{SADD}), or \code{lock_session}/\code{finalize_session} in the application class; example 28.}
Browsers issue parallel requests: several tabs, background fetches, a double-clicked button. With the plugin's load-modify-save cycle, two requests that both change the session overwrite each other (Figure~\ref{fig:sessionrace}). Nothing fails and nothing is logged.
\fig{sessionrace}{Two workers serving the same session: the later save silently discards the earlier change.}
The laboratory sends eight simultaneous requests in one session to four Starman workers, each adding one item to a cart, with three implementations:
\begin{longtable}{L{.22\linewidth}L{.20\linewidth}L{.50\linewidth}}\toprule Pattern & Result & Trade-off \\ \midrule\endhead
Naive: cart in the session blob & 3 of 8 kept & Silent data loss under ordinary browser behavior \\
Atomic: Redis hash outside the session, \code{HINCRBY} & 8 of 8 & No waiting; the data needs its own key with a TTL tied to the session \\
Per-session lock around load and save & 8 of 8 & Requests in the session run one at a time (0.94\,s for 8 $\times$ 0.1\,s) \\
\bottomrule\end{longtable}
\fig{sessionfixes}{The three cart implementations measured by the laboratory.}
\note{Default to the atomic pattern: the session holds identity, a CSRF token and small preferences, and anything that several requests change goes into its own Redis structure (hash, set, counter) or into the database, updated with single atomic commands. Use the per-session lock only for routes that must change the session blob itself, and keep those requests short.}
The lock must cover both the load and the save. \code{lock_session} is called before the action's first \code{$c->session} access, and the application overrides \code{finalize_session} to release the lock \emph{after} the plugin has written the session. The lease has a TTL longer than any request, and release is owner-checked, as in Chapter~\ref{ch:redis}. If a worker dies while holding the lock, it expires on its own.
See the complete implementation in Appendix~\ref{app:catalyst}.

The same file sets \code{STDERR->autoflush(1)} after \code{setup}. Catalyst's default UTF-8 encoding applies \code{binmode(STDERR, ':encoding(UTF-8)')}, which makes STDERR \emph{buffered}. Without the fix, worker log lines reached the container log late, in batches, and were lost when a worker was killed. The laboratory's log assertion failed until this line was added.

\section{Idempotency keys for unsafe requests}
\index{Idempotency}
\strategy{idempotency}{Idempotency key}{A cloakroom ticket for a request. Hand in the same ticket twice and you get the same coat back, not a second coat. While the attendant is still fetching your coat, a second ticket holder is told to wait.}
{a client might send an unsafe request more than once: payments, orders, filings, account changes.}
{the operation is naturally idempotent (setting a value, deleting by ID) or read-only.}
{\code{SET key pending NX EX 120}, then store the response with \code{SET XX EX 86400}; scope keys per user and fingerprint the body; example 28.}
A payment, order or filing must not execute twice because a user double-clicked or a client retried after a timeout. The client sends an \code{Idempotency-Key} header. The server claims that key atomically with \code{SET NX}, performs the side effect once, and stores the response so that duplicates receive the same answer (Figure~\ref{fig:idemstates}).
\fig{idemstates}{Idempotency-key lifecycle. Scope keys per user, and fingerprint the request so a reused key with a different body is rejected.}
The details matter:
\begin{itemize}
\item \textbf{Two TTLs.} A pending claim expires after two minutes; a completed one is kept for 24 hours. With a single 24-hour TTL, a worker killed between the side effect and storing the response would block every retry for a day.
\item \textbf{Unknown outcomes stay blocked.} The laboratory kills a worker with \code{SIGKILL} right after the charge. The client sees an error, the charge has happened, and an immediate retry receives 409 rather than a second charge. Pass the same key to the payment provider so that its own deduplication covers the window after the claim expires.
\item \textbf{Failure before the effect} deletes the claim, so the client can retry at once.
\end{itemize}
\tested{Six simultaneous submissions produced exactly one 201 and five 409s. A later retry replayed the stored 201 with \code{Idempotent-Replayed: true}, and the customer was charged once. A reused key with a different amount received 422. After the injected crash, the retry received 409, no second charge occurred, and the pending key expired 120 seconds later.}

\section{Login rate limiting}
\strategy{ratelimit}{Rate limiting}{A turnstile that lets five people through per minute, with a counter that resets when the minute is up.}
{you need to cap attempts or requests per time window: logins, password resets, expensive endpoints, outbound API quotas.}
{you want to bound how many things run \emph{at the same time}; that is a semaphore (Section~\ref{str:semaphore}).}
{an atomic \code{INCR} plus first-time \code{PEXPIRE} in one Lua script (fixed window, example 28); a GCRA script when bursts at window edges matter.}
\fig{ratelimit}{A fixed-window rate limit: the counter and its expiry are created together, so the window always ends.}
The login action counts attempts per username in a fixed one-minute window. \code{INCR} and the first \code{PEXPIRE} run in one script, so a crash can never leave a counter without a TTL, which would amount to a permanent lockout. In production, limit per client address as well as per account, return \code{Retry-After}, and consider GCRA or a sliding window if bursts at window boundaries matter. Account-based lockout can be used to lock out a victim deliberately, so prefer slowing attempts down to hard lockouts for important accounts.

\section{API concurrency and security}
Refresh rotation, nonce consumption, rate limiting and revocation are
concurrent state transitions. Appendix~\ref{app:api} explains their
security contracts and server-side operations.

\section{Deployment on Starman and Kubernetes}
\fig{sessionarch}{Stateless application pods with shared session and coordination state. Session and cache data are split by eviction policy.}
\begin{itemize}
\item \textbf{No sticky sessions.} With a shared store, any worker in any pod can serve any request.
\item \textbf{Connections after fork.} With \code{--preload-app}, the application loads once in the master and is shared copy-on-write. Nothing may open a connection at load time; per-process lazy handles (\code{SessionLab::Redis}) make this safe.
\item \textbf{Separate instances by eviction policy.} \code{maxmemory-policy} applies to a whole Redis instance, not to a logical database number. A cache wants \code{allkeys-lru}. Sessions and coordination keys want \code{noeviction}, because evicting a session logs a user out and evicting a lease breaks mutual exclusion.
\item \textbf{Failover semantics.} Asynchronous replication can lose the last writes. For sessions this means a few logouts or a lost cart change, which is acceptable. For leases it is a correctness issue (Chapter~\ref{ch:redis}). Enable AOF with \code{everysec} on the session instance.
\item \textbf{Security.} Restrict the application's ACL user to its key patterns (\code{sess:*}, \code{cart:*}, \code{idem:*}, \code{ratelimit:*}, \code{user-sessions:*}, \code{slock:*}, and for the API \code{apikey:*}, \code{user-apikeys:*}, \code{rt:*}, \code{rtfam:*}, \code{jwt:deny:*}, \code{nonce:*}), and use TLS outside a trusted network.
\item \textbf{Monitoring.} \code{evicted_keys} on the session instance must stay at zero. Also watch memory per session, session count, lock waits, and 409, 429 and 503 rates. For the API, alert on refresh outcomes: a rising \code{reuse} count means stolen tokens or a client that does not refresh single-flight.
\end{itemize}

\section{Laboratory guide}
The application and verification drivers are collected in Appendix~\ref{app:catalyst}.
Read the session storage and atomic update patterns here first, then use
the appendix as a source reference.

\chapter{Building concurrent Catalyst applications with Memcached}
\label{ch:memcached-capstone}
\chapterbrief{Separate disposable cache coordination from durable authorization and transactional invariants.}
This companion to the Redis Catalyst laboratory uses four Starman workers, Memcached and InnoDB. Memcached supplies bounded CAS retries and best-effort rate counters. The database owns bearer-token authorization and idempotent order records. No cache entry is authoritative for authentication or financial effects.

\begin{figure}[H]\centering
\begin{tikzpicture}[node distance=9mm and 12mm]
\node[box] (http) {Concurrent HTTP clients};
\node[box,below=of http] (app) {Catalyst / four Starman workers};
\node[box,below left=of app,text width=45mm] (cache) {Memcached\\CAS and rate counters};
\node[okbox,below right=of app,text width=45mm] (db) {InnoDB\\Tokens and unique orders};
\draw[arr] (http)--(app);
\draw[arr] (app)--(cache);
\draw[arr] (app)--(db);
\end{tikzpicture}
\caption{Disposable acceleration and durable authority are separate resources.}
\end{figure}

\section{Authentication and the API contract}
The application serves its OpenAPI 3.1 contract at \code{/openapi.json}. Opaque 256-bit bearer credentials are provisioned by the laboratory driver; only SHA-256 digests, scopes and expiration times are stored in the database. Read credentials cannot mutate counters or orders. Revocation removes the credential from the authoritative store. A request already authorized before revocation may complete. Use TLS outside the loopback laboratory and never log bearer credentials. This example does not implement an OAuth authorization server or browser cookie sessions.

\section{CAS and durable idempotency}
The counter reads a CAS token and attempts a conditional replacement. Contention triggers a bounded retry with jitter; an absent key is initialized with add-if-absent. Eviction can reset the counter, so it is not a durable total. The rate counter is likewise advisory and can reset on eviction. An unavailable cache returns an error rather than silently bypassing the rate check.

The order endpoint uses a unique database key scoped to the authenticated principal. Concurrent identical submissions return the same stored amount; reuse with a different amount returns a conflict. The database row is the simulated effect, not a real payment. External effects require a transactional outbox and downstream idempotency.

\section{Executable laboratory}
Run example 40 against disposable MySQL and Memcached services. The driver starts four workers, checks the served API contract, exercises missing credentials and insufficient scope, submits eight concurrent increments and eight duplicate orders, deletes cache entries and verifies the durable order remains unique. It then revokes the token and checks that subsequent authentication fails. A successful run prints an explicit final success line; consult the release verification record for observed results.

\CodeFile{sql}{examples/catalyst-memcached/schema.sql}{Durable authorization and idempotency schema}
\CodeFile{perl}{examples/catalyst-memcached/lib/MemcachedLab/Controller/Root.pm}{Catalyst API controller and cache coordination}
\CodeFile{json}{examples/catalyst-memcached/openapi.json}{Served OpenAPI contract}
\CodeFile{perl}{examples/40-catalyst-memcached.pl}{Concurrent HTTP regression driver}

\chapter{Scheduled jobs and resumable processing pipelines}
\label{ch:enterprise-pipelines}
\chapterbrief{Build bounded, recoverable scheduled and multi-stage systems.}
\index{Batch processing}
\index{Checkpoints}
\index{Job claims}
\index{Cron}

Consider an import of measurement files or transaction records, followed by
validation, enrichment, aggregation and publication. The input can arrive
twice; a worker can die halfway through a file; a downstream service can
temporarily fail. The design must preserve completeness and avoid duplicate
business effects without assuming that a job executes exactly once.

\strategy{app-pipelines}{Durable job claims}{A work ticket names a task and its current attempt.}{work must survive worker failure.}{a short synchronous operation needs no durable queue.}{DBI, worker processes}

\section{Pipeline decomposition}
\begin{figure}[!htbp]\centering
\begin{tikzpicture}[node distance=5mm]
\node[box,text width=2.3cm] (i) {Input\\manifest};
\node[box,right=of i,text width=2.3cm] (v) {Validate\\and stage};
\node[okbox,right=of v,text width=2.3cm] (q) {Durable\\job queue};
\node[box,right=of q,text width=2.3cm] (o) {Publish\\results};
\draw[arr](i)--(v); \draw[arr](v)--(q);
\node[badbox,below=of v,text width=2.3cm] (e) {Quarantine\\invalid input};
\node[okbox,below=of q,text width=2.3cm] (w) {Bounded\\worker pool};
\draw[badarr](v)--(e); \draw[arr](q)--(w); \draw[arr](w.east)-|(o.south);
\node[lbl,below=of w,text width=5cm] {Claim, process, checkpoint; retry with the same input identity};
\end{tikzpicture}
\caption{A durable processing pipeline separates admission, execution and publication.}
\end{figure}
Assign a stable source identity such as a producer ID plus record ID, or
a file manifest ID plus record offset. Validate the identity and payload
together. Stage data in chunks, quarantine invalid records, and use a unique
constraint on the canonical input identity. If a replay changes the payload
for the same identity, surface a conflict rather than silently overwrite it.

Partition work by a key that owns the invariant: account, device, customer
or dataset. Independent partitions can run concurrently; dependent updates
must be ordered or guarded. Use bounded queues and cap the number of workers
by the downstream connection and write budget. A checkpoint is valid only
when the effects it describes are durable. A high-water mark alone is
unsafe when records finish out of order: track completed ranges or individual
identities until the contiguous prefix is complete.

\section{Prevent overlapping scheduled runs}
Host-local coordination is appropriate when all launchers execute on the
same host and share a stable lock file. The Linux util-linux \code{flock}
command can return a distinct conflict code; preserve it in monitoring.
\CodeFile{bash}{snippets/enterprise-cron.sh}{Host-local cron wrapper (deployment excerpt)}
Do not delete or rotate the lock file while a process may hold it.
Do not let the protected command detach or close the inherited lock
descriptor. A missed run should be observable; do not confuse it with a
successful import. Across hosts use durable claims or a suitable shared
coordinator rather than unrelated local lock files.

\section{Durable claims, attempts and deadlines}
\index{API!SKIP LOCKED}\index{Job claims}
\begin{figure}[!htbp]\centering
\begin{tikzpicture}[node distance=9mm]
\node[box,text width=2.1cm] (r) {Ready};
\node[okbox,right=of r,text width=2.5cm] (w) {Running\\attempt $n$};
\node[okbox,right=of w,text width=2.1cm] (d) {Done};
\node[badbox,below=of w,text width=2.5cm] (x) {Dead letter};
\draw[arr](r)--(w) node[midway,lbl,above]{claim};
\draw[arr](w)--(d) node[midway,lbl,above]{complete};
\draw[badarr](w.south west) to[bend left=20] node[lbl,below left]{expire / retry} (r.south);
\draw[badarr](w)--(x) node[midway,lbl,right]{budget spent};
\end{tikzpicture}
\caption{A claim identifies an attempt, not permanent ownership of a job.}
\end{figure}
A worker claims one row in a short transaction, performs expensive work
outside it, then conditionally completes the matching attempt. A reclaimed
job has a newer attempt number. The earlier worker must not be allowed
to mark the replacement's work complete. Enforce uniqueness or fencing at
the effect destination as well; changing job status alone cannot prevent
a paused worker from executing a side effect.

\CodeFile{sql}{snippets/outbox-claim.sql}{MySQL claim and conditional completion (SQL sequence)}
The placeholders represent bound values from the selected row. Use an
InnoDB job table with a primary key, status, attempt, lease expiry and
indexes appropriate to the claim predicates; a production scheduler also
needs a retry budget, next eligible time and dead-letter handling.
Use \code{SKIP LOCKED} for queue-like access, not for general consistent
reads. Check each conditional update's affected-row count. If no row is
selected, roll back and wait using a bounded polling interval.

\section{Operational examples}
\begin{longtable}{L{.26\linewidth}L{.34\linewidth}L{.32\linewidth}}
\toprule Scenario & Concurrency boundary & Replay contract \\ \midrule\endhead
File import & Manifest and input identity & Reject divergent duplicate payloads \\
Aggregation & Partitioned dataset and period & Recompute or atomically replace a version \\
Report generation & Consistent read snapshot & Publish completed output by version \\
Export & Durable delivery job & Retry using stable destination identity \\
\bottomrule
\end{longtable}
For scheduled jobs on Kubernetes, \code{concurrencyPolicy: Forbid} prevents
overlap only among Jobs belonging to the same CronJob [66]. It is not a
distributed business lock, and scheduling is approximate. Jobs must remain
idempotent even when manually started or recreated.
\CodeFile{yaml}{snippets/enterprise-cronjob.yaml}{Kubernetes scheduling policy (illustrative manifest)}


\chapter{Reliable cross-system integration}
\label{ch:enterprise-outbox}
\chapterbrief{Record durable intent and safely replay delivery.}
\par\smallskip\noindent\textbf{Abbreviations used in this chapter:} DBI --- Database Interface; HTTP --- Hypertext Transfer Protocol; SQL --- Structured Query Language; ID --- identifier.\par\smallskip

\index{Transactional outbox}
\index{Inbox}
\index{Idempotency}
\index{Replay}

A local SQL transaction cannot roll back an HTTP request or a message
already accepted by another system. Holding a local transaction open while
calling the external system does not solve this problem; it also increases
lock duration. Instead, separate durable intent from delivery.

\strategy{app-outbox}{Transactional outbox}{Write the shipping instruction on the same receipt as the sale.}{a committed change must trigger a remote effect.}{the effect can be committed in the same local transaction.}{DBI, unique keys, durable queues}

\section{Transactional outbox and inbox}
\index{Transactional outbox}\index{Idempotency}
\begin{figure}[!htbp]\centering
\begin{tikzpicture}[node distance=10mm]
\node[okbox,text width=5cm] (s) {Source database\\business row + outbox event\\one transaction};
\node[box,right=of s,text width=3cm] (r) {Relay\\claim and send};
\node[okbox,below=of r,text width=3cm] (d) {Destination\\unique event ID};
\node[box,left=of d,text width=5cm] (a) {Source acknowledgement\\conditional on claim};
\draw[arr](s)--(r); \draw[arr](r)--(d) node[midway,lbl,right]{stable ID};
\draw[arr](d)--(a); \draw[arr](a)--(s);
\node[badbox,below=of a,text width=5cm] (f) {Crash after remote commit?\\Retry with the same ID};
\draw[badarr](f.east)-|(d.south);
\end{tikzpicture}
\caption{The outbox preserves intent; destination deduplication resolves ambiguous delivery.}
\end{figure}
Insert an outbox event alongside the business row in the same transaction
[71]. A relay claims committed events, delivers them, and records completion
after acknowledgement. If the response is lost, the relay cannot know
whether the remote operation happened. Reuse the event ID on retries.
The receiving application should commit its deduplication record (inbox)
and effect in a single transaction. A unique inbox row inserted before a
nontransactional external effect does not, by itself, guarantee exactly
one effect.

\section{Failure windows and recovery}
\begin{longtable}{L{.34\linewidth}L{.58\linewidth}}
\toprule Failure point & Recovery \\ \midrule\endhead
Before source commit & Neither business state nor intent is committed; replay input. \\
After source commit & Outbox remains pending; relay later. \\
After destination commit, before acknowledgement & Retry same event ID; destination returns prior result or deduplicates. \\
After acknowledgement & No delivery needed; retain audit and dedupe records for the replay horizon. \\
\bottomrule
\end{longtable}
External APIs vary: some support idempotency keys, others offer lookup by a
client reference, and some offer neither. For the last case, ambiguous
outcomes may require reconciliation or human review. Compensation is a new
business operation, not a rollback of history. Preserve ordering by
aggregate when required, and record failures with correlation identifiers
rather than credentials or full sensitive payloads.

\section{Laboratory: two independent databases}
The following laboratory uses separate local SQLite databases, one for
business state and outbox and one for the receiving effect. It injects a
failure after the destination commits but before the source acknowledges,
then replays the same event. It also verifies source rollback and rejects
reuse of an input identity with a different value. The relay is intentionally
single-worker; multiple relays need the claim protocol of
Chapter~\ref{ch:enterprise-pipelines}.
\CodeFile{perl}{examples/34-replay-outbox.pl}{Example 34: rollback, replay and destination deduplication}


\chapter{Capstone: recoverable measurement processing}\label{ch:capstone}
\chapterbrief{Combine bounded admission, process workers, durable claims, transactional results and replay-safe delivery.}
\index{Capstone}\index{Backpressure}\index{Batch processing}\index{Transactional outbox}
This synthetic application receives a stream of measurement records. It
validates each source identity, computes a deterministic transformation,
commits results with an outbox record, and delivers each effect into a
separate database. SQLite makes the laboratory portable; it is not a
substitute for the MySQL row-locking examples.

\section{Architecture and invariants}
\begin{figure}[!htbp]\centering
\begin{tikzpicture}[node distance=7mm]
\node[box,text width=10cm] (a) {Input stream: stable record ID and value};
\node[okbox,below=of a,text width=10cm] (b) {Bounded admission: validate identity and stage jobs};
\node[box,below left=8mm and 0mm of b.south,anchor=north east,text width=4.3cm] (c) {Worker A\\claim and compute};
\node[box,below right=8mm and 0mm of b.south,anchor=north west,text width=4.3cm] (d) {Worker B\\claim and compute};
\node[okbox,below=8mm of $(c.south)!0.5!(d.south)$,text width=10cm] (e) {One transaction: result, completion and outbox event};
\node[box,below left=8mm and 0mm of e.south,anchor=north east,text width=4.3cm] (f) {Aggregate\\committed results};
\node[okbox,below right=8mm and 0mm of e.south,anchor=north west,text width=4.3cm] (g) {Relay to independent\\deduplicating destination};
\draw[arr](a)--(b); \draw[arr](b)--(c); \draw[arr](b)--(d);
\draw[arr](c)--(e); \draw[arr](d)--(e); \draw[arr](e)--(f); \draw[arr](e)--(g);
\end{tikzpicture}
\caption{Admission bounds unprocessed records; effects and progress commit together. Delivery is independently recoverable.}
\end{figure}
\strategy{capstone}{Replay-safe processing pipeline}{A warehouse records incoming packages and work tickets, then records each completed package and its shipping instruction together.}{input can repeat, workers can fail and delivery crosses a transaction boundary.}{durable staging costs more than recovery is worth.}{DBI, DBD::SQLite, forked workers, durable attempt IDs and unique constraints.}
The application enforces four invariants: identical input identities
cannot represent different values; a stale claim cannot finish a new
attempt; a completed result and its delivery intent are atomic; the
destination stores at most one effect per stable event ID. The source
can remain pending after a successful remote commit, so acknowledgement
is not equivalent to delivery.

\section{Admission, processing and aggregation}
Ingest a bounded batch at a time, then drain it with a fixed worker pool.
The input is generated incrementally, rather than loaded as an unbounded
array. A production file adapter should use a stable manifest and record
offset. Publish completed files with rename within the same filesystem,
with file and directory synchronization if crash durability is required
[79]. Atomic rename alone is not a durability guarantee.
Use a short database transaction to claim work, compute outside it,
then conditionally commit the matching attempt. SQLite immediate
transactions serialize writers [77,78]. On InnoDB replace this claim
with indexed locking reads and \code{SKIP LOCKED}; retry engine-specific
deadlocks. Never hold a transaction across a remote call.

\section{Failure injection and recovery}
\begin{longtable}{L{.33\linewidth}L{.59\linewidth}}
\toprule Injected condition & Expected result \\ \midrule\endhead
Worker exits after claim & Expired attempt can be reclaimed.\\
Old worker completes late & Conditional completion rejects it.\\
Duplicate input & Existing payload retained; mismatch rejected.\\
Destination commits before source ack & Retry stable event ID; no extra effect.\\
Re-run completed input & Existing committed results remain unchanged.\\
\bottomrule
\end{longtable}
The laboratory advances the recorded expiry explicitly to avoid
sleep-based races. Real deployments require a consistent time source,
bounded lease renewal, explicit retry/dead-letter policy and supervisor
restart behavior. This fixture uses finite generated work and no external
network. It demonstrates the protocol rather than a production scheduler
or exactly-once transport. Its admission bound applies to processing
jobs, not the entire outbox backlog; a deployed system should also bound
or throttle delivery backlog independently.

\section{Run and inspect the laboratory}
\CodeFile{bash}{snippets/capstone-run.sh}{Run the self-verifying capstone and its independent tests}
The complete implementation appears in Appendix~\ref{app:capstone}.
For production, partition by the resource whose updates must be ordered,
bound queue depth and worker connection budgets, and expose queue age,
attempt counts and outbox lag. Report throughput only with the dataset,
hardware, worker count and measured contention documented.

\partstart{VI / Production engineering}{Engineering for production}{Evolve legacy systems, verify invariants and operate bounded, observable services.}
\chapter{Modernizing legacy locking and database access}
\label{ch:enterprise-migration}
\chapterbrief{Migrate locks while old and new writers coexist.}
\par\smallskip\noindent\textbf{Abbreviations used in this chapter:} DBI --- Database Interface; DDL --- data definition language.\par\smallskip

\index{Migration}
\index{MyISAM}
\index{Metadata locks}

Legacy applications often contain multiple generations of synchronization:
file locks around scripts, table locks around bulk writes, named advisory
locks around cross-table operations, and transactional row locks in newer
code. Removing a legacy lock without identifying every writer can create
a concurrency bug that did not previously exist.

\strategy{app-migration}{Incremental invariant migration}{Install the new guard before taking away the old one.}{multiple generations of writers share state.}{the contracts cannot coexist safely.}{DBI, InnoDB, regression tests}

\section{Inventory before migration}
\begin{figure}[!htbp]\centering
\begin{tikzpicture}[node distance=5mm]
\node[box,text width=10cm] (a) {Inventory writers, storage engines, locks and transaction owners};
\node[box,below=of a,text width=10cm] (b) {Measure contention and capture regression fixtures};
\node[okbox,below=of b,text width=10cm] (c) {Add guards, constraints and replay-safe identities};
\node[okbox,below=of c,text width=10cm] (d) {Canary new writers; compare and reconcile outcomes};
\node[box,below=of d,text width=10cm] (e) {Retire old locking only when every writer obeys the new contract};
\draw[arr](a)--(b); \draw[arr](b)--(c); \draw[arr](c)--(d); \draw[arr](d)--(e);
\end{tikzpicture}
\caption{Incremental migration preserves the invariant across mixed generations of code.}
\end{figure}
Inspect table engines, foreign keys, unique constraints, isolation levels,
queries, cron launchers and administrative scripts. MyISAM is nontransactional;
changing a DBI transaction boundary cannot make a MyISAM update roll back.
A move to InnoDB changes locking, index and operational requirements [68].
Test migration and rollback procedures on representative disposable data.

\section{Table locks and connection-local state}
MySQL \code{LOCK TABLES}, DDL and some administrative statements can
interact with transaction boundaries through implicit commits. Consult the
server's version-specific documentation [11,67]. Avoid combining
table-locking legacy code with ORM-managed transactions without a precise
contract. Named advisory locks, temporary tables, session variables and
isolation overrides are connection-local. A pooled connection must be reset
before reuse, or explicitly dedicated to the operation.

Long-running reports can retain snapshots and delay purge or contend with
writers, depending on the query and isolation. Prefer appropriate indexes,
bounded report windows and consistent snapshots. A replica may be useful
for reports that tolerate lag, but not for correctness-critical read-after-write
decisions unless the required consistency is enforced.

\warn{An online schema change can still require metadata locks. Test the
DDL algorithm and lock requirements for the exact server and operation;
do not assume that an online label means zero blocking.}

\section{Migration acceptance criteria}
Before removing a lock, list every code path that can mutate the protected
resource. Establish an invariant test, deploy the new guard while preserving
compatibility, and compare results with representative load. Record a
rollback strategy that does not reintroduce mixed incompatible writers.
Avoid a big-bang migration from local to distributed locking when a database
constraint can solve the problem more directly.


\chapter{Concurrency verification in application test suites}

\label{ch:enterprise-testing}
\chapterbrief{Force interleavings, isolate fixtures and verify recovery.}
\index{Concurrency tests}
\index{Test2}
\index{Fault injection}
\index{API!Test2::Harness}

Concurrency tests should prove an invariant, not merely demonstrate that
two workers finished. Random sleeps can help stress a system but are not
a reliable method of forcing a specific race. Use rendezvous points,
independent connections, explicit deadlines and failure injection.

\strategy{app-testing}{Deterministic race tests}{Stop workers at a known point before releasing them. A rendezvous makes the disputed interleaving repeatable instead of relying on lucky scheduling.}{a regression depends on a particular ordering of operations.}{measuring throughput or latency; use repeated workload trials instead.}{Thread::Queue, Test2 and process fixtures; assert the invariant and check every worker outcome.}

\section{Force the interleaving}
\begin{figure}[!htbp]\centering
\begin{tikzpicture}[x=1.35cm,y=.8cm]
\node[anchor=east] at (0,2){Worker A}; \node[anchor=east] at (0,0){Worker B};
\draw[lane](0,2)--(7.5,2); \draw[lane](0,0)--(7.5,0);
\foreach \y in {0,2}{
  \node[box] at (1,\y){Read 0};
  \node[box] at (3,\y){Ready};
  \node[badbox] at (5.8,\y){Write 1};
  \draw[arr](1.65,\y)--(2.3,\y);
  \draw[arr](3.65,\y)--(5.0,\y);
}
\draw[teal,thick](4.4,-.45)--(4.4,2.45);
\node[lbl,above] at (4.4,2.45){Release both workers};
\node[badbox] at (7.3,1){Final 1\\expected 2};
\draw[badarr](6.45,2)--(6.9,1.45);
\draw[badarr](6.45,0)--(6.9,.55);
\end{tikzpicture}
\caption{A deterministic test makes both readers observe the same old value before either writes.}
\end{figure}
The laboratory first protects individual reads and writes with locks,
which is insufficient for a compound update. A pair of queues forces
the lost-update interleaving. It then locks the complete read-modify-write
and verifies the exact final count.
\CodeFile{perl}{examples/35-deterministic-race.pl}{Example 35: reproducible lost update and corrected critical section}

\section{Parallel test isolation}
Test2::Harness can run multiple test files concurrently [69].
Parallelizing a test runner is separate from testing concurrency inside a
single test. Give parallel tests separate schemas or uniquely scoped data,
independent temporary directories and distinct service namespaces. Do not
flush a shared cache or reset a shared schema while another test is running.
Serialize tests that cannot yet be isolated, but treat that as a migration
constraint rather than a correctness guarantee.

\begin{longtable}{L{.31\linewidth}L{.61\linewidth}}
\toprule Test scenario & Invariant to assert \\ \midrule\endhead
Opposing transactions & Conserved totals; no negative balance; bounded retries \\
Worker crash after claim & Job becomes recoverable; old attempt cannot complete it \\
Lost remote response & Same event ID results in one durable effect \\
Duplicate input & Original payload preserved; divergent duplicate rejected \\
Process fork & Child uses a new connection; parent remains operational \\
Rolling restart & Admission stops; committed work remains recoverable \\
\bottomrule
\end{longtable}
Use deadlines at both the test and job level so a failed barrier cannot
hang CI. Run stress tests repeatedly with logged random seeds in addition
to deterministic tests. A test against SQLite is useful for application
composition but cannot prove MySQL-specific locking or deadlock behavior.
Run actual engine integration tests for those claims.

\section{Testing strategy}
Test correctness with invariants, not timing assumptions. Start competing workers together, vary worker counts and scheduling, inject random delays, and assert final totals, concurrency bounds and the absence of duplicates. Run under a watchdog so a deadlock fails the build instead of stalling it. Add fault injection: throw while holding a lock, kill a process holding a permit, interrupt a database session, provoke a deadlock with reversed row order, pause a lease holder past its TTL (\code{kill -STOP}), restart Redis without persistence, and replay duplicate deliveries. Benchmark throughput, p95/p99 wait, memory per worker and tail latency under realistic contention, and compare a simple mutex against alternatives before adopting a complex primitive.
\CodeFile{bash}{snippets/guide-10.sh}{Laboratory regression suite}
\tip{Make races reproducible before fixing them. The laboratory widens race windows with a delay parameter (\code{delay=0.3}) and fires requests from forked processes; a race that fails every run is one you can prove fixed.}


\section{A Test2 integration fixture}
The Test2 fixture in Appendix~\ref{app:capstone} runs separate child
processes and checks exit status and declared invariants. Temporary
databases are created by the capstone, so multiple instances do not
share state. Run it with \code{prove} or the optional
\code{yath test -j2 examples/t/} runner [80]. The integration suite runs
service-backed examples serially where namespaces are shared.

\chapter{Production deployment and observability}
\label{ch:enterprise-operations}
\chapterbrief{Budget capacity, drain workers and diagnose contention.}
\index{Observability}
\index{Connection budgets}
\index{Graceful shutdown}
\index{Metrics}

Concurrency controls must remain valid as worker counts, replicas,
schedulers and service availability change. A configuration that works on
a single development process may fail when a production deployment forks
workers or rolls out multiple replicas.

\strategy{app-operations}{Bounded admission}{A waiting room limits work entering a busy service.}{downstream resources have finite capacity.}{limits are guessed without measuring load.}{worker pools, metrics, deadlines}

\section{Budget workers and connections}
For a hypothetical deployment of four replicas, six web workers per
replica and one live database connection per worker, web traffic alone can
consume 24 database connections. Add batch workers, administrative sessions,
pool overflow and rollout surge to the budget. A connection limit is not
a throughput target. Start with bounded concurrency and measure queueing
and database saturation before increasing parallelism.

Use the same resource ownership discipline with Apache, FastCGI and PSGI
workers: inspect the actual server lifecycle, avoid inheriting live clients,
and distinguish process-local from cluster-wide locks. Maintain bounded
request and job deadlines. Never hold an ordinary synchronization lock
across an unbounded external call.

\section{Graceful termination and durable recovery}
\begin{figure}[!htbp]\centering
\begin{tikzpicture}[node distance=7mm]
\node[box,text width=2.7cm] (t) {Termination requested};
\node[okbox,right=of t,text width=2.7cm] (s) {Stop accepting work};
\node[okbox,right=of s,text width=2.7cm] (d) {Drain bounded in-flight work};
\node[box,below=of d,text width=2.7cm] (e) {Commit, close, exit};
\node[badbox,below=of s,text width=2.7cm] (r) {Deadline expires\\reclaim later};
\draw[arr](t)--(s); \draw[arr](s)--(d); \draw[arr](d)--(e);
\draw[badarr](d.south west)--(r.north east);
\end{tikzpicture}
\caption{Graceful shutdown is an optimization; durable recovery also handles forced termination.}
\end{figure}
On a shutdown request, mark the process unready, stop claiming jobs,
finish bounded in-flight work, then close resources. A Kubernetes
termination grace period is finite [72]; forced termination must not
require a final cleanup handler for correctness. Expired attempts,
idempotent effects and durable checkpoints provide recovery when a
worker cannot finish.

\section{Measure symptoms, not just CPU}
\begin{longtable}{L{.33\linewidth}L{.59\linewidth}}
\toprule Signal & Meaning and action \\ \midrule\endhead
Lock wait duration & Histogram; increasing tail latency identifies contention. \\
Lock hold duration & Histogram; investigate slow critical sections or transactions. \\
Job queue age & Oldest eligible job; detects starvation or insufficient capacity. \\
Attempt and retry counters & Separate conflicts, transient failures and exhausted budgets. \\
Database connections & Active, idle and waiting; check saturation and leaked sessions. \\
Outbox age and backlog & Delivery delay or remote outage; alert on sustained lag. \\
\bottomrule
\end{longtable}
Trace the business operation, transaction, claim and delivery separately.
Record request ID, job ID, attempt, worker PID and elapsed time in structured
logs. Keep high-cardinality identifiers out of metric labels; use bounded
operation and outcome categories instead [73]. Never log tokens, passwords
or private payloads. Correlate traces and logs with database lock diagnostics.

\tip{Start with a small dashboard: queue age, throughput, retry rate,
lock wait percentiles, connection usage and failed jobs. Use these signals
to choose the next concurrency improvement rather than tuning blindly.}

\section{Architecture review checklist}
\begin{itemize}
\item Identify every writer and the resource that enforces its invariant.
\item Document transaction scope, lock order and retry classification.
\item Bound worker counts, queues, connection use and operation deadlines.
\item Define durable identities, replay rules and ambiguous-outcome recovery.
\item Verify crash, timeout, retry and duplicate delivery scenarios.
\item Establish deployment, rollback and observable failure criteria.
\end{itemize}
\section{Metrics and debugging}
Instrument acquisition wait and hold duration separately, using monotonic time. Count acquisitions, contention, timeouts, renewals, lost leases, rejected fences, retries and errors. Keep label cardinality bounded (lock class, backend, operation), never customer IDs or individual resource names. Structured logs may carry a correlation ID, lock name and fencing token. Trace waiting separately from useful work so that tail latency can be attributed.
Example 26 provides \code{with_timed_lock}, a wrapper that records wait and hold time for any critical section. Each thread keeps its own samples, so measurement adds no shared contention, and percentiles are computed after \code{join}.
\CodeFile{perl}{examples/26-lock-metrics.pl}{\nolinkurl{26-lock-metrics.pl}}
\CodeFile{text}{snippets/guide-08.txt}{Output on the verification host (4 threads, 1 ms of work under the lock)}
Read this the way you would read production data: waiting (3~ms at the median) dominates useful hold time (1~ms) because four threads queue for one lock, so the fix is shorter or finer-grained critical sections, not faster work. The maximum hold time is nearly three times the p99, a typical sign of a scheduler preemption while the lock was held.
\CodeFile{bash}{snippets/guide-09.sh}{Linux troubleshooting tools}
\tip{Export lock \emph{wait} time and lock \emph{hold} time as separate histograms (example 26). Rising wait with flat hold means contention; rising hold means the critical section itself got slower.}

\section{Deployment checklist}
\begin{itemize}
\item Identify the synchronization scope and state the invariant; decide whether the lock is for efficiency or correctness.
\item Document lock identity, acquisition order, maximum wait, cancellation and cleanup behavior.
\item Ensure all participants cooperate; separate independent resources rather than locking globally.
\item Never share DBI handles or Redis connections across workers; validate storage engine, isolation and query plans.
\item Set per-session lock wait timeouts consistent with retry budgets; disable transparent reconnects on coordination connections.
\item For leases: random owner tokens, owner-checked renewal and release, bounded heartbeats, and destination-enforced fencing for correctness-critical effects.
\item Give coordination keys a dedicated Redis/Valkey instance with \code{noeviction}, a restricted ACL and TLS on untrusted networks.
\item Design for process death, pauses, connection loss, retry safety and shutdown before deploying.
\item Record dependency versions and test each supported Perl, OS and server combination, including legacy hosts.
\item Observe contention, lost leases, rejected fences and queue age; optimize from measurements.
\end{itemize}

% ================================================================ glossary



\section{Repeatable contention measurement}
\index{Benchmarking}
Example 39 compares one lock against independent striped counters using
three trials of each at four workers. It validates the final count in
every trial and reports elapsed wall time, including thread creation
and join. This is not a pure lock-latency benchmark or proof of linear
scaling. Record interpreter version, CPU, workload, worker count and
scheduler conditions. Warm workers before timing for more stable
microbenchmarks; vary the contention distribution and report variation
and tail latency rather than only a mean. The full benchmark is in
Appendix~\ref{app:capstone}.
\section{Connection pool exhaustion}
\index{Connection pools}
A pool is a semaphore around a finite set of live resources. When all
slots are busy, bound the number and duration of waiters rather than
allocating more connections. Never hold a checked-out connection while
waiting for a second slot unless the pool contract explicitly reserves
enough capacity. Separate admission timeout, query timeout and total
operation deadline. Export active, idle and waiting counts, and record
which operations retain a handle during external I/O.

\clearpage\appendix
\renewcommand{\currentpart}{Technical appendices}
\chapter{Native synchronization implementation}\label{app:native}

\par\smallskip\noindent\textbf{Abbreviations used in this chapter:} CPU --- central processing unit.\par\smallskip
\section{Futexes, atomics and memory ordering}
\strategy{futex}{Atomic operations and futexes}{Underneath every lock is a single memory word that the CPU can test and change in one indivisible step. If the word says ``free'', a thread flips it to ``taken'' without asking the operating system. Only when it must wait does it ask the kernel to put it to sleep.}
{you need to understand lock performance, read \code{strace} output, or write native (XS or C) code.}
{writing Perl: never build locks from raw atomics or futex calls on Perl scalars; use \code{lock} or a vetted module.}
{indirectly, through \code{lock}, which rests on pthread mutexes and futexes; example 21 is the C implementation.}
Linux futexes operate on an aligned 32-bit word in memory. A native lock attempts an atomic state transition in user space and enters the kernel only when contention requires sleeping. \code{FUTEX_WAIT} sleeps only if the word still holds an expected value, checked atomically by the kernel, which closes the lost-wakeup window between checking and sleeping. \code{FUTEX_WAKE} wakes waiters; it does not transfer ownership, so a woken thread must compete again [6].

Example 21 is a complete, working futex mutex: the classic three-state design used, with refinements, inside glibc's \code{pthread_mutex_t}, and therefore underneath Perl's \code{lock}. State 0 is unlocked, 1 is locked with no waiters, and 2 is locked with possible waiters. An uncontended lock and unlock is one atomic instruction each and makes no system call. A contended locker sets the state to 2 before sleeping, so the unlocker knows it must issue a wake. The acquire ordering on lock and the release ordering on unlock are what make the protected writes visible to the next owner: atomicity alone would not.
\fig{futexstates}{States of the three-state futex mutex in example 21. Only red transitions enter the kernel.}
\CodeFile{c}{examples/21-futex-mutex.c}{\texttt{21-futex-mutex.c} (C11; build with \texttt{cc -O2 -pthread})}
\tested{Eight threads performed 8{,}000{,}000 increments of a deliberately non-atomic counter with no lost update, in repeated runs and under ThreadSanitizer. Between 3\% and 6\% of acquisitions took the kernel slow path.}
The same behavior is visible from Perl without writing any C. Tracing the 40{,}000 \code{lock} calls of example 01 shows how rarely the kernel is involved:
\CodeFile{bash}{snippets/guide-02.sh}{Observing futex use behind Perl's lock()}
Fifty-four system calls for forty thousand acquisitions: the rest completed in user space. A rising futex count under load is therefore a direct measure of lock contention.

Do not use raw futex calls against Perl scalars: Perl SV storage is not a stable, aligned atomic integer, and interpreter cloning complicates addresses. Use the core facilities or a vetted native library. Relaxed atomics provide atomicity without ordering; only acquire/release (or stronger) ordering publishes protected data. Spinlocks waste CPU when the owner is descheduled and are a poor default for Perl-level critical sections. Robust native mutexes can report owner death (\code{EOWNERDEAD}), but recovery still requires repairing the protected state.

% ================================================================ chapter 3

\chapter{API security and concurrent state transitions}\label{app:api}

\par\smallskip\noindent\textbf{Abbreviations used in this chapter:} API --- application programming interface.\par\smallskip
\par\smallskip\noindent\textbf{Abbreviations used in this section:} ACL --- Access control list; API --- Application programming interface; CSPRNG --- Cryptographically secure pseudorandom number generator; HMAC --- Hash-based message authentication code; ID --- Identifier; JSON --- JavaScript Object Notation; JWT --- JSON Web Token; SHA-256 --- Secure Hash Algorithm with a 256-bit digest; TLS --- Transport Layer Security; TTL --- Time to live; YAML --- YAML Ain't Markup Language (recursive name).\par\smallskip
% API credentials on Redis: included from catalyst-sessions.tex
\section{API credentials on Redis}\label{sec:api-tokens}
Browsers authenticate with a session cookie. API clients, such as mobile apps, single-page applications, partner integrations and command-line tools, send a credential in a header with every request. The OpenAPI document describes those credentials as \emph{security schemes}; Redis is where the server checks, rotates, limits and revokes them. Every concurrency problem in this handbook turns up here too: two tabs refreshing the same token are a race, a retried signed request is a replay, and a per-key quota is a shared counter.

Table~\ref{tab:credentials} compares the five credentials the laboratory implements. They complement one another: a browser logs in with a session, creates an API key for a script, or obtains a short-lived access token and a refresh token for a mobile client.

\begin{table}[!htbp]\centering\small
\caption{Credentials, their lifetimes and where Redis takes part.}\label{tab:credentials}
\begin{tabular}{L{.14\linewidth}L{.14\linewidth}L{.12\linewidth}L{.27\linewidth}L{.17\linewidth}}\toprule
Credential & Held by & Lifetime & Redis keys and checks & OpenAPI scheme \\ \midrule
Session cookie & browser & idle timeout & \code{sess:<id>}, per-session lock & \code{apiKey}, \code{in: cookie} \\
API key & scripts, integrations & months; rotated & \code{apikey:<sha256>}; per-key rate counter & \code{apiKey}, \code{in: header} \\
Access token (JWT) & API client & 5 minutes & none on the happy path; \code{jwt:deny:<jti>}, family check & \code{http}, \code{bearer}, \code{JWT} \\
Refresh token & API client, stored securely & 30 days; rotated on every use & \code{rt:{fam}:<sha256>}, \code{rtfam:{fam}} & issued by \code{/token}; \code{oauth2} if an authorization server issues it \\
Signed request & partner systems & one request & \code{nonce:<kid>:<nonce>} (\code{SET NX}) & \code{apiKey} header, described \\
\bottomrule\end{tabular}
\end{table}

\subsection*{Describing the schemes in OpenAPI 3.2}
The laboratory's \code{openapi.yaml} declares four schemes and states the requirement on every operation. An empty requirement list (\code{security: []}) marks an operation as public, which is clearer than relying on a global default. Alternatives in a requirement list are ORed; schemes inside one requirement object are ANDed.
\CodeFile{yaml}{snippets/api-credentials-01.yaml}{Security schemes and one operation from \texttt{openapi.yaml}}
\version{OpenAPI 3.2.0 (September 2025) [51] added three security features: the OAuth~2.0 device authorization flow (\code{deviceAuthorization}), \code{oauth2MetadataUrl} for discovering an authorization server's metadata, and a \code{deprecated} flag that tells clients to move off an old scheme, such as a legacy API-key header during a migration to tokens. An OpenAPI document describes credentials; it does not enforce them. \code{OpenAPI::Modern} [53] can validate requests and responses against the document in tests and middleware. The laboratory validates its document with \code{openapi-spec-validator} and checks that every exercised route is documented with an explicit requirement; it does not use \code{OpenAPI::Modern} at run time.}
\tip{Put the OpenAPI document under test. Example 30 loads it with \code{YAML::XS} and fails if a route it calls is missing or has no \code{security} entry. A forgotten requirement then shows up as a failing test, not as an unprotected endpoint.}

% ---------------------------------------------------------------- API keys
\subsection*{API keys: store the hash, never the key}
\index{API!sha256_hex@\texttt{sha256\_hex}}
\strategy{apikey}{Hashed API keys}{A hotel key card. The front desk does not keep a copy of your card, only a record of which card number opens which room. If someone photographs the desk's records, they still cannot open the door.}
{scripts, integrations and service accounts need a long-lived credential that a person can revoke.}
{the client is a browser (use the session) or can perform an OAuth flow (use short-lived tokens).}
{\code{Crypt::URandom} or \code{/dev/urandom} for 32 bytes, \code{Digest::SHA::sha256_hex} for storage, \code{HSET apikey:<hash>} with \code{EXPIRE}; example 30.}
The server generates the key from a CSPRNG, shows it once, and stores only its SHA-256 (Figure~\ref{fig:apikeyflow}). On every request it hashes the presented key and looks the hash up. Because the lookup is by hash, there is no stored secret to compare and no timing side channel to defend.
\fig{apikeyflow}{API key issue and use. The plaintext key exists only in the response that creates it and in the client's configuration.}
\begin{itemize}
\item \textbf{No salt, no slow hash.} Argon2 and bcrypt exist to protect low-entropy secrets such as passwords from guessing. A 256-bit random key cannot be guessed, so a single SHA-256 is correct, and it keeps authentication to one cheap lookup.
\item \textbf{Recognizable prefix.} \code{slk_} followed by 64 hex digits lets log filters and secret scanners spot a leaked key, and tells support staff which system a key belongs to.
\item \textbf{Shown once, never cached.} The creating response carries \code{Cache-Control: no-store}. Lost keys are replaced, never recovered.
\item \textbf{Revocation and rotation.} \code{user-apikeys:<user>} lists a user's key hashes, so ``revoke all'' is one loop. Allow two live keys per integration so the old one can be retired after the new one is deployed. Expire unused keys.
\item \textbf{Limits per key.} The fixed-window script from the login limiter (Strategy~\ref{str:ratelimit}) runs per key hash. Responses advertise the quota with \code{RateLimit-Policy} and \code{RateLimit} headers in the syntax of the current IETF draft [54], \code{"perkey";q=10;w=60}, plus \code{Retry-After} on 429.
\end{itemize}

% ---------------------------------------------------------------- JWT
\subsection*{Access tokens: short-lived JWTs, verified strictly}
\index{API!Crypt::JWT}
\strategy{jwtrevocation}{Short-lived token with a revocation list}{A day pass at a theme park. Staff at each ride only check the stamp and the date, which is fast. If a pass is reported stolen, its number goes on a short list at the gates until the day is over; after that the pass is useless anyway.}
{many requests must be authorized cheaply, ideally without a lookup, and you still need logout to take effect at once.}
{tokens must live for hours or days (revocation lists grow and stolen tokens stay useful) or every request needs fresh server-side state anyway.}
{\code{Crypt::JWT} \code{encode_jwt}/\code{decode_jwt} with \code{accepted_alg}, \code{verify_exp => 1}, \code{verify_iss}, \code{verify_aud}; \code{SET jwt:deny:<jti> 1 EX <remaining>}; example 30.}
A JWT access token is a signed claim set that a server can verify without asking anyone. That is its advantage and its danger: a valid signature says nothing about whether the token was meant for \emph{this} API, as \emph{this} kind of token, or whether it is still current. Verify all of it.
\CodeFile{perl}{snippets/api-credentials-02.pl}{Access-token verification in \texttt{SessionLab::Tokens}}
\warn{In \code{Crypt::JWT} [52] the default for \code{verify_exp} checks \code{exp} only \emph{if the claim is present}. A token minted without \code{exp}, by a bug in another service or by anyone who holds the key, is accepted forever. Pass \code{verify_exp => 1}. The library does not check the \code{typ} header either, so the explicit access-token type of RFC~9068 (\code{at+jwt}) [50] is checked by hand. Without that check, any other JWT signed with the same key, such as an ID token, would be accepted as an access token.}
Revocation costs two Redis reads: \code{EXISTS jwt:deny:<jti>}, set at logout, and \code{EXISTS rtfam:{fam}}, which disappears when the token's refresh family is revoked. The denylist entry expires when the token would have expired anyway, so the list never grows beyond five minutes of logouts. Drop the family check if you can tolerate up to five minutes of access after revocation; keep it if not.
\note{The laboratory signs with HS256, a shared secret, because one application both issues and verifies. When other services verify your tokens, switch to an asymmetric algorithm (ES256, EdDSA or RS256) with a \code{kid} header and a published JWKS, so verifiers never hold a key that can mint tokens. RFC~8725, the JWT best current practice [49], covers algorithm pinning and explicit typing; an update to it is in progress at the IETF.}

% ---------------------------------------------------------------- refresh
\subsection*{Refresh-token rotation and the concurrency trap}
\index{Refresh-token rotation}
\strategy{refreshrotation}{Refresh-token rotation with reuse detection}{A cinema ticket that is torn when you go in. If someone shows up later with the other half of a torn ticket, the cinema knows a copy exists and cancels every ticket bought on that booking, including yours, until you prove who you are again.}
{long-lived sessions for API clients that cannot keep a secret perfectly: mobile apps, single-page applications, desktop tools.}
{the client can use sender-constrained tokens (mutual TLS or DPoP, RFC~9449 [57]), which make a stolen token useless by itself; or the client is a confidential server with a secure store.}
{one Lua script per refresh; successor = \code{HMAC(secret, old token)}; family hash tag \code{{fam}}; grace window for duplicates; example 30.}
RFC~9700, the OAuth~2.0 security best current practice published in January 2025 [48], requires refresh tokens for public clients to be either sender-constrained or rotated. Rotation means every refresh returns a new refresh token and invalidates the old one. If the old one is ever presented again, someone has a copy: the server revokes the entire chain, called the \emph{token family}, and the user signs in again.

The concurrency trap is that legitimate clients present the same refresh token more than once. Two browser tabs notice an expired access token at the same moment; a mobile app retries after a timeout whose response was lost. Figure~\ref{fig:refreshrace} shows what a read-then-write implementation does with three simultaneous refreshes: every request reads ``active'' before any of them writes, and each mints its own successor. The family now has three live branches, and reuse detection can no longer tell a thief's copy from a second tab. A strict implementation that revokes on the second presentation has the opposite bug: it logs out every user who has two tabs open.
\fig{refreshrace}{Six simultaneous refreshes in the laboratory: the naive version left three live refresh tokens in one family; the atomic version left exactly one, and every caller received it.}
The fix has three parts:
\begin{enumerate}
\item \textbf{One atomic decision.} A Lua script checks the family, reads the token's state, and either rotates it, replays the recorded successor, or revokes the family, all in one step. Redis runs scripts one at a time, so duplicate presentations are serialized exactly as the single-command patterns in Chapter~\ref{ch:redis} are.
\item \textbf{A deterministic successor.} The successor is \code{HMAC(secret, presented token)}, so every duplicate can be handed the same token without the server ever storing a plaintext refresh token. Only the hashes live in Redis.
\item \textbf{A short grace window.} A duplicate presentation within the window receives the same successor (\emph{replay}); after the window, it is treated as theft (\emph{reuse}). The script reads Redis \code{TIME} so that application clocks never matter.
\end{enumerate}
The Perl controller computes the successor and hashes, then passes three keys and three arguments to its Redis wrapper. This application-side excerpt is from the complete controller in Section~\ref{sec:catalyst-lab}.
\CodeFile{perl}{snippets/api-credentials-04.pl}{Perl controller excerpt -- submit the atomic operation}
The next listing runs inside Redis, not in the Catalyst worker. The numbered keys and arguments correspond directly to the invocation above. Its return value selects the subsequent Perl response path.
\CodeFile{lua}{snippets/api-credentials-03.lua}{Redis operation -- refresh-token rotation}
\fig{refreshstates}{The life of one refresh token. Every token in a family shares the hash tag, so the script's keys live in one Redis Cluster slot.}
The refresh token itself carries its family ID (\code{<fam>.<64 hex digits>}), so the server can build the \code{{fam}}-tagged keys before any lookup. Deleting \code{rtfam:{fam}} is the single switch that revokes everything: the next refresh returns \code{revoked}, and the family check makes outstanding access tokens fail at once.
\warn{The grace window is a deliberate trade-off. Within it, an attacker who replays a stolen token receives the same successor as the real client, so detection waits until one of them rotates again. Keep the window short; identity providers expose it as a setting (Okta's default is 30 seconds, configurable from 0 to 60; Amazon Cognito allows 0 to 60 seconds; Auth0 calls it the reuse interval) [58--60]. Never extend it to hide client bugs.}
\tip{Make clients refresh single-flight. In an async Perl client, guard the refresh call with a \code{Future::Mutex} (Section~\ref{sec:async}) and have every waiting request reuse the result; in a browser, share one in-flight promise across tabs with a \code{BroadcastChannel} or a Web Lock. The server's grace window then becomes a safety net, not the main mechanism.}

% ---------------------------------------------------------------- logout
\tested{Example 30 exercised 33 checks under four Starman workers, three times. Only API-key hashes reached Redis; altered keys and invalid token claims were rejected; rate limits returned 429. Atomic refresh produced one successor for six concurrent requests, whereas the naive path left three live tokens.}

\subsection*{Logout and revocation}
Logout must stop three things: the current access token, every future refresh, and any access token already handed to another tab. The laboratory does all three with two commands: \code{SET jwt:deny:<jti> 1 EX <seconds until exp>} and \code{DEL rtfam:{fam}}. A password change or an administrator's ``sign out everywhere'' deletes every family listed for the user, just as \code{logout_everywhere} deletes every session (Section~\ref{sec:catalyst-lab}). Because both keys expire on their own, revocation leaves no residue to clean up.

% ---------------------------------------------------------------- signed
\subsection*{Signed requests and replay protection}
\index{Nonce}\index{Replay protection}
\strategy{signedrequest}{Signed request with a single-use nonce}{A cheque with a serial number. The bank checks the signature first, then the date, and finally that it has never cashed that serial number before. A photocopy of a cashed cheque is worthless.}
{a partner calls you over networks or proxies you do not control, the body must not be altered, and a captured request must not be replayable.}
{TLS between two systems you operate already provides integrity; mutual TLS or RFC~9421 HTTP Message Signatures are available and agreed with the partner.}
{\code{Digest::SHA::hmac_sha256_hex} over a canonical string; a constant-time comparison; \code{SET nonce:<kid>:<nonce> 1 NX EX 600}; example 30.}
The client signs the method, path, timestamp, nonce and raw body with a per-key signing secret. The server checks them in a strict order (Figure~\ref{fig:signedverify}): first everything that needs no writes, then the signature, and only then records the nonce with \code{SET NX}. A replay inside the timestamp window finds the nonce and gets 409; a replay outside it fails the timestamp check. The nonce TTL is twice the window so that a nonce cannot expire while its timestamp is still acceptable.
\fig{signedverify}{Verification order for a signed request. The first write happens only after the signature has been verified.}
\warn{Record the nonce \emph{after} verifying the signature, never before. Otherwise anyone who can see or guess a nonce can send a forged request carrying it, and the legitimate request that follows is rejected as a replay. Compute the signature over the raw request bytes, before any JSON parsing or normalization, and compare signatures in constant time.}
\note{For new partner APIs prefer a standard: RFC~9421 HTTP Message Signatures [56] defines the canonical string, the covered components and the header syntax. The ordering and nonce rules above apply unchanged.}

% ---------------------------------------------------------------- idempotency
\tested{Reuse outside the 1.5-second grace window revoked the family. Logout denied the token for its remaining lifetime. Signed-request replay (409), stale timestamps (401) and forged signatures (401) were rejected; an invalid signature did not consume the nonce.}

\subsection*{Idempotency for API clients}
API clients retry far more often than people double-click, so the idempotency design in Section~\ref{str:idempotency} matters even more here. The OpenAPI document declares the \code{Idempotency-Key} header as a reusable parameter on every unsafe operation, so client generators send it automatically.
\version{The IETF \code{Idempotency-Key} header draft (draft-07) [55] expired in April 2026 without becoming an RFC. The header remains a widely used convention, popularized by payment APIs. Document the exact semantics, such as key format, retention period and the response to a mismatched body, in the API description rather than relying on the draft. The RateLimit header fields draft was revised in May 2026 (draft-11) and is still in progress.}

% ---------------------------------------------------------------- keys
\subsection*{Key layout and access control}
\begin{table}[!htbp]\centering\small
\caption{Redis keys for API credentials. Braces mark Redis Cluster hash tags.}\label{tab:apikeys}
\begin{tabular}{L{.29\linewidth}L{.13\linewidth}L{.17\linewidth}L{.28\linewidth}}\toprule
Key & Type & TTL & Contents \\ \midrule
\code{apikey:<sha256>} & hash & 90 days & user, created \\
\code{user-apikeys:<user>} & set & none & key hashes, for listing and revocation \\
\code{rtfam:{fam}} & hash & 30 days, sliding & user, current token hash; deleting it revokes the family \\
\code{rt:{fam}:<sha256>} & hash & 30 days & state, used\_at, successor, user \\
\code{jwt:deny:<jti>} & string & until token exp & present while a logged-out token is still unexpired \\
\code{nonce:<kid>:<nonce>} & string & 600 s & single-use marker \\
\code{ratelimit:api:<kid>} & string & 60 s & request count in the current window \\
\bottomrule\end{tabular}
\end{table}
All of these keys belong on the \code{noeviction} instance alongside sessions: evicting \code{rtfam} logs users out, and evicting a denylist or nonce entry silently re-enables a revoked token or a replay. The ACL in Section~\ref{sec:redis-security} extends naturally to the new prefixes.


\chapter{Complete Catalyst application laboratory}\label{app:catalyst}

\par\smallskip\noindent\textbf{Abbreviations used in this chapter:} API --- application programming interface; PSGI --- Perl Web Server Gateway Interface.\par\smallskip
The following listings retain the full implementation accompanying Chapter~\ref{ch:catalyst}.
\section{Session storage and application coordination}
\CodeFile{perl}{examples/catalyst-session/lib/SessionLab/Plugin/Session/Store/RedisSafe.pm}{\nolinkurl{catalyst-session/lib/SessionLab/Plugin/Session/Store/RedisSafe.pm}}
\CodeFile{perl}{examples/catalyst-session/lib/SessionLab/Redis.pm}{\nolinkurl{catalyst-session/lib/SessionLab/Redis.pm}}
\CodeFile{perllua}{examples/catalyst-session/lib/SessionLab.pm}{\nolinkurl{catalyst-session/lib/SessionLab.pm}}
\section{The laboratory}\label{sec:catalyst-lab}
\subsection{Session controller and routing}
The controller implements the session patterns in this chapter. The cart actions differ only in where the data lives and whether a lock is held.
\CodeFile{perllua}{examples/catalyst-session/lib/SessionLab/Controller/Root.pm}{\nolinkurl{catalyst-session/lib/SessionLab/Controller/Root.pm}}
\needspace{11\baselineskip}\subsection{PSGI entry point}
\CodeFile{perl}{examples/catalyst-session/app.psgi}{\nolinkurl{catalyst-session/app.psgi}}
\subsection{Session integration driver}
The driver starts Starman with four workers on a free port. It acts as several independent browsers, fires concurrent requests from forked processes that share one session cookie, inspects Redis directly, and cleans up only its own key namespaces. Its \code{END} block localizes \code{$?} so that stopping the server cannot mask a failure as success.
\CodeFile{perl}{examples/28-catalyst-sessions.pl}{\nolinkurl{28-catalyst-sessions.pl}}
\CodeFile{text}{snippets/catalyst-sessions-01.txt}{Driver output on the verification host (abridged)}

\subsection{API credential helpers}
The API controller and the token helpers implement Section~\ref{sec:api-tokens}. Secrets come from the environment (\code{SESSIONLAB_JWT_KEY}, \code{SESSIONLAB_REFRESH_KEY}, \code{SESSIONLAB_SIGNING_KEY}); the defaults are for the laboratory only.
\CodeFile{perl}{examples/catalyst-session/lib/SessionLab/Tokens.pm}{\nolinkurl{catalyst-session/lib/SessionLab/Tokens.pm}}
\needspace{11\baselineskip}\subsection{API controller and atomic Redis operations}
\CodeFile{perllua}{examples/catalyst-session/lib/SessionLab/Controller/API.pm}{\nolinkurl{catalyst-session/lib/SessionLab/Controller/API.pm}}
\subsection{API integration driver}
The driver logs in through the browser flow, creates an API key and token pairs, then attacks its own API: concurrent refreshes from forked processes, tokens it forges with the laboratory key, replayed and forged signatures. It starts the server with a 1.5-second grace window so that the reuse test runs quickly.
\CodeFile{perl}{examples/30-api-tokens.pl}{\nolinkurl{30-api-tokens.pl}}
\CodeFile{text}{snippets/catalyst-sessions-02.txt}{Example 30 output on the verification host}

\chapter{Processing and verification laboratories}\label{app:capstone}
\section{Complete measurement pipeline}
\CodeFile{perl}{examples/36-processing-capstone.pl}{Example 36: recoverable measurement pipeline}
\section{Nested transaction and savepoint laboratory}
\CodeFile{perl}{examples/37-savepoints.pl}{Example 37: nested DBIx::Class transactions and partial rollback}
\section{Bounded retries and cooperative cancellation}
\CodeFile{perl}{examples/38-retry-cancel.pl}{Example 38: classified retries and cancellation}
% Start this long listing on a fresh page to avoid a split continuation crossing the top margin.
\clearpage\section{Contention measurement}
\CodeFile{perl}{examples/39-contention-benchmark.pl}{Example 39: compare serialized and striped counters}
% Keep the short benchmark continuation and Test2 introduction together when space permits.
\needspace{12\baselineskip}\section{Test2 integration}
\CodeFile{perl}{examples/t/processing.t}{Test2 capstone verification in an isolated fixture}

\chapter{Example catalog and API lookup}\label{app:catalog}
All executable examples check their invariants. Example 15 is a
protocol transcript, not a program. The table links to each printed
listing and names its primary prerequisites; the README describes
installation and disposable service setup.
\begin{longtable}{L{.06\linewidth}L{.41\linewidth}L{.31\linewidth}L{.08\linewidth}}
\toprule No. & Source and expected invariant & Requirements & Page\\ \midrule\endhead
01 & \nolinkurl{01-mutex.pl}\newline Scoped mutex & Threaded Perl & \pageref{listing:examples/01-mutex.pl}\\
02 & \nolinkurl{02-semaphore.pl}\newline Permit bound & Threaded Perl & \pageref{listing:examples/02-semaphore.pl}\\
03 & \nolinkurl{03-condition.pl}\newline Predicate rendezvous & Threaded Perl & \pageref{listing:examples/03-condition.pl}\\
04 & \nolinkurl{04-queue.pl}\newline Work queue & Threaded Perl & \pageref{listing:examples/04-queue.pl}\\
05 & \nolinkurl{05-rwlock.pl}\newline Educational reader/writer lock & Threaded Perl & \pageref{listing:examples/05-rwlock.pl}\\
06 & \nolinkurl{06-flock.pl}\newline File lock & Unix & \pageref{listing:examples/06-flock.pl}\\
07 & \nolinkurl{07-sysv.pl}\newline System V semaphore & Linux IPC & \pageref{listing:examples/07-sysv.pl}\\
08 & \nolinkurl{08-pipe.pl}\newline Pipe IPC & Unix & \pageref{listing:examples/08-pipe.pl}\\
09 & \nolinkurl{09-barrier.pl}\newline Educational barrier & Threaded Perl & \pageref{listing:examples/09-barrier.pl}\\
10 & \nolinkurl{10-mysql-advisory.pl}\newline Named locks & MySQL, DBI & \pageref{listing:examples/10-mysql-advisory.pl}\\
11 & \nolinkurl{11-mysql-transfer.pl}\newline Conserved transfer & MySQL, DBI & \pageref{listing:examples/11-mysql-transfer.pl}\\
12 & \nolinkurl{12-mysql-queue.pl}\newline Queue claims & MySQL, DBI & \pageref{listing:examples/12-mysql-queue.pl}\\
13 & \nolinkurl{13-redis-lease.pl}\newline Lease owner checks & Redis client & \pageref{listing:examples/13-redis-lease.pl}\\
14 & \nolinkurl{14-memcached-lease.pl}\newline Cache CAS & Memcached client & \pageref{listing:examples/14-memcached-lease.pl}\\
15 & \nolinkurl{15-memcached-meta.txt}\newline Meta transcript & Memcached protocol & \pageref{listing:examples/15-memcached-meta.txt}\\
16 & \nolinkurl{16-redis-semaphore.pl}\newline Distributed permits & Redis client & \pageref{listing:examples/16-redis-semaphore.pl}\\
17 & \nolinkurl{17-redis-stream-worker.pl}\newline Stream recovery & Redis 6.2+ & \pageref{listing:examples/17-redis-stream-worker.pl}\\
18 & \nolinkurl{18-fenced-write.pl}\newline Fenced destination & Redis, MySQL & \pageref{listing:examples/18-fenced-write.pl}\\
19 & \nolinkurl{19-async-await.pl}\newline Async mutex & Future, IO::Async & \pageref{listing:examples/19-async-await.pl}\\
20 & \nolinkurl{20-scope-release.pl}\newline Cleanup guard & Perl & \pageref{listing:examples/20-scope-release.pl}\\
21 & \nolinkurl{21-futex-mutex.c}\newline Native mutex & Linux, C compiler & \pageref{listing:examples/21-futex-mutex.c}\\
22 & \nolinkurl{22-config-snapshot.pl}\newline Immutable snapshot & Threaded Perl & \pageref{listing:examples/22-config-snapshot.pl}\\
23 & \nolinkurl{23-lock-striping.pl}\newline Striped locks & Threaded Perl & \pageref{listing:examples/23-lock-striping.pl}\\
24 & \nolinkurl{24-try-defer.pl}\newline Experimental cleanup & Perl 5.36+ & \pageref{listing:examples/24-try-defer.pl}\\
25 & \nolinkurl{25-flock-deadline.pl}\newline Bounded file lock & Linux & \pageref{listing:examples/25-flock-deadline.pl}\\
26 & \nolinkurl{26-lock-metrics.pl}\newline Wait/hold metrics & Threaded Perl & \pageref{listing:examples/26-lock-metrics.pl}\\
27 & \nolinkurl{27-mce-loop.pl}\newline Parallel process loop & MCE & \pageref{listing:examples/27-mce-loop.pl}\\
28 & \nolinkurl{28-catalyst-sessions.pl}\newline Session integration & Catalyst, Redis & \pageref{listing:examples/28-catalyst-sessions.pl}\\
29 & \nolinkurl{29-dbic-locking.pl}\newline ORM locking & DBIx::Class, MySQL & \pageref{listing:examples/29-dbic-locking.pl}\\
30 & \nolinkurl{30-api-tokens.pl}\newline API integration & Catalyst, Redis, JWT & \pageref{listing:examples/30-api-tokens.pl}\\
31 & \nolinkurl{31-redis-watch-lease.pl}\newline WATCH lease & Redis client & \pageref{listing:examples/31-redis-watch-lease.pl}\\
32 & \nolinkurl{32-worker-lifecycle.pl}\newline Worker ownership & Moose & \pageref{listing:examples/32-worker-lifecycle.pl}\\
33 & \nolinkurl{33-layered-transaction.pl}\newline Rollback composition & DBIx::Class, SQLite & \pageref{listing:examples/33-layered-transaction.pl}\\
34 & \nolinkurl{34-replay-outbox.pl}\newline Duplicate-safe delivery & DBI, SQLite & \pageref{listing:examples/34-replay-outbox.pl}\\
35 & \nolinkurl{35-deterministic-race.pl}\newline Forced lost update & Threaded Perl & \pageref{listing:examples/35-deterministic-race.pl}\\
36 & \nolinkurl{36-processing-capstone.pl}\newline Recoverable pipeline & DBI, SQLite & \pageref{listing:examples/36-processing-capstone.pl}\\
37 & \nolinkurl{37-savepoints.pl}\newline Nested transactions & DBIx::Class, SQLite/MySQL & \pageref{listing:examples/37-savepoints.pl}\\
38 & \nolinkurl{38-retry-cancel.pl}\newline Classified retries & Perl & \pageref{listing:examples/38-retry-cancel.pl}\\
39 & \nolinkurl{39-contention-benchmark.pl}\newline Contention trials & Threaded Perl & \pageref{listing:examples/39-contention-benchmark.pl}\\
\bottomrule
\end{longtable}
\section{Frequently consulted APIs}
The alphabetical index includes an API group. These links provide
entry points for frequently consulted operations.
\begin{longtable}{L{.32\linewidth}L{.59\linewidth}}
\toprule API & Relevant material\\ \midrule\endhead
\code{lock}, \code{threads::shared} & Shared state (example 01, p.~\pageref{listing:examples/01-mutex.pl})\\
\code{Thread::Queue} & Work distribution (example 04, p.~\pageref{listing:examples/04-queue.pl})\\
\code{flock} & Process locking (example 06, p.~\pageref{listing:examples/06-flock.pl})\\
\code{txn_do}, \code{auto_savepoint} & Layered transactions (Chapter~\ref{ch:enterprise-transactions})\\
\code{FOR UPDATE}, \code{SKIP LOCKED} & Database work queues (Section~\ref{sec:dbqueue})\\
\code{WATCH}, \code{MULTI}, \code{EXEC} & Owner-checked lease (example 31, p.~\pageref{listing:examples/31-redis-watch-lease.pl})\\
\code{XAUTOCLAIM} & Consumer groups (Section~\ref{sec:streams})\\
\code{Future}, \code{await} & Event loops (Section~\ref{sec:async})\\
\bottomrule
\end{longtable}

\chapter{Glossary}\label{ch:glossary}
\section*{Abbreviations and initialisms}
\addcontentsline{toc}{section}{Abbreviations and initialisms}
\begin{longtable}{L{.20\linewidth}L{.72\linewidth}}
\toprule Term & Expansion and meaning\\\midrule\endhead
ACL & Access control list: permissions assigned to a principal.\\
ACID & Atomicity, consistency, isolation and durability: transaction properties.\\
API & Application programming interface: a contract for interacting with software.\\
CAS & Compare-and-swap (compare-and-set in some interfaces): conditional replacement using an expected value or version.\\
CI / CD & Continuous integration / continuous delivery or deployment.\\
CPU & Central processing unit.\\
CSRF & Cross-site request forgery.\\
CSPRNG & Cryptographically secure pseudorandom number generator.\\
DBI / DBD & Database Interface / Database Driver: Perl database interface and driver namespaces.\\
DDL / DML & Data definition language / data manipulation language.\\
FIFO / LIFO & First in, first out / last in, first out.\\
FTS & Full-text search.\\
GCC & GNU Compiler Collection; GNU is a recursive acronym for GNU's Not Unix.\\
HMAC & Hash-based message authentication code.\\
HTTP / HTTPS & Hypertext Transfer Protocol / HTTP over TLS.\\
ID & Identifier.\\
I/O & Input/output.\\
IPC & Interprocess communication.\\
JSON & JavaScript Object Notation.\\
JWT & JSON Web Token.\\
MCE & Many-Core Engine, a Perl parallel processing framework.\\
MVCC & Multi-version concurrency control.\\
OAuth & Open Authorization, commonly used name for the authorization framework.\\
OIDC & OpenID Connect, an identity layer over OAuth 2.0.\\
OS & Operating system.\\
PDF & Portable Document Format.\\
PEL & Pending entries list, in Redis Streams.\\
PID & Process identifier.\\
POSIX & Portable Operating System Interface.\\
PSGI & Perl Web Server Gateway Interface.\\
RDB / AOF & Redis database snapshot / append-only file.\\
REST & Representational state transfer.\\
RPC & Remote procedure call.\\
SHA-256 & Secure Hash Algorithm with a 256-bit digest.\\
SQL & Structured Query Language.\\
SSE & Server-sent events.\\
TCP / UDP & Transmission Control Protocol / User Datagram Protocol.\\
TLS & Transport Layer Security.\\
TTL & Time to live.\\
URI / URL & Uniform Resource Identifier / Uniform Resource Locator.\\
UUID & Universally unique identifier.\\
WAL & Write-ahead log.\\
YAML & YAML Ain't Markup Language (recursive name).\\
ZAB & ZooKeeper Atomic Broadcast.\\
\bottomrule
\end{longtable}
\section*{Technical terms}

\begin{description}[style=unboxed,leftmargin=0pt,labelsep=.5em,itemsep=3pt,parsep=0pt,font=\sffamily\bfseries\color{navy}]\small
\item[Advisory lock] A lock that only works if every participant checks it; nothing stops code that ignores it. Perl's \code{lock}, \code{flock} and MySQL \code{GET_LOCK} are advisory.
\item[At-least-once delivery] Under the transport's durability, retention and acknowledgement assumptions, an unacknowledged message is eligible for redelivery. Duplicates are possible; consumers must make effects idempotent. Data loss outside those assumptions is not excluded.
\item[Atomic operation] A change that happens completely or not at all, with no observable halfway state, such as a single Redis command or one SQL \code{UPDATE}.
\item[Attempt ID] A generation associated with a particular job claim; a late worker cannot complete a replacement worker's claim.
\item[Backoff and jitter] Waiting longer after each failed attempt (backoff), with a random component (jitter), so that competing retries do not collide again in lockstep.
\item[Barrier] A meeting point where a group of threads waits until all members have arrived.
\item[CAS (compare-and-swap)] Update a value only if its expected value or version still matches; otherwise do nothing and report a conflict. Memcached compares a server-issued version token.
\item[Checkpoint] Durable progress recorded with the effects it describes, allowing a restarted pipeline to resume without losing committed work.
\item[Consensus] A protocol (Raft, ZAB) that lets a group of servers agree on one ordered history even when some fail; etcd and ZooKeeper are built on it.
\item[Consumer group] A Redis Streams feature that delivers each entry to one member of a group and tracks unacknowledged entries in a pending entries list (PEL).
\item[Critical section] The code that must not run in two places at once because it touches shared state.
\item[Deadlock] Two or more parties each waiting for a lock the other holds, so none can proceed.
\item[Eviction] A cache deleting keys to free memory. Harmless for caches, harmful for sessions and locks.
\item[Failover] Promoting a replica to primary after the primary fails; with asynchronous replication the newest writes can be lost.
\item[Fencing token] A number that increases with every new lock holder and is checked by the protected resource, which rejects older numbers.
\item[Gap lock, next-key lock] InnoDB locks on the space between index records (gap) or on a record plus the gap before it (next-key); they block inserts into ranges.
\item[Grace window] A short period after a refresh token is rotated during which presenting it again returns the same successor instead of revoking the family; it absorbs concurrent and retried refreshes.
\item[HMAC] Hash-based message authentication code: a keyed construction (for example HMAC-SHA256). Only holders of the key can compute it, so it proves that a message came from them and was not altered.
\item[Idempotent] Safe to repeat: doing it twice has the same effect as doing it once.
\item[Isolation level] How much of other transactions' work a transaction can see (READ COMMITTED, REPEATABLE READ, and so on), and therefore which locks it takes.
\item[JWT] JSON Web Token: a signed set of claims (user, audience, expiry) that a server can verify without a lookup. Used here as a short-lived access token.
\item[Lease] A lock with an expiry time. It frees itself if the holder disappears, but can also expire while the holder is still working.
\item[Livelock] Workers keep reacting to each other and retrying without making progress.
\item[Lost update] Two read-modify-write operations interleave and one overwrites the other's change.
\item[Monotonic clock] A clock that never jumps backwards (\code{CLOCK_MONOTONIC}); use it for timeouts and durations, never wall-clock time.
\item[Mutex] A lock that only one holder can own at a time (mutual exclusion).
\item[MVCC] Multi-version concurrency control: readers see a consistent snapshot without blocking writers.
\item[Nonce] A value used only once. A server that remembers recent nonces can reject a replayed request.
\item[Open file description] The kernel object created by \code{open}; \code{flock} locks belong to it, and \code{fork} shares it.
\item[OpenAPI security scheme] The part of an OpenAPI description that says how clients authenticate (API key, HTTP bearer, OAuth~2.0, mutual TLS); operations then state which schemes they require.
\item[Optimistic and pessimistic concurrency] Optimistic: proceed without locking and detect conflicts at write time. Pessimistic: lock first so conflicts cannot happen.
\item[Outbox] An event record committed in the same transaction as a business change, then relayed independently.
\item[Preforking server] A server (Starman, uWSGI) that forks a pool of worker processes; each worker has its own memory and its own locks.
\item[Race condition] A bug whose outcome depends on the timing of concurrent operations.
\item[Refresh token] A long-lived credential that a client exchanges for new access tokens. With rotation, each use returns a new one and invalidates the old.
\item[Semaphore] A counter of available permits; acquiring takes one, releasing returns one, and acquirers wait at zero.
\item[Starvation] A participant is postponed indefinitely even though others make progress.
\item[Token family] All refresh tokens descended by rotation from one sign-in. Reuse of any old member revokes the whole family.
\item[TTL] Time to live: how long a key or lease exists before it expires automatically.
\end{description}

% ================================================================ chapter 12
\chapter{References and verification}\label{ch:refs}
\section{Current release verification}\label{sec:current-verification}
\IfFileExists{generated/release-verification.tex}{\textbf{Current run:} 2026-11-25T00:24:39+00:00. \href{https://github.com/interrupt21h/Modern-Perl-Concurrency/actions/runs/45517252287}{GitHub Actions run 45517252287}.
Trigger revision: \nolinkurl{fa780acf211ab86c314a1979d8ff79803907cf0f}. The final release commit and PDF checksums are recorded with the downloadable assets.
\begin{longtable}{L{.30\linewidth}L{.62\linewidth}}\toprule Component & Recorded value \\ \midrule\endhead
Primary suite & 39 passed; 0 skipped\\
Additional checks & SQLite savepoints; Test2 via prove and yath; seven Valkey integration examples; Perl::Tidy idempotence\\
Operating system & Ubuntu 24.04.5 LTS\\
Perl & v5.38.2; ithreads define\\
MySQL & 8.4.13\\
Redis & 7.0.15\\
Valkey & 8.1.12\\
Memcached & 1.6.24\\
SQLite & 3.42.0\\
DBI & 1.643\\
DBIx::Class & 0.082843\\
DBD::SQLite & 1.74\\
DBD::mysql & 4.052\\
Redis & 2.000\\
Cache::Memcached::Fast & 0.28\\
Catalyst & 5.90131\\
Moose & 2.2207\\
MCE & 1.889\\
Future & 0.54\\
Test2::V0 & 0.000159\\
\bottomrule\end{longtable}
The manifest records the test phase before compilation. The release is published only after subsequent PDF preflight and packaging succeed. It does not establish compatibility with untested versions.
}{No recorded release manifest is present in this source checkout. Run the complete publishing workflow to generate current evidence; the historical record below is not a substitute.}

\section{Historical verification record}
The following integration record is retained from Edition 1.2. Edition 1.3 added example 31; Edition 2.0 added application architecture and examples 32--35. Edition 2.1 reorganized the publication and added the processing capstone and further laboratories. Edition 2.1.1 consolidates the presentation and records current release evidence separately. The historical results below are not a fresh integration run. See the revision test log for current checks.

Primary documentation was consulted on November 24, 2026. Perl documentation resolved to 5.44.0; MySQL documentation to 8.4 and 9.7; Redis documentation to 8.4; Valkey documentation to 9.0. The examples were executed with \code{run-tests.sh} three consecutive times after the final change on Linux with threaded Perl 5.38.2, GCC (example 21), MCE 1.889, Catalyst 5.90131 with \code{Catalyst::Plugin::Session} 0.43 and Starman 0.4017 (examples 28 and 30), \code{Crypt::JWT} 0.035 and \code{YAML::XS} 0.89 (example 30), \code{openapi-spec-validator} 0.9.0 for the OpenAPI document, DBIx::Class 0.082843 (example 29), Sereal 5.004, DBI 1.643, \code{DBD::mysql} 4.052, MySQL 8.0.46, Redis 7.0.15 with \code{Redis.pm} 2.000, Memcached 1.6.24 with \code{Cache::Memcached::Fast} 0.28, IO::Async 0.803 and \code{Future::AsyncAwait} 0.66. Examples 13 and 16 were additionally run under the least-privilege ACL user in Section~\ref{sec:redis-security}. Native Redis 8.4 and Valkey 9.0 commands are documented from primary sources but were not available on the verification host. These historical results establish neither a throughput baseline nor a fairness guarantee; see the measured contention trials in Part VI.
\begin{longtable}{L{.30\linewidth}L{.62\linewidth}}\toprule Example & Verified invariant \\ \midrule\endhead
01--09, 20 & Exact totals, permit bounds, rendezvous, queue completeness, lock release on exception \\
21 & 8 threads $\times$ 10\textsuperscript{6} non-atomic increments, no lost update; clean under ThreadSanitizer \\
22--23 & Snapshots always self-consistent and monotonic; striped transfers conserve the total under concurrent audits \\
24--25 & Lock released on return, exception and nesting; \code{flock} timeout honored and acquisition after release \\
26--27 & Wait/hold samples complete and plausible; MCE result exact with every chunk gathered once \\
29 & DBIx::Class: 200 opposing transfers under \code{FOR UPDATE} preserve the total; 80 optimistic edits kept with conflicts retried; \code{SKIP LOCKED} claims never duplicated \\
28 & Catalyst under Starman (4 workers), 29 checks: random IDs, cookie flags, atomic TTL, Sereal, fixation, lost update reproduced and fixed two ways, idempotency under concurrency and crash, rate limit, revocation, forged record refused \\
30 & API under Starman, 33 checks: only key hashes stored; per-key 429 with RateLimit headers; \code{alg: none}, missing \code{exp} and missing \code{typ} refused; naive rotation forked a family, atomic rotation gave one successor to six callers; reuse revoked the family; logout denylist; signed-request replay, staleness and forgery refused without burning the nonce; OpenAPI coverage \\
10--12 & Advisory lock ownership; balance conservation under 4$\times$100 concurrent transfers; 220 jobs completed exactly once; lease reclaim, dead-letter, stale completion rejected \\
13 & Four competing processes: no overlapping critical sections, strictly increasing fences \\
14 & Four competing processes serialized (best effort) \\
15 & Protocol transcript replayed against Memcached 1.6.24 \\
16 & Ten processes, three permits: peak concurrency never above three \\
17 & 200 entries; crashed consumer's pending entries reclaimed; each effect applied once \\
18 & Stale fence rejected by MySQL; stale owner cannot release successor's lease; counter regression fails safe \\
19 & Concurrency bound respected; no lost update with \code{Future::Mutex} \\
\bottomrule\end{longtable}
\section{Application architecture laboratories}
Examples 32--39 use independently constructed synthetic data. Example 32 checks PID-local Moose resource ownership; 33 checks transaction rollback and duplicate-request behavior with DBIx::Class and SQLite; 34 checks durable source/outbox consistency and replay against a separate SQLite destination; 35 forces a lost-update interleaving and verifies the corrected critical section. SQLite tests demonstrate application transaction boundaries, not MySQL row-lock or isolation semantics. Examples 36--39 cover replay-safe processing, savepoints, classified retry/cancellation and contention measurement. Consult the release build log for actual test results and skipped dependencies.

\section{References}
\begingroup\small\setlength{\parskip}{2pt}
References are grouped by subject; original numeric citations are preserved.
\subsection{Perl, Unix and asynchronous execution}
\refdoc{[1] Perl threading tutorial}{https://perldoc.perl.org/perlthrtut}
\refdoc{[2] Shared variables and condition synchronization}{https://perldoc.perl.org/threads::shared}
\refdoc{[3] Thread semaphores}{https://perldoc.perl.org/Thread::Semaphore}
\refdoc{[4] Thread queues}{https://perldoc.perl.org/Thread::Queue}
\refdoc{[5] Linux flock(2)}{https://man7.org/linux/man-pages/man2/flock.2.html}
\refdoc{[6] Linux futex(2)}{https://man7.org/linux/man-pages/man2/futex.2.html}
\refdoc{[7] System V semaphore wrapper}{https://perldoc.perl.org/IPC::Semaphore}
\refdoc{[8] POSIX semaphore overview}{https://man7.org/linux/man-pages/man7/sem_overview.7.html}
\refdoc{[9] MCE worker loops}{https://metacpan.org/pod/MCE::Loop}
\refdoc{[10] IO::Async}{https://metacpan.org/pod/IO::Async}
\refdoc{[16] MCE shared objects}{https://metacpan.org/pod/MCE::Shared}
\refdoc{[17] Linux record and OFD locks}{https://man7.org/linux/man-pages/man2/fcntl_locking.2.html}
\refdoc{[27] threads module documentation}{https://perldoc.perl.org/threads}
\refdoc{[28] Try/catch and defer blocks (perlsyn)}{https://perldoc.perl.org/perlsyn}
\refdoc{[29] Future::AsyncAwait}{https://metacpan.org/pod/Future::AsyncAwait}
\refdoc{[30] Future::Mutex}{https://metacpan.org/pod/Future::Mutex}
\refdoc{[46] perl5440delta}{https://perldoc.perl.org/perl5440delta}
\refdoc{[76] Future cancellation semantics}{https://metacpan.org/pod/Future}
\refdoc{[79] Linux rename semantics}{https://man7.org/linux/man-pages/man2/rename.2.html}
\refdoc{[83] Time::HiRes monotonic clock}{https://perldoc.perl.org/Time::HiRes}
\subsection{Database transactions and storage}
\refdoc{[11] InnoDB locking}{https://dev.mysql.com/doc/refman/8.4/en/innodb-locking.html}
\refdoc{[12] InnoDB locking reads}{https://dev.mysql.com/doc/refman/8.4/en/innodb-locking-reads.html}
\refdoc{[13] Handling InnoDB deadlocks}{https://dev.mysql.com/doc/refman/8.4/en/innodb-deadlocks-handling.html}
\refdoc{[14] MySQL named locking functions}{https://dev.mysql.com/doc/refman/8.4/en/locking-functions.html}
\refdoc{[15] DBI documentation}{https://metacpan.org/pod/DBI}
\refdoc{[18] DBD::mysql driver and changes}{https://metacpan.org/dist/DBD-mysql/changes}
\refdoc{[31] DBD::MariaDB}{https://metacpan.org/pod/DBD::MariaDB}
\refdoc{[32] MySQL release model: Innovation and LTS}{https://dev.mysql.com/doc/refman/9.7/en/mysql-releases.html}
\refdoc{[47] PostgreSQL advisory locks}{https://www.postgresql.org/docs/current/explicit-locking.html}
\refdoc{[65] DBIx::Class storage and nested transactions}{https://metacpan.org/pod/DBIx::Class::Storage}
\refdoc{[67] InnoDB error handling}{https://dev.mysql.com/doc/refman/8.4/en/innodb-error-handling.html}
\refdoc{[68] Converting MyISAM tables}{https://dev.mysql.com/doc/refman/8.4/en/converting-tables-to-innodb.html}
\refdoc{[70] SQLite UPSERT}{https://www.sqlite.org/lang_upsert.html}
\refdoc{[75] DBIx::Class DBI storage and savepoints}{https://metacpan.org/pod/DBIx::Class::Storage::DBI}
\refdoc{[77] SQLite transaction semantics}{https://www.sqlite.org/lang_transaction.html}
\refdoc{[78] SQLite WAL and concurrency}{https://www.sqlite.org/wal.html}
\refdoc{[81] SQLite savepoint and release semantics}{https://www.sqlite.org/lang_savepoint.html}
\refdoc{[82] DBD::SQLite transaction semantics}{https://metacpan.org/pod/DBD::SQLite}
\subsection{Redis, Valkey and Memcached}
\refdoc{[19] Redis distributed locks}{https://redis.io/docs/latest/develop/clients/patterns/distributed-locks/}
\refdoc{[20] Redis SET}{https://redis.io/docs/latest/commands/set/}
\refdoc{[21] Perl Redis client}{https://metacpan.org/pod/Redis}
\refdoc{[22] Memcached basic protocol}{https://docs.memcached.org/protocols/basic/}
\refdoc{[23] Memcached use cases}{https://docs.memcached.org/userguide/usecases/}
\refdoc{[24] Perl Cache::Memcached::Fast}{https://metacpan.org/pod/Cache::Memcached::Fast}
\refdoc{[25] Memcached meta protocol}{https://docs.memcached.org/protocols/meta/}
\refdoc{[26] Memcached protocol specification}{https://github.com/memcached/memcached/blob/master/doc/protocol.txt}
\refdoc{[33] Redis DELEX}{https://redis.io/docs/latest/commands/delex/}
\refdoc{[34] What's new in Redis 8.4}{https://redis.io/docs/latest/develop/whats-new/8-4/}
\refdoc{[35] Valkey DELIFEQ}{https://valkey.io/commands/delifeq/}
\refdoc{[36] Valkey SET (IFEQ)}{https://valkey.io/commands/set/}
\refdoc{[37] Redis WAITAOF}{https://redis.io/docs/latest/commands/waitaof/}
\refdoc{[38] Redis XAUTOCLAIM}{https://redis.io/docs/latest/commands/xautoclaim/}
\refdoc{[39] Redis ACL}{https://redis.io/docs/latest/operate/oss_and_stack/management/security/acl/}
\refdoc{[40] Redis licenses}{https://redis.io/legal/licenses/}
\refdoc{[41] Redis::Fast}{https://metacpan.org/pod/Redis::Fast}
\refdoc{[42] Net::Async::Redis}{https://metacpan.org/pod/Net::Async::Redis}
\refdoc{[61] Redis transactions and WATCH semantics}{https://redis.io/docs/latest/develop/using-commands/transactions/}
\refdoc{[62] Redis EXEC reply semantics}{https://redis.io/docs/latest/commands/exec/}
\subsection{Distributed coordination and production operations}
\refdoc{[43] M. Kleppmann, How to do distributed locking}{https://martin.kleppmann.com/2016/02/08/how-to-do-distributed-locking.html}
\refdoc{[44] Kubernetes Leases}{https://kubernetes.io/docs/concepts/architecture/leases/}
\refdoc{[45] etcd API guarantees and revisions}{https://etcd.io/docs/v3.5/learning/api/}
\refdoc{[66] Kubernetes CronJob controller}{https://kubernetes.io/docs/concepts/workloads/controllers/cron-jobs/}
\refdoc{[71] Transactional outbox pattern}{https://docs.aws.amazon.com/prescriptive-guidance/latest/cloud-design-patterns/transactional-outbox.html}
\refdoc{[72] Kubernetes Pod termination}{https://kubernetes.io/docs/concepts/workloads/pods/pod-lifecycle/}
\refdoc{[73] OpenTelemetry metrics concepts}{https://opentelemetry.io/docs/concepts/signals/metrics/}
\refdoc{[74] Kubernetes leader election and fencing limitations}{https://pkg.go.dev/k8s.io/client-go/tools/leaderelection}
\subsection{Application frameworks, API security and verification}
\refdoc{[48] RFC 9700: Best Current Practice for OAuth 2.0 Security}{https://www.rfc-editor.org/rfc/rfc9700}
\refdoc{[49] RFC 8725: JSON Web Token Best Current Practices}{https://www.rfc-editor.org/rfc/rfc8725}
\refdoc{[50] RFC 9068: JWT Profile for OAuth 2.0 Access Tokens}{https://www.rfc-editor.org/rfc/rfc9068}
\refdoc{[51] OpenAPI Specification 3.2.0}{https://spec.openapis.org/oas/v3.2.0.html}
\refdoc{[52] Crypt::JWT}{https://metacpan.org/pod/Crypt::JWT}
\refdoc{[53] OpenAPI::Modern}{https://metacpan.org/pod/OpenAPI::Modern}
\refdoc{[54] RateLimit header fields for HTTP (IETF draft)}{https://datatracker.ietf.org/doc/draft-ietf-httpapi-ratelimit-headers/}
\refdoc{[55] The Idempotency-Key HTTP header field (IETF draft)}{https://datatracker.ietf.org/doc/draft-ietf-httpapi-idempotency-key-header/}
\refdoc{[56] RFC 9421: HTTP Message Signatures}{https://www.rfc-editor.org/rfc/rfc9421}
\refdoc{[57] RFC 9449: OAuth 2.0 Demonstrating Proof of Possession (DPoP)}{https://www.rfc-editor.org/rfc/rfc9449}
\refdoc{[58] Okta: refresh token rotation and grace period}{https://developer.okta.com/docs/guides/refresh-tokens/main/}
\refdoc{[59] Amazon Cognito: RefreshTokenRotationType}{https://docs.aws.amazon.com/cognito-user-identity-pools/latest/APIReference/API_RefreshTokenRotationType.html}
\refdoc{[60] Auth0: refresh token rotation}{https://auth0.com/docs/secure/tokens/refresh-tokens/refresh-token-rotation}
\refdoc{[63] Moose attributes}{https://metacpan.org/pod/Moose::Manual::Attributes}
\refdoc{[64] Catalyst component lifecycle}{https://metacpan.org/pod/Catalyst::Component}
\refdoc{[69] Test2 harness}{https://metacpan.org/pod/Test2::Harness}
\refdoc{[80] Test2::V0}{https://metacpan.org/pod/Test2::V0}
\endgroup

\clearpage
\renewcommand{\indexname}{Subject and API index}
\phantomsection
\addcontentsline{toc}{chapter}{Subject and API index}
\printindex
\end{document}
